% ECONOMIC GEOLOGY: The Soil — Geology Upper Sixth — Cours signé OnBuch+
% Official MINESEC syllabus. Compiler avec : tectonic GEO04-economic-geology-the-soil.tex
\newcommand{\DOCMATIERE}{Geology}
\newcommand{\DOCNIVEAU}{Upper Sixth}
\newcommand{\DOCMODULE}{Upper Sixth — Geology}
\newcommand{\DOCLECON}{4}
\newcommand{\DOCTITRE}{ECONOMIC GEOLOGY: The Soil}
% OnBuch+ — préambule « cours signés » ENRICHI (compilé avec tectonic / XeLaTeX)
% Système visuel « Pop » d'OnBuch+ : fond crème (jamais blanc), bordures encre
% épaisses, ombres « dures » décalées, titres Archivo Black en capitales.
% Version MATHÉMATIQUES : graphes (pgfplots/tikz), boîtes pédagogiques
% (définition, propriété, méthode, attention, le savais-tu, exercice résolu,
% exercice, TP, bilan), page de garde et signature OnBuch+ ancrée en bas.
%
% Macros attendues dans main.tex AVANT \input{preamble} :
%   \DOCMATIERE  \DOCNIVEAU  \DOCMODULE  \DOCLECON (numéro)  \DOCTITRE
\documentclass[11pt]{article}

\usepackage{fontspec}
\usepackage[english]{babel}
\usepackage[a4paper,margin=2cm,top=3.4cm,bottom=2.8cm,headheight=1.4cm,footskip=1.3cm]{geometry}
\usepackage{amsmath,amssymb,amsfonts}
\usepackage{mathtools} % \xmapsto, \coloneqq... souvent utilisés par les modèles
\usepackage{cancel} % simplifications barrées (bilans de forces)
% \abs{z} (module, valeur absolue) et \norm{u} (norme), utilisés par les modèles
\providecommand{\abs}{}\let\abs\relax\DeclarePairedDelimiter{\abs}{\lvert}{\rvert}
\providecommand{\norm}{}\let\norm\relax\DeclarePairedDelimiter{\norm}{\lVert}{\rVert}
\usepackage[table]{xcolor} % [table] : fournit aussi \rowcolors (alternance de lignes)
\usepackage{pagecolor}
\usepackage{tikz}
\usetikzlibrary{shapes.symbols,circuits.logic.US,circuits.logic.IEC,arrows.meta,positioning,calc,shapes.geometric,shapes.misc,shapes.arrows,decorations.pathmorphing,decorations.pathreplacing,decorations.markings,patterns,shadows,mindmap,trees,fit,backgrounds,matrix,intersections,angles,quotes}
\usepackage{pgfplots}
\usepackage[version=4]{mhchem} % équations nucléaires \ce{^{235}_{92}U}
\usepackage[european,siunitx]{circuitikz} % schémas électriques
\pgfplotsset{compat=1.18}
\usepgfplotslibrary{fillbetween,groupplots} % aires entre courbes (calcul intégral)
\usepackage{enumitem}
\usepackage{fancyhdr}
% [most] charge toutes les bibliothèques tcolorbox (listings, minted, raster,
% documentation...) et gonfle inutilement la mémoire TeX sur un document long
% (~300 boîtes) : on ne charge que ce qui est réellement utilisé.
\usepackage[skins,breakable]{tcolorbox}
\usepackage{graphicx}
\usepackage{colortbl}
\usepackage{booktabs}
\usepackage{tabularx}
\usepackage{diagbox} % case d'en-tête diagonale des tableaux à double entrée
% Colonnes extensibles centrée / gauche / droite (C, L, R), souvent utilisées par les modèles
\newcolumntype{C}{>{\centering\arraybackslash}X}
\newcolumntype{L}{>{\raggedright\arraybackslash}X}
\newcolumntype{R}{>{\raggedleft\arraybackslash}X}
\usepackage{array}
\usepackage{multicol}
\usepackage{siunitx}
% parse-numbers=false : accepte aussi \SI{2,5 \times 10^{-3}}{...}
\sisetup{locale=UK,output-decimal-marker={.},group-minimum-digits=4,parse-numbers=false}
\DeclareSIUnit\ohm{\ensuremath{\symup{\Omega}}} % U+2126 absent de la police
\DeclareSIUnit\division{div} % réglages d'oscilloscope (ms/div, V/div)
\DeclareSIUnit\year{yr}\DeclareSIUnit\annum{yr}\DeclareSIUnit\Ma{Ma}\DeclareSIUnit\tr{tr} % tours (vitesse de rotation, ex. \SI{1500}{\tr\per\minute})
\usepackage{qrcode}
% \degree : commande fréquemment utilisée par les modèles pour un angle ou une
% température en dehors de \SI{}{} (non fournie par les paquets chargés).
\providecommand{\degree}{^{\circ}}
% \qed (fin de démonstration), utilisé par les modèles sans amsthm
\providecommand{\qed}{\hfill$\square$}
% \barre : double barre des valeurs interdites dans un tableau de variations (tkz-tab)
\providecommand{\barre}{\ensuremath{\|}}
% environnement proof (amsthm non chargé) utilisé par les modèles
\ifdefined\proof\else\newenvironment{proof}[1][Proof]{\par\noindent\textit{#1.}\ }{\qed\par}\fi
\usepackage{lastpage}
\usepackage{hyperref}
\hypersetup{hidelinks}

% ============================================================
% Palette « Pop » OnBuch+ — mêmes hex que la classe P (plus_ui.dart)
% ============================================================
\definecolor{poporange}{HTML}{F59321}
\definecolor{popdark}{HTML}{E07A0C}
\definecolor{popcream}{HTML}{FFF6E8}
\definecolor{popink}{HTML}{1C1714}
\definecolor{poppurple}{HTML}{7A5AE0}
\definecolor{popgreen}{HTML}{2E9E6A}
\definecolor{poppink}{HTML}{E0457B}
\definecolor{popblue}{HTML}{2F7DE1}
\definecolor{popgold}{HTML}{C98A12}
\definecolor{popmuted}{HTML}{8A7F76}
\definecolor{popline}{HTML}{EADFD2}
\definecolor{popgray}{HTML}{8A7F76}
\colorlet{popgrey}{popgray}
% Alias tolérants (le modèle nomme parfois les couleurs différemment)
\colorlet{popwhite}{white}
\colorlet{popblack}{popink}
\colorlet{popred}{poppink}
\colorlet{popyellow}{popgold}
% Teintes claires pour fonds de tableaux / zones
\colorlet{popblueL}{popblue!10!white}
\colorlet{poporangeL}{poporange!14!white}
\colorlet{popgreenL}{popgreen!12!white}
\colorlet{poppurpleL}{poppurple!12!white}
\colorlet{poppinkL}{poppink!12!white}
\colorlet{popgoldL}{popgold!14!white}

\pagecolor{popcream}

% ============================================================
% Typographie
% ============================================================
\defaultfontfeatures{Ligatures=TeX}
\setmainfont{PlusJakartaSans-Medium.ttf}[Path=../../../fonts/,BoldFont=PlusJakartaSans-Bold.ttf]
% Maths en sans-serif (Fira Math) avec chiffres et lettres droites de Plus
% Jakarta, pour que les formules se fondent dans le texte.
\usepackage[math-style=ISO,bold-style=ISO]{unicode-math}
\setmathfont{FiraMath-Regular.otf}[Path=../../../fonts/]
\setmathfont{PlusJakartaSans-Medium.ttf}[Path=../../../fonts/,range={up/num,up/{Latin,latin},\mathup}]
% (icomma retiré : décimales avec point en anglais)
% Caractères mathématiques Unicode absents des polices de texte (Plus Jakarta,
% Archivo Black) : rendus par leur équivalent mathématique, en texte comme en
% formule — sinon ils s'affichent en carré vide (ex. « Arithmétique dans ℤ »).
\usepackage{newunicodechar}
\newunicodechar{ℝ}{\ensuremath{\mathbb{R}}}
\newunicodechar{ℤ}{\ensuremath{\mathbb{Z}}}
\newunicodechar{ℂ}{\ensuremath{\mathbb{C}}}
\newunicodechar{ℕ}{\ensuremath{\mathbb{N}}}
\newunicodechar{ℚ}{\ensuremath{\mathbb{Q}}}
\newunicodechar{𝔻}{\ensuremath{\mathbb{D}}}
\newunicodechar{α}{\ensuremath{\alpha}}
\newunicodechar{β}{\ensuremath{\beta}}
\newunicodechar{γ}{\ensuremath{\gamma}}
\newunicodechar{δ}{\ensuremath{\delta}}
\newunicodechar{ε}{\ensuremath{\varepsilon}}
\newunicodechar{θ}{\ensuremath{\theta}}
\newunicodechar{λ}{\ensuremath{\lambda}}
\newunicodechar{μ}{\ensuremath{\mu}}
\newunicodechar{σ}{\ensuremath{\sigma}}
\newunicodechar{τ}{\ensuremath{\tau}}
\newunicodechar{φ}{\ensuremath{\varphi}}
\newunicodechar{ψ}{\ensuremath{\psi}}
\newunicodechar{ω}{\ensuremath{\omega}}
\newunicodechar{Σ}{\ensuremath{\Sigma}}
% majuscules produites par \MakeUppercase dans les titres (α-aminés -> Α)
\newunicodechar{Α}{\ensuremath{\alpha}}
\newunicodechar{Β}{\ensuremath{\beta}}
\newunicodechar{Γ}{\ensuremath{\Gamma}}
\newunicodechar{Δ}{\ensuremath{\Delta}}
\newunicodechar{Ε}{\ensuremath{\varepsilon}}
\newunicodechar{Θ}{\ensuremath{\Theta}}
\newunicodechar{Λ}{\ensuremath{\Lambda}}
\newunicodechar{Μ}{\ensuremath{\mu}}
\newunicodechar{Τ}{\ensuremath{\tau}}
\newunicodechar{Φ}{\ensuremath{\Phi}}
\newunicodechar{Ψ}{\ensuremath{\Psi}}
\newunicodechar{Ω}{\ensuremath{\Omega}}
\newunicodechar{↦}{\ensuremath{\mapsto}}
\newunicodechar{∈}{\ensuremath{\in}}
\newunicodechar{∉}{\ensuremath{\notin}}
\newunicodechar{∧}{\ensuremath{\wedge}}
\newunicodechar{⇔}{\ensuremath{\Leftrightarrow}}
\newunicodechar{⟺}{\ensuremath{\Longleftrightarrow}}
\newunicodechar{⇒}{\ensuremath{\Rightarrow}}
\newunicodechar{⟹}{\ensuremath{\Longrightarrow}}
\newunicodechar{∘}{\ensuremath{\circ}}
\newunicodechar{∪}{\ensuremath{\cup}}
\newunicodechar{∩}{\ensuremath{\cap}}
\newunicodechar{≡}{\ensuremath{\equiv}}
\newunicodechar{⊂}{\ensuremath{\subset}}
\newunicodechar{⊥}{\ensuremath{\perp}}
\newunicodechar{∃}{\ensuremath{\exists}}
\newunicodechar{∀}{\ensuremath{\forall}}
\newunicodechar{∛}{\ensuremath{\sqrt[3]{\,}}}
\newunicodechar{∜}{\ensuremath{\sqrt[4]{\,}}}
\newunicodechar{⃗}{\ensuremath{{}^{\to}}}
\newunicodechar{ⁿ}{\ifmmode^{n}\else\textsuperscript{\ensuremath{n}}\fi}
\newunicodechar{ˣ}{\ifmmode^{x}\else\textsuperscript{\ensuremath{x}}\fi}
\newunicodechar{ᵇ}{\ifmmode^{b}\else\textsuperscript{\ensuremath{b}}\fi}
\newunicodechar{ʳ}{\ifmmode^{r}\else\textsuperscript{\ensuremath{r}}\fi}
\newunicodechar{ᵘ}{\ifmmode^{u}\else\textsuperscript{\ensuremath{u}}\fi}
\newunicodechar{ʸ}{\ifmmode^{y}\else\textsuperscript{\ensuremath{y}}\fi}
\newunicodechar{ᵃ}{\ifmmode^{a}\else\textsuperscript{\ensuremath{a}}\fi}
\newunicodechar{ᵉ}{\ifmmode^{e}\else\textsuperscript{\ensuremath{e}}\fi}
\newunicodechar{ᵏ}{\ifmmode^{k}\else\textsuperscript{\ensuremath{k}}\fi}
\newunicodechar{ᵖ}{\ifmmode^{p}\else\textsuperscript{\ensuremath{p}}\fi}
\newunicodechar{ᴺ}{\ifmmode^{N}\else\textsuperscript{\ensuremath{N}}\fi}
\newunicodechar{ᶜ}{\ifmmode^{c}\else\textsuperscript{\ensuremath{c}}\fi}
\newunicodechar{ʰ}{\ifmmode^{h}\else\textsuperscript{\ensuremath{h}}\fi}
\newunicodechar{ᵠ}{\ifmmode^{q}\else\textsuperscript{\ensuremath{q}}\fi}
\newunicodechar{ₙ}{\ifmmode_{n}\else\textsubscript{\ensuremath{n}}\fi}
\newunicodechar{ₐ}{\ifmmode_{a}\else\textsubscript{\ensuremath{a}}\fi}
\newunicodechar{ᵢ}{\ifmmode_{i}\else\textsubscript{\ensuremath{i}}\fi}
\newunicodechar{ᵣ}{\ifmmode_{r}\else\textsubscript{\ensuremath{r}}\fi}
\newunicodechar{ₖ}{\ifmmode_{k}\else\textsubscript{\ensuremath{k}}\fi}
\newunicodechar{ᵦ}{\ifmmode_{\beta}\else\textsubscript{\ensuremath{\beta}}\fi}
\newunicodechar{ₓ}{\ifmmode_{x}\else\textsubscript{\ensuremath{x}}\fi}
\newfontfamily\popdisplay{ArchivoBlack-Regular.ttf}[Path=../../../fonts/]
\newfontfamily\poplogo{SpaceGrotesk-Bold.ttf}[Path=../../../fonts/]
\newfontfamily\popsemi{PlusJakartaSans-SemiBold.ttf}[Path=../../../fonts/]
\newfontfamily\popxbold{PlusJakartaSans-ExtraBold.ttf}[Path=../../../fonts/]
\newfontfamily\popxboldls{PlusJakartaSans-ExtraBold.ttf}[Path=../../../fonts/,LetterSpace=3.0]

\setlist[enumerate]{leftmargin=1.7em,itemsep=5pt,topsep=4pt}
\setlist[itemize]{leftmargin=1.7em,itemsep=4pt,topsep=4pt,label={\color{poporange}\textbullet}}
\setlength{\parindent}{0pt}
\setlength{\parskip}{5pt}
\renewcommand{\arraystretch}{1.25}

% pgfplots : style Pop par défaut
\pgfplotsset{
  every axis/.append style={axis line style={popink,line width=1pt},tick style={popink},
    grid style={popline,line width=0.5pt},label style={font=\small},tick label style={font=\footnotesize}},
  popcurve/.style={poporange,line width=1.8pt,no markers,smooth},
}
% Même style disponible dans un \draw TikZ simple (hors pgfplots)
\tikzset{popcurve/.style={poporange,line width=1.8pt}}
\tikzset{popgrid/.style={popline,very thin}}
\colorlet{popcurve}{poporange} % le nom du style est parfois utilisé comme couleur

% ============================================================
% Pill / squiggle
% ============================================================
\newcommand{\poppill}[3]{%
  \tikz[baseline=-0.55ex]\node[draw=popink,line width=1.2pt,rounded corners=2.6pt,%
    fill=#2,inner xsep=6.5pt,inner ysep=3pt]{%
    \popxboldls\fontsize{7}{7}\selectfont\color{#3}\MakeUppercase{#1}};%
}
\newcommand{\popsquiggle}[1]{%
  \tikz[baseline=-0.4ex]{
    \draw[#1,line width=2.6pt,line cap=round]
      plot [smooth] coordinates {(0,0.09) (0.34,0.01) (0.68,0.21) (1.02,0.01) (1.36,0.10)};
  }%
}
% Mot-clé surligné (marqueur orange sous le texte)
\newcommand{\cle}[1]{{\popxbold\color{popdark}#1}}

% ============================================================
% En-tête / pied
% ============================================================
\pagestyle{fancy}
\fancyhf{}
\renewcommand{\headrulewidth}{0pt}
\renewcommand{\footrulewidth}{0pt}
\lhead{\raisebox{-0.15ex}{\poplogo\fontsize{13}{13}\selectfont\color{popink}OnBuch\color{poporange}+}}
\rhead{\poppill{\DOCMATIERE\ \textbullet\ \DOCNIVEAU\ \textbullet\ Chapter \DOCLECON}{white}{popink}}
\lfoot{\popxbold\fontsize{7.6}{7.6}\selectfont\color{popmuted}\MakeUppercase{Signed course OnBuch+}\ \textbullet\ {\popsemi\DOCTITRE}}
\rfoot{\tikz[baseline=-0.6ex]\node[rounded corners=8.5pt,fill=popink,minimum height=17pt,minimum width=17pt,inner xsep=5pt,inner sep=0]{\popxbold\fontsize{8}{8}\selectfont\color{popcream}\thepage\,/\,\pageref*{LastPage}};}

% ============================================================
% Titres
% ============================================================
\newcommand{\chaptitre}[2]{%
  \begin{tcolorbox}[enhanced,colback=poporange,colframe=popink,boxrule=2.6pt,arc=11pt,
      drop shadow=popink,left=15pt,right=15pt,top=12pt,bottom=11pt,before skip=3pt,after skip=18pt]
    \poppill{#2}{popink}{popcream}\par\vspace{9pt}
    {\raggedright\hyphenpenalty=10000\exhyphenpenalty=10000\popdisplay\fontsize{22}{24}\selectfont\color{popcream}\MakeUppercase{#1}\par}
  \end{tcolorbox}
}
% Section numérotée : pastille orange avec le numéro + titre Archivo
\newcounter{popsec}
\newcommand{\coursec}[1]{%
  \stepcounter{popsec}\setcounter{popsub}{0}%
  \par\vspace{16pt}\noindent
  \begin{minipage}{\linewidth}
  \tikz[baseline=-0.7ex]\node[draw=popink,line width=1.6pt,fill=poporange,rounded corners=4pt,minimum size=22pt,inner sep=0]{\popdisplay\fontsize{12}{12}\selectfont\color{popcream}\thepopsec};%
  \hspace{8pt}{\popdisplay\fontsize{15}{17}\selectfont\color{popink}\MakeUppercase{#1}}\par
  \vspace{4pt}\noindent{\color{poporange}\rule{36pt}{3.6pt}}
  \end{minipage}\par\nopagebreak\vspace{8pt}
}
\newcounter{popsub}
\newcommand{\courssub}[1]{%
  \stepcounter{popsub}%
  \par\vspace{9pt}\noindent{\popxbold\fontsize{11.5}{13}\selectfont\color{popink}{\color{poporange}\thepopsec.\thepopsub}\hspace{6pt}#1}\par\nopagebreak\vspace{4pt}
}

% ============================================================
% Boîtes pédagogiques — même recette : fond blanc, bordure épaisse colorée,
% ombre dure encre, badge plein. Titre optionnel [#1].
% ============================================================
\tcbset{popbase/.style={breakable,enhanced,colback=white,boxrule=2.3pt,arc=9pt,
  shadow={3pt}{-3pt}{0pt}{popink},left=14pt,right=14pt,top=11pt,bottom=11pt,before skip=11pt,after skip=15pt,
  fontupper=\popsemi\fontsize{10.6}{14.5}\selectfont\color{popink},
  fontlower=\popsemi\fontsize{10.6}{14.5}\selectfont\color{popink}}}
% Coche dessinée (le glyphe ✓ n'existe pas dans les polices du document)
\newcommand{\popcheck}[1]{\tikz[baseline=-0.5ex]\draw[#1,line width=1.6pt,line cap=round,line join=round](-0.13,0.0)--(-0.04,-0.09)--(0.13,0.1);}
\providecommand{\coche}{\popcheck{popgreen}}
\providecommand{\resultat}[1]{\par\smallskip\noindent\fcolorbox{popgreen}{popgreenL}{\parbox{\dimexpr\linewidth-2\fboxsep-2\fboxrule\relax}{\textbf{Result.} #1}}\par\smallskip}
\newcommand{\popboxhead}[3]{\poppill{#1}{#2}{white}\ifx\relax#3\relax\else\hspace{6pt}{\popxbold\fontsize{10.6}{12}\selectfont\color{popink}#3}\fi\par\vspace{7pt}}

\newtcolorbox{aretenir}[1][]{popbase,colframe=popgreen,before upper={\popboxhead{Key points}{popgreen}{#1}}}
\newtcolorbox{exemplebox}[1][]{popbase,colframe=poporange,before upper={\popboxhead{Exemple}{poporange}{#1}}}
\newtcolorbox{definition}[1][]{popbase,colframe=popblue,before upper={\popboxhead{Definition}{popblue}{#1}}}
\newtcolorbox{propriete}[1][]{popbase,colframe=poppurple,before upper={\popboxhead{Property}{poppurple}{#1}}}
\newtcolorbox{methode}[1][]{popbase,colframe=popgold,before upper={\popboxhead{Method}{popgold}{#1}}}
\newtcolorbox{attention}[1][]{popbase,colframe=poppink,before upper={\popboxhead{Warning}{poppink}{#1}}}
\newtcolorbox{experience}[1][]{popbase,colframe=popblue,colback=popblueL,before upper={\popboxhead{Experiment}{popblue}{#1}}}
% « Le savais-tu ? » : carte violette pleine (contexte, culture, Cameroun)
\newtcolorbox{savaistu}[1][]{popbase,colback=poppurple,colframe=popink,
  fontupper=\popsemi\fontsize{10.4}{14.2}\selectfont\color{white},
  before upper={\poppill{Did you know?}{white}{poppurple}\ifx\relax#1\relax\else\hspace{6pt}{\popxbold\color{white}#1}\fi\par\vspace{7pt}}}
% Exercice résolu : énoncé en haut, \tcblower puis la solution sur fond crème
\newcounter{popexo}\newcounter{popexr}
\newtcolorbox{exoresolu}[1][]{popbase,colframe=poporange,colbacklower=popcream,
  segmentation style={popink,line width=1.2pt,dashed},
  before upper={\stepcounter{popexr}\popboxhead{Worked example \thepopexr}{poporange}{#1}},
  before lower={\poppill{Solution}{popink}{popcream}\par\vspace{6pt}}}
% Exercice (non résolu en place) — la correction est dans la partie Corrigés
\newtcolorbox{exercice}[1][]{popbase,colframe=popink,
  before upper={\stepcounter{popexo}\popboxhead{Exercise \thepopexo}{popink}{#1}}}
% Corrigé
\newtcolorbox{corrige}[1][]{popbase,colframe=popgreen,colback=popgreenL,
  before upper={\popboxhead{Solution}{popgreen}{#1}}}
% Objectifs / prérequis / bilan
\newtcolorbox{objectifs}{popbase,colframe=popink,colback=poporangeL,
  before upper={\popboxhead{Chapter objectives}{popink}{}}}
\newtcolorbox{prerequis}{popbase,colframe=popink,colback=popblueL,
  before upper={\popboxhead{Prerequisites}{popblue}{}}}
\newtcolorbox{bilan}{popbase,colframe=popink,colback=white,boxrule=2.8pt,
  before upper={\popboxhead{Fiche bilan}{poporange}{L'essentiel à connaître pour l'examen}}}
% Figure encadrée avec légende
\newtcolorbox{popfigure}[1][]{enhanced,breakable=false,colback=white,colframe=popline,boxrule=1.4pt,arc=8pt,
  left=10pt,right=10pt,top=10pt,bottom=8pt,before skip=10pt,after skip=12pt,halign=center,
  lower separated=false,halign lower=center,
  fontlower=\popsemi\fontsize{9}{11.5}\selectfont\color{popmuted},
  #1}
\newcommand{\legende}[1]{\par\vspace{6pt}{\popsemi\fontsize{9}{11.5}\selectfont\color{popmuted}#1}}

% ============================================================
% Signature OnBuch+
% ============================================================
\newcommand{\poplogoinline}[1]{{\poplogo\fontsize{#1}{#1}\selectfont\color{popink}OnBuch\color{poporange}+}}
\newcommand{\signatureonbuch}{%
  \begin{tcolorbox}[enhanced,colback=popink,colframe=popink,boxrule=2.2pt,arc=10pt,
      drop shadow=poporange,left=14pt,right=14pt,top=10pt,bottom=10pt,before skip=0pt,after skip=0pt]
    \tikz[baseline=-0.7ex]\node[circle,fill=popgreen,draw=popcream,line width=1.3pt,minimum size=22pt,inner sep=0]{\popcheck{white}};%
    \hspace{9pt}\begin{minipage}[c]{0.62\linewidth}
      {\popxbold\fontsize{12}{13}\selectfont\color{popcream}Signed\ \textbullet\ {\poplogo OnBuch\color{poporange}+}}\par\vspace{2pt}
      {\raggedright\popsemi\fontsize{8}{10.5}\selectfont\color{popline}\MakeUppercase{OnBuch+ teaching team}\\\MakeUppercase{Follows the official MINESEC syllabus}\par}
    \end{minipage}\hfill
    {\poplogo\fontsize{22}{22}\selectfont\color{popcream}OnBuch\color{poporange}+}
  \end{tcolorbox}%
}

% ============================================================
% Page de garde
% #1 titre · #2 matière · #3 niveau · #4 module · #5 n° leçon · #6 durée ·
% #7 description / objectifs (texte ou itemize)
% ============================================================
\newcommand{\pagegarde}[7]{%
  \thispagestyle{empty}
  \newgeometry{a4paper,margin=1.7cm,top=1.6cm,bottom=1.4cm}%
  \begin{tikzpicture}[remember picture,overlay]
    \fill[poporange!18] ([xshift=-3.2cm,yshift=-3.6cm]current page.north east) circle (4.6cm);
    \fill[poppurple!14] ([xshift=2.4cm,yshift=7.6cm]current page.south west) circle (3.8cm);
  \end{tikzpicture}%
  {\poplogo\fontsize{30}{30}\selectfont\color{popink}OnBuch\color{poporange}+}\hfill
  \raisebox{0.6ex}{\poppill{Signed course \textbullet\ Premium}{popink}{popcream}}\par
  \vspace{1pt}\popsquiggle{poppurple}\par
  \vspace{8pt}
  \begin{tcolorbox}[enhanced,colback=poporange,colframe=popink,boxrule=2.8pt,arc=13pt,
      drop shadow=popink,left=18pt,right=18pt,top=12pt,bottom=13pt,after skip=10pt]
    \poppill{#2 \textbullet\ #3}{popink}{popcream}\hspace{5pt}\poppill{#4}{white}{popink}\par\vspace{8pt}
    {\popdisplay\fontsize{14}{14}\selectfont\color{popink}CHAPTER #5}\par\vspace{4pt}
    {\raggedright\hyphenpenalty=10000\exhyphenpenalty=10000\popdisplay\fontsize{25}{27}\selectfont\color{popcream}\MakeUppercase{#1}\par}
    \vspace{8pt}
    {\popxbold\fontsize{9.5}{9.5}\selectfont\color{popink}\MakeUppercase{Suggested time: #6}}
  \end{tcolorbox}
  \begin{tcolorbox}[enhanced,colback=white,colframe=popink,boxrule=2.2pt,arc=10pt,
      drop shadow=popink,left=16pt,right=16pt,top=10pt,bottom=10pt,after skip=8pt]
    \poppill{In this chapter}{poppurple}{white}\par\vspace{5pt}
    \popsemi\fontsize{9.3}{12.6}\selectfont\color{popink}#7
  \end{tcolorbox}
  \noindent\begin{minipage}[t]{0.487\linewidth}
    \begin{tcolorbox}[enhanced,colback=popgreen,colframe=popink,boxrule=2.2pt,arc=10pt,
        drop shadow=popink,left=11pt,right=11pt,top=10pt,bottom=10pt,height=3.1cm,valign=center]
      \tcbox[colback=white,colframe=popink,boxrule=1.3pt,arc=5pt,boxsep=0pt,nobeforeafter,
        left=3pt,right=3pt,top=3pt,bottom=3pt,valign=center]{\qrcode[height=1.9cm]{https://play.google.com/store/apps/details?id=com.luvvix.onbuch&hl=en}}%
      \hspace{8pt}\begin{minipage}[c]{0.5\linewidth}
        {\popdisplay\fontsize{11}{12.5}\selectfont\color{white}\MakeUppercase{Download OnBuch}}\par\vspace{4pt}
        {\raggedright\popsemi\fontsize{8.4}{11}\selectfont\color{white}Free \textbullet\ Google Play\\AI tutor, past papers, results\par}
      \end{minipage}
    \end{tcolorbox}
  \end{minipage}\hfill
  \begin{minipage}[t]{0.487\linewidth}
    \begin{tcolorbox}[enhanced,colback=poppurple,colframe=popink,boxrule=2.2pt,arc=10pt,
        drop shadow=popink,left=11pt,right=11pt,top=10pt,bottom=10pt,height=3.1cm,valign=center]
      \tikz[baseline=-0.7ex]\node[circle,fill=white,draw=popink,line width=1.3pt,minimum size=1.6cm,inner sep=0]{\popdisplay\fontsize{24}{24}\selectfont\color{poppurple}+};%
      \hspace{8pt}\begin{minipage}[c]{0.6\linewidth}
        {\popdisplay\fontsize{11}{12.5}\selectfont\color{white}\MakeUppercase{Go OnBuch+}}\par\vspace{4pt}
        {\raggedright\popsemi\fontsize{8.4}{11}\selectfont\color{white}All signed courses, unlimited AI tutor, PDF sheets — OnBuch+ tab in the app\par}
      \end{minipage}
    \end{tcolorbox}
  \end{minipage}
  \vspace{8pt}
  \popcontenu
  \vfill
  \signatureonbuch
  \restoregeometry
  \newpage
}

% Rangée « Ce cours contient » (page de garde)
\newcommand{\poptile}[3]{% #1 couleur #2 gros texte #3 légende
  \begin{tcolorbox}[enhanced,colback=#1,colframe=popink,boxrule=2pt,arc=9pt,shadow={2.5pt}{-2.5pt}{0pt}{popink},
      left=6pt,right=6pt,top=9pt,bottom=9pt,halign=center,valign=center,height=1.75cm,width=\linewidth,nobeforeafter]
    {\popdisplay\fontsize{12.5}{13}\selectfont\color{white}\MakeUppercase{#2}}\par\vspace{3pt}
    {\centering\popsemi\fontsize{7.8}{9.5}\selectfont\color{white}#3\par}
  \end{tcolorbox}}
\newcommand{\popcontenu}{%
  \par\noindent{\popxboldls\fontsize{7.5}{7.5}\selectfont\color{popmuted}THIS SIGNED COURSE CONTAINS}\par\vspace{7pt}
  \noindent\begin{minipage}[t]{0.235\linewidth}\poptile{popblue}{Course}{complete, illustrated, step by step}\end{minipage}\hfill
  \begin{minipage}[t]{0.235\linewidth}\poptile{poporange}{Methods}{and worked examples}\end{minipage}\hfill
  \begin{minipage}[t]{0.235\linewidth}\poptile{poppink}{Exercises}{exam style + solutions}\end{minipage}\hfill
  \begin{minipage}[t]{0.235\linewidth}\poptile{popgreen}{Summary sheet}{the essentials to remember}\end{minipage}\par}

% Sommaire
\newtcolorbox{sommaire}{enhanced,colback=white,colframe=popink,boxrule=2.2pt,arc=10pt,
  drop shadow=popink,left=18pt,right=18pt,top=13pt,bottom=13pt,before skip=6pt,after skip=20pt,
  before upper={\poppill{Contents}{poppurple}{white}\par\vspace{9pt}\popsemi\fontsize{10.6}{14.5}\selectfont\color{popink}}}

% Signature de fin de cours — poussée tout en bas de la dernière page
\newcommand{\findecours}{%
  \par\vfill\nopagebreak
  \begin{center}\popsquiggle{poporange}\end{center}\vspace{4pt}
  \signatureonbuch
}
\AtBeginDocument{\def\llbracket{\mathopen{[\mkern-3.5mu[}}\def\rrbracket{\mathclose{]\mkern-3.5mu]}}} % intervalles d'entiers (glyphe ⟦ absent de la police)
% Glyphes absents de Fira Math : redéfinis à partir de glyphes disponibles
\AtBeginDocument{%
  \def\setminus{\mathbin{\backslash}}\let\smallsetminus\setminus
  \def\vdots{\mathord{\vbox{\baselineskip3.2pt\lineskiplimit0pt\kern4pt\hbox{.}\hbox{.}\hbox{.}}}}%
  \def\ddots{\mathinner{\mkern1mu\raise6.5pt\hbox{.}\mkern2mu\raise3.6pt\hbox{.}\mkern2mu\raise0.7pt\hbox{.}\mkern1mu}}%
  \def\triangle{\mathord{\scalebox{0.9}{$\symup{\Delta}$}}}%
  \def\nabla{\mathord{\raisebox{\depth}{\scalebox{1}[-1]{$\symup{\Delta}$}}}}%
  \def\checkmark{\popcheck{popdark}}%
  \def\star{\mathbin{\ast}}\def\bigstar{\mathord{\ast}}%
  \def\diamond{\mathbin{\scalebox{0.6}{$\Diamond$}}}%
  \def\bigcup{\mathop{\mathchoice{\raisebox{-0.3em}{\scalebox{1.7}{$\cup$}}}{\scalebox{1.25}{$\cup$}}{\cup}{\cup}}\displaylimits}%
  \def\bigcap{\mathop{\mathchoice{\raisebox{-0.3em}{\scalebox{1.7}{$\cap$}}}{\scalebox{1.25}{$\cap$}}{\cap}{\cap}}\displaylimits}%
  \def\lesseqqgtr{\mathrel{\lessgtr}}\def\bigoplus{\mathop{\oplus}\displaylimits}%
}
\providecommand{\ssi}{if and only if} % abréviation fréquente
\AtBeginDocument{\providecommand{\wideparen}{\overparen}} % arc de cercle

% Fonctions usuelles que les modèles utilisent sans que amsmath les fournisse
\providecommand{\arsinh}{\operatorname{arsinh}}\providecommand{\arcosh}{\operatorname{arcosh}}
\providecommand{\artanh}{\operatorname{artanh}}\providecommand{\arsech}{\operatorname{arsech}}
\providecommand{\arcsch}{\operatorname{arcsch}}\providecommand{\arcoth}{\operatorname{arcoth}}
\providecommand{\arcsinh}{\operatorname{arsinh}}\providecommand{\arccosh}{\operatorname{arcosh}}
\providecommand{\arctanh}{\operatorname{artanh}}\providecommand{\sech}{\operatorname{sech}}
\providecommand{\csch}{\operatorname{csch}}\providecommand{\arcsec}{\operatorname{arcsec}}
\providecommand{\arccsc}{\operatorname{arccsc}}\providecommand{\arccot}{\operatorname{arccot}}
\providecommand{\sgn}{\operatorname{sgn}}\providecommand{\lcm}{\operatorname{lcm}}
\providecommand{\Var}{\operatorname{Var}}\providecommand{\Cov}{\operatorname{Cov}}
\providecommand{\permil}{\ensuremath{{}^{0}\!/_{00}}}\providecommand{\textperthousand}{\ensuremath{{}^{0}\!/_{00}}}

%% FIGFIX-BEGIN v5 — correctifs globaux des figures (docs/onbuch-generation/tools/figfix)
\usepackage{adjustbox}\usepackage{etoolbox}\tcbuselibrary{raster}
% toute figure TikZ/pgfplots plus large que la ligne est réduite à la largeur disponible
\BeforeBeginEnvironment{tikzpicture}{\begin{adjustbox}{max width=\linewidth}}
\AfterEndEnvironment{tikzpicture}{\end{adjustbox}}
% légendes pgfplots lisibles par-dessus les courbes
\pgfplotsset{every axis/.append style={legend style={fill=white,fill opacity=0.92,text opacity=1,draw=black!15,inner sep=3pt}}}
% étiquettes d'arêtes (syntaxe edge["..."]) avec fond blanc
\tikzset{every edge quotes/.append style={fill=white,inner sep=1.2pt}}
%% FIGFIX-END
\begin{document}

\pagegarde{ECONOMIC GEOLOGY: The Soil}{Geology}{Upper Sixth}{Upper Sixth — Geology}{4}{7 periods}{This chapter studies the soil as a product of rock weathering and biological activity, describing its formation, its layered profile, its physical and chemical characteristics, and its economic importance. It ends with the sustainable management and protection of soils, a major environmental and agricultural issue in Cameroon.
\vspace{4pt}
\begin{multicols}{2}\small
\begin{itemize}
\item By the end of this chapter I can define soil and distinguish it from the parent rock.
\item By the end of this chapter I can explain the roles of climate, organisms, relief, parent material and time in soil formation.
\item By the end of this chapter I can describe and draw a labelled soil profile with its horizons (O, A, E, B, C, R).
\item By the end of this chapter I can identify the main types of soil layers and relate them to Cameroonian environments.
\item By the end of this chapter I can describe the physical characteristics of soil (texture, structure, porosity, colour, depth).
\item By the end of this chapter I can describe the chemical and biological characteristics of soil (pH, humus, mineral nutrients, soil organisms).
\end{itemize}\end{multicols}}

\begin{sommaire}
\begin{multicols}{2}
\begin{enumerate}
\item Soil Formation
\item The Soil Profile and Types of Soil Layers
\item Physical Characteristics of Soil
\item Chemical and Biological Characteristics of Soil
\item Soil Usage, Sustainable Management and Protection
\item Integration activity
\item Methods and worked examples
\item Exercises
\item Solutions to the exercises
\item Summary sheet
\end{enumerate}
\end{multicols}
\end{sommaire}

\begin{objectifs}
\begin{itemize}
\item By the end of this chapter I can define soil and distinguish it from the parent rock.
\item By the end of this chapter I can explain the roles of climate, organisms, relief, parent material and time in soil formation.
\item By the end of this chapter I can describe and draw a labelled soil profile with its horizons (O, A, E, B, C, R).
\item By the end of this chapter I can identify the main types of soil layers and relate them to Cameroonian environments.
\item By the end of this chapter I can describe the physical characteristics of soil (texture, structure, porosity, colour, depth).
\item By the end of this chapter I can describe the chemical and biological characteristics of soil (pH, humus, mineral nutrients, soil organisms).
\item By the end of this chapter I can explain the agricultural and economic uses of soil in Cameroon.
\item By the end of this chapter I can propose methods of sustainable soil management and soil protection against erosion and degradation.
\item By the end of this chapter I can interpret a soil profile diagram or data table in a GCE-style structured question.
\end{itemize}
\end{objectifs}

\begin{prerequis}
\begin{itemize}
\item Rocks and the rock cycle (igneous, sedimentary, metamorphic rocks) from Lower Sixth.
\item Weathering and erosion as external geological processes.
\item Basic chemistry: pH, ions, and simple chemical reactions.
\item Map reading and relief interpretation (slopes, plains, plateaus).
\item General knowledge of Cameroonian physical geography (climate zones, vegetation, relief).
\end{itemize}
\end{prerequis}

\newpage

\coursec{Soil Formation}

Every rainy season in Bamenda, farmers watch helplessly as red water carries the topsoil of their maize fields down the slopes, while in the lowlands of Ngaoundéré, soils crack and harden before the planting season even ends. Yet a few kilometres away, other farmers on the same slopes harvest bumper crops year after year. What makes one soil fertile and another exhausted? Why do the volcanic soils of Mount Cameroon produce some of the best bananas in the world? To answer these questions, we must first understand what soil really is and how it comes into existence. This section takes you beneath your feet: we shall define soil rigorously, distinguish it from regolith and bedrock, study the five factors that control its formation, and follow the stages by which bare rock is slowly transformed into a living, fertile soil.

\courssub{What is Soil?}

\textbf{An observation to start with.}\par
Dig a small pit in a school garden in Buea. The uppermost material is dark, soft, full of rootlets and tiny living creatures, and it smells of decaying leaves. A few tens of centimetres deeper, the material becomes lighter in colour, harder, and contains broken fragments of rock. Deeper still, you strike solid, unbroken rock. Only the upper, dark, living layer deserves the name \cle{soil}. The rest is weathered rock debris and solid bedrock. This simple observation contains the whole definition of soil.

\begin{definition}[Soil]
\cle{Soil} is the thin, loose surface layer of the Earth's crust produced by the \textbf{weathering of rocks} and the \textbf{decomposition of organic matter}, which is capable of \textbf{supporting plant life}. Soil is therefore a \emph{natural body}, a mixture of mineral particles, organic matter (humus), water, air, and living organisms.
\end{definition}

\begin{definition}[Parent material and bedrock]
The \cle{parent material} (or parent rock) is the geological material — solid \textbf{bedrock} or loose transported deposits — from which a soil develops through weathering and the addition of organic matter. \cle{Bedrock} is the solid, unweathered rock found beneath the soil and the weathered debris.
\end{definition}

\begin{attention}[Soil is not just sand or sediment!]
A very common mistake at the GCE is to call any loose material at the Earth's surface a ``soil''. This is wrong. Beach sand at Limbe or river sediment along the Sanaga is loose material, but it is \emph{not} a true soil, because it lacks \textbf{organic matter (humus)}, has \textbf{no organised layers (horizons)}, and \textbf{cannot support plant life} on its own. A true soil must (i) result from weathering \emph{plus} biological activity, (ii) contain humus, and (iii) be able to support plants.
\end{attention}

\textbf{Soil, regolith and bedrock: three terms to distinguish.}\par
Geologists distinguish three levels at the Earth's surface:

\begin{center}
\begin{tabularx}{\linewidth}{|l|X|X|}
\hline
\rowcolor{popblueL} \textbf{Term} & \textbf{Description} & \textbf{Supports plants?} \\
\hline
Soil & Upper weathered layer enriched in humus and organised into horizons & Yes \\
\hline
Regolith & Loose, fragmented, weathered rock material (including soil) with little or no organic content & Poorly or not at all \\
\hline
Bedrock & Solid, unweathered rock beneath the soil and regolith & No \\
\hline
\end{tabularx}
\end{center}

The \cle{regolith} is thus the general name for \emph{all} loose material covering the bedrock; soil is only its uppermost, biologically active part. The relationship can be summarised as:
\[
\text{bedrock} \xrightarrow{\ \text{weathering}\ } \text{regolith} \xrightarrow{\ +\ \text{organic matter, time}\ } \text{soil}.
\]

\courssub{The Five Soil-Forming Factors}

\textbf{Why do soils differ from place to place?}\par
The red lateritic soils of Ebolowa, the black volcanic soils of the slopes of Mount Cameroon, and the sandy soils of the Far North near Maroua are all soils, yet they look and behave completely differently. The Russian-American soil scientist Hans Jenny summarised the controls on soil formation in five factors, easily remembered with the mnemonic \cle{CLORPT}: \textbf{CL}imate, \textbf{O}rganisms, \textbf{R}elief, \textbf{P}arent material, \textbf{T}ime.

\begin{methode}[The CLORPT mnemonic]
To remember and present the five soil-forming factors in an examination:
\begin{enumerate}
\item \textbf{CL} — \textbf{Climate}: rainfall and temperature control weathering and leaching.
\item \textbf{O} — \textbf{Organisms}: plants, animals and microbes add and mix organic matter.
\item \textbf{R} — \textbf{Relief}: slope and elevation control drainage and erosion.
\item \textbf{P} — \textbf{Parent material}: the rock or deposit from which the soil forms.
\item \textbf{T} — \textbf{Time}: soil formation is a slow process needing hundreds to thousands of years.
\end{enumerate}
In an essay, define each factor in one sentence and illustrate it with a Cameroonian example.
\end{methode}

\begin{popfigure}
\begin{center}
\begin{tikzpicture}[font=\small,
factor/.style={draw=popblue, fill=popblueL, rounded corners, minimum width=2.7cm, minimum height=0.9cm, align=center},
rock/.style={draw=popdark, fill=gray!25, rounded corners, minimum width=2.6cm, minimum height=1.1cm, align=center},
soilbox/.style={draw=popgreen, fill=popgreenL, rounded corners, minimum width=2.6cm, minimum height=1.1cm, align=center},
arrow/.style={-latex, thick, popblue}]
\node[factor] (cl) at (0,3.2) {Climate};
\node[factor] (org) at (3.4,3.2) {Organisms};
\node[factor] (rel) at (6.8,3.2) {Relief};
\node[factor] (par) at (10.2,3.2) {Parent material};
\node[factor] (tim) at (13.6,3.2) {Time};
\node[rock, align=center] (rk) at (3.4,0){Parent rock\\(bedrock)};
\node[soilbox, align=center] (so) at (10.2,0){Mature soil\\(with horizons)};
\draw[arrow] (cl.south) -- ++(0,-0.5) -| (rk.north);
\draw[arrow] (org.south) -- ++(0,-0.5) -| (rk.north);
\draw[arrow] (rel.south) -- ++(0,-0.5) -| (rk.north);
\draw[arrow] (par.south) -- ++(0,-0.5) -| (so.north);
\draw[arrow] (tim.south) -- ++(0,-0.5) -| (so.north);
\draw[-latex, very thick, poporange] (rk.east) -- node[above, align=center]{weathering\\+ humification} (so.west);
\end{tikzpicture}
\end{center}
\legende{The five soil-forming factors (CLORPT) acting on parent rock to produce a mature soil.}
\end{popfigure}

\courssub{Role of Parent Material}

The parent material is the ``raw ingredient'' of the soil: it determines the \textbf{mineral composition}, the \textbf{texture} (proportions of sand, silt and clay) and often the \textbf{colour} of the resulting soil. Two soils formed under the same climate but on different rocks will be quite different.

\begin{exemplebox}[Parent rocks and their soils in Cameroon]
\begin{itemize}
\item \textbf{Granite} (Adamawa Plateau): rich in quartz, which resists weathering, granite gives \textbf{sandy, well-drained, often acidic and rather poor soils}.
\item \textbf{Basalt} (Mount Cameroon, Cameroon Volcanic Line): rich in iron, magnesium and calcium minerals that weather easily, basalt gives \textbf{dark, clayey, very fertile soils} — the famous volcanic soils that produce bananas, tea and palm oil around Buea and Limbe.
\item \textbf{Limestone}: dissolves by chemical weathering, leaving little residue, so it gives \textbf{shallow, stony, calcareous soils} rich in calcium.
\end{itemize}
\end{exemplebox}

\courssub{Role of Climate}

Climate acts through \textbf{rainfall} and \textbf{temperature}. Heat accelerates chemical reactions: as a useful rule, chemical weathering roughly doubles in speed for every rise of about \SI{10}{\ensuremath{{}^{\circ}\mathrm{C}}}. Abundant rainfall supplies the water needed for \textbf{hydrolysis} and \textbf{dissolution}, and it washes (leaches) soluble minerals downwards.

\begin{itemize}
\item In the \textbf{humid tropics} of southern Cameroon (Ebolowa, Kribi, Yaoundé), high temperatures (about \SI{25}{\ensuremath{{}^{\circ}\mathrm{C}}}) and heavy rainfall (often above \SI{1500}{mm} per year) cause intense weathering and strong \cle{leaching}: soluble bases (calcium, magnesium, potassium, sodium) are washed away, leaving behind iron and aluminium oxides. The result is the deep, red, \textbf{lateritic soils} (ferralsols) typical of the South.
\item In the \textbf{drier north} (Maroua region), lower rainfall means weaker leaching; soils are thinner, less weathered, and may harden and crack in the dry season.
\end{itemize}

\begin{savaistu}[Why are lateritic soils red?]
The red colour of the soils of southern Cameroon comes from \textbf{iron oxides} (such as haematite) that remain after intense leaching has removed most other minerals. The same iron-rich material, when hardened, forms \emph{laterite}, traditionally cut into blocks for building houses.
\end{savaistu}

\courssub{Role of Organisms}

Without living organisms there would be no soil — only regolith. Organisms contribute in two essential ways:

\begin{enumerate}
\item \textbf{They add organic matter.} Dead leaves, roots and animal remains are decomposed by \textbf{bacteria and fungi} into \cle{humus}, the dark, stable organic material that gives topsoil its colour, fertility and ability to retain water.
\item \textbf{They mix and aerate the soil.} \textbf{Plant roots} penetrate and break up the material; \textbf{earthworms} swallow soil and deposit it as fertile casts; \textbf{termites} — extremely abundant in Cameroonian savannas — transport mineral particles from depth to the surface, building mounds and improving aeration and drainage.
\end{enumerate}

\courssub{Role of Relief and Time}

\textbf{Relief (topography)} controls how water moves over and through the ground, and therefore how soil is eroded or accumulated:

\begin{itemize}
\item On \textbf{steep slopes} (for example the flanks of Mount Cameroon or the Bamenda escarpment), rainwater runs off rapidly, eroding material as fast as it forms: soils remain \textbf{thin and immature}.
\item On \textbf{flat or gently undulating land} (plateaus, valley bottoms), water infiltrates and weathered material accumulates: soils become \textbf{deep and mature}.
\item \textbf{Slope aspect} also matters: slopes facing the sun are warmer and drier than shaded slopes, which affects weathering and vegetation.
\end{itemize}

\textbf{Time} is the silent factor. Weathering and humus accumulation are extremely slow: the formation of just a few centimetres of topsoil may require \textbf{hundreds to thousands of years}. A \cle{young soil} shows little layering and resembles its parent rock; a \cle{mature soil} displays well-developed horizons and characteristics dominated by climate and organisms rather than by the parent rock. This is why the loss of topsoil by erosion, as on the Bamenda slopes, is practically irreversible on a human timescale.

\courssub{Stages of Soil Formation}

Soil formation is a continuous process, but for clarity it can be divided into four stages:

\begin{enumerate}
\item \textbf{Weathering of the bedrock.} Physical disintegration (temperature changes, freeze-thaw in highlands, root pressure) and chemical decomposition (hydrolysis, oxidation, dissolution) break the solid rock into loose regolith.
\item \textbf{Accumulation of organic matter.} Pioneer plants (lichens, mosses, grasses) colonise the regolith; their remains, decomposed by microbes, form the first humus, which darkens the upper layer.
\item \textbf{Development of horizons.} Water moving downwards transports clay and dissolved substances, so distinct layers (\cle{horizons}) appear: an organic-rich topsoil, a leached zone, and an accumulation zone above the weathered parent material.
\item \textbf{Mature soil.} After hundreds to thousands of years, a deep, well-layered soil in equilibrium with its climate and vegetation is established.
\end{enumerate}

\begin{popfigure}
\begin{center}
\begin{tikzpicture}[font=\small]
% Stage 1
\begin{scope}[shift={(0,0)}]
\fill[gray!45] (0,0) rectangle (2.6,1.4);
\node at (1.3,0.7) {solid bedrock};
\node[below] at (1.3,0) {\textbf{1. Bare rock}};
\end{scope}
% Stage 2
\begin{scope}[shift={(3.6,0)}]
\fill[gray!45] (0,0) rectangle (2.6,1.0);
\fill[gray!20] (0,1.0) rectangle (2.6,1.4);
\draw[popgreen, thick] (0.5,1.4) -- (0.5,1.8);
\draw[popgreen, thick] (1.3,1.4) -- (1.3,1.85);
\draw[popgreen, thick] (2.1,1.4) -- (2.1,1.8);
\node at (1.3,0.5) {bedrock};
\node at (1.3,1.2) {\footnotesize regolith};
\node[below] at (1.3,0) {\textbf{2. Weathering}};
\end{scope}
% Stage 3
\begin{scope}[shift={(7.2,0)}]
\fill[gray!45] (0,0) rectangle (2.6,0.7);
\fill[gray!20] (0,0.7) rectangle (2.6,1.1);
\fill[popgold!50] (0,1.1) rectangle (2.6,1.4);
\draw[popgreen, thick] (0.5,1.4) -- (0.5,1.8);
\draw[popgreen, thick] (1.3,1.4) -- (1.3,1.85);
\draw[popgreen, thick] (2.1,1.4) -- (2.1,1.8);
\node at (1.3,0.35) {\footnotesize bedrock};
\node at (1.3,0.9) {\footnotesize regolith};
\node at (1.3,1.25) {\footnotesize humus};
\node[below] at (1.3,0) {\textbf{3. Organic matter}};
\end{scope}
% Stage 4
\begin{scope}[shift={(10.8,0)}]
\fill[gray!45] (0,0) rectangle (2.6,0.5);
\fill[gray!20] (0,0.5) rectangle (2.6,0.9);
\fill[poporange!40] (0,0.9) rectangle (2.6,1.2);
\fill[popgold!60] (0,1.2) rectangle (2.6,1.5);
\draw[popgreen, thick] (0.5,1.5) -- (0.5,1.9);
\draw[popgreen, thick] (1.3,1.5) -- (1.3,1.95);
\draw[popgreen, thick] (2.1,1.5) -- (2.1,1.9);
\node at (1.3,0.25) {\footnotesize bedrock};
\node at (1.3,0.7) {\footnotesize parent mat.};
\node at (1.3,1.05) {\footnotesize subsoil};
\node at (1.3,1.35) {\footnotesize topsoil};
\node[below] at (1.3,0) {\textbf{4. Mature soil}};
\end{scope}
\draw[-latex, very thick, poporange] (2.7,0.8) -- (3.5,0.8);
\draw[-latex, very thick, poporange] (6.3,0.8) -- (7.1,0.8);
\draw[-latex, very thick, poporange] (9.9,0.8) -- (10.7,0.8);
\end{tikzpicture}
\end{center}
\legende{Schematic sequence of the four stages of soil formation, from bare bedrock to a mature, layered soil.}
\end{popfigure}

\begin{exoresolu}[Explaining the fertility of Mount Cameroon soils]
\textbf{Statement.} The slopes of Mount Cameroon support some of the most productive banana plantations in Africa. Using your knowledge of soil-forming factors, explain why the soils there are so fertile.
\tcblower
\textbf{Solution.} We apply the CLORPT factors one by one.
\begin{itemize}
\item \textbf{Parent material:} Mount Cameroon is a basaltic volcano on the Cameroon Volcanic Line. Basalt is rich in ferromagnesian minerals and calcium-rich feldspars, which weather readily and release plant nutrients (iron, magnesium, calcium, potassium). The resulting soil is dark, clayey and naturally fertile.
\item \textbf{Climate:} the mountain receives very heavy rainfall and enjoys warm temperatures, so chemical weathering is rapid and a thick soil mantle develops quickly on young lava flows and ash deposits.
\item \textbf{Organisms:} luxuriant natural vegetation and plantation cover return abundant organic matter, and decomposers maintain a humus-rich topsoil.
\item \textbf{Relief:} on gentle lower slopes and lava plateaus, deep soils accumulate; only on the steepest flanks are soils thin.
\item \textbf{Time:} repeated eruptions continually supply fresh ash and lava, which weather rapidly in the humid climate, constantly renewing soil fertility.
\end{itemize}
\textbf{Conclusion:} the exceptional fertility results from the combination of a nutrient-rich basaltic parent material with a hot, wet climate and abundant biological activity — a perfect illustration of the soil-forming factors.
\end{exoresolu}

\begin{attention}[Exam tip — one factor, one example]
GCE Advanced Level questions frequently ask candidates to \emph{``explain how ONE named factor influences soil formation''}. Do not list all five factors: select the factor asked for, explain the \textbf{mechanism} (for example, rainfall $\rightarrow$ leaching $\rightarrow$ loss of bases $\rightarrow$ lateritic soil), and always support your answer with a \textbf{Cameroonian example} (lateritic soils of Ebolowa for climate, volcanic soils of Mount Cameroon for parent material, thin soils on the Bamenda escarpment for relief).
\end{attention}

\begin{exercice}[Soil formation around a granite inselberg]
Near Ngaoundéré, a granite inselberg rises above a flat plain. (a) Compare the soils you would expect on the steep slopes of the inselberg and on the surrounding plain, justifying your answer with two soil-forming factors. (b) Predict the general texture of the soils derived from the granite.
\end{exercice}

\begin{corrige}[Exercise 1]
(a) \textbf{Relief:} on the steep slopes, rapid runoff erodes weathered material as fast as it forms, so soils are thin, stony and immature; on the flat plain, material and water accumulate, giving deep, mature soils. \textbf{Time:} because slope materials are constantly removed, slope soils never have time to develop horizons, whereas plain soils have remained in place long enough to become mature. (b) Granite is rich in quartz, a mineral very resistant to weathering, so the soils will be predominantly \textbf{sandy}, well drained and rather poor in nutrients.
\end{corrige}

\begin{aretenir}[Key points of this section]
\begin{itemize}
\item Soil is the thin, loose, \textbf{living} surface layer formed by rock weathering and organic decomposition; it must contain humus and support plant life — unlike regolith or sediment.
\item Five factors control soil formation, summarised by \textbf{CLORPT}: Climate, Organisms, Relief, Parent material, Time.
\item Parent material sets the mineral base: granite $\rightarrow$ sandy soils; basalt $\rightarrow$ dark fertile clayey soils; limestone $\rightarrow$ shallow calcareous soils.
\item Hot, wet climates (southern Cameroon) cause intense leaching and red lateritic soils.
\item Organisms produce humus and mix the soil; steep slopes keep soils thin while flat lands accumulate deep soils; maturity requires hundreds to thousands of years.
\item Soil forms in four stages: weathering of bedrock $\rightarrow$ accumulation of organic matter $\rightarrow$ development of horizons $\rightarrow$ mature soil.
\end{itemize}
\end{aretenir}

\coursec{The Soil Profile and Types of Soil Layers}

In the previous section we saw how soil is born: bedrock is attacked by weathering, broken into mineral fragments, mixed with organic matter from dead plants and animals, and slowly transformed by water, organisms and time. But soil is not a uniform mixture. If you dig a deep pit in a farm near Bafoussam or along the slopes of Mount Cameroon and look carefully at the wall of the pit, you will notice something striking: the soil is arranged in \emph{distinct horizontal bands}, each with its own colour, texture and composition. Dark at the top, reddish or yellowish in the middle, pale and stony near the bottom. This layered arrangement is not accidental --- it is the visible record of the soil-forming processes we studied earlier. Reading these layers is like reading the biography of the soil. In this section we learn to read it.

\begin{aretenir}[Learning objectives for this section]
\begin{itemize}
\item Define and distinguish \cle{soil profile} and \cle{soil horizon}.
\item Identify, describe and label the master horizons (O, A, E, B, C, R) and common subordinate distinctions.
\item Explain the processes of \cle{eluviation} and \cle{illuviation} and link them to the correct horizons.
\item Compare soil profiles from different climatic zones of Cameroon (humid forest vs. Sahelian north).
\item Interpret a soil profile diagram or photograph in an examination context.
\end{itemize}
\end{aretenir}

\courssub{The Soil Profile and Soil Horizons}

\begin{definition}[Soil profile and soil horizon]
A \cle{soil profile} is a vertical section cut through the soil, from the surface down to the unaltered bedrock, showing the succession of horizontal layers that make up the soil. Each distinguishable horizontal layer is called a \cle{soil horizon}. Horizons are distinguished from one another by differences in colour, texture, structure, consistence, and chemical or organic composition.
\end{definition}

The horizons are labelled with capital letters, from top to bottom: \cle{O}, \cle{A}, \cle{E}, \cle{B}, \cle{C} and \cle{R}. Not every soil shows all six horizons --- a young soil on a recent lava flow of Mount Cameroon may show only a thin A over C and R, while an ancient soil of the southern forest plateau may display the complete sequence. The combination of horizons present, their thickness, and the nature of the boundaries between them (abrupt, clear, gradual, or diffuse) defines the \emph{type} of the soil profile.

\begin{attention}[Common mistake]
A frequent error in examinations is to place the dark, humus-rich layer at the \emph{bottom} of the profile. Organic matter is produced by living plants and animals at the surface, so it is \textbf{always found at or near the top} of the profile (O and A horizons). The bottom of the profile is occupied by the parent material and bedrock, which contain almost no organic matter.
\end{attention}

\courssub{Description of the Master Horizons}

\textbf{The O horizon --- the organic layer}\par
At the very top lies the \cle{O horizon} (O for \emph{organic}). It consists of \emph{litter} (freshly fallen leaves, twigs, dead roots and animal remains) together with \emph{humus}, the dark, partly decomposed organic matter produced by the action of bacteria, fungi, earthworms and termites. In the humid forest zone of the South and Centre regions of Cameroon, the litter layer is renewed constantly because vegetation grows all year round, but it is often thin (1--5 cm) because decomposition is extremely rapid under hot, wet conditions. In cooler or waterlogged climates the O horizon may be much thicker (peat). Subdivisions: Oi (slightly decomposed), Oe (moderately decomposed), Oa (highly decomposed/humus).

\textbf{The A horizon --- the topsoil}\par
Beneath the O horizon lies the \cle{A horizon}, or \cle{topsoil}. It is a mineral layer, but it is dark brown or black because humus is intimately mixed with the mineral particles (sand, silt and clay). It is the horizon of greatest biological activity: roots, earthworms, termites and micro-organisms abound. It is also the most fertile layer and the one exploited by crops --- which is why farmers treasure it. At the same time, the A horizon is a zone of \emph{eluviation}: rainwater percolating downwards dissolves soluble minerals and carries fine clay particles away to deeper levels. This downward washing is called \cle{leaching}. Common subordinate distinctions: Ap (ploughed/disturbed), Ah (humus accumulation), Ahe (humus + eluviation).

\textbf{The E horizon --- the leached layer}\par
In some soils, especially under forest and in cool, humid climates, a distinct \cle{E horizon} (E for \emph{eluviation}) develops between A and B. It is a light-coloured --- grey, whitish or pale yellow --- layer from which iron oxides, aluminium compounds and clay have been almost completely washed out, leaving behind resistant minerals such as quartz. It is the hallmark of strongly leached soils called \emph{podzols}. In the deeply weathered tropical soils of Cameroon the E horizon is usually absent or poorly developed, because iron oxides are not removed but accumulate below. When present, it is often thin and may be designated Ae or E.

\textbf{The B horizon --- the subsoil}\par
The \cle{B horizon}, or \cle{subsoil}, is the zone of \emph{illuviation}: it receives and accumulates the materials washed down from A and E. Clay particles, iron and aluminium oxides, and sometimes calcium carbonate, gypsum or silica, are deposited here. Because it is enriched in iron oxides, the B horizon of tropical soils is typically \emph{red, reddish-brown or yellow} --- the famous red colour of the lateritic soils seen in road cuttings around Yaoundé and Ebolowa. The B horizon is denser, more compact and poorer in humus and roots than the topsoil. In lateritic profiles it may harden into an iron-rich crust (\emph{laterite} or ironpan) when exposed at the surface. Common subordinate distinctions: Bt (clay accumulation), Bw (weathering/structure development), Bs (sesquioxide accumulation), Bk (carbonate accumulation).

\textbf{The C horizon --- the weathered parent material}\par
Below the subsoil lies the \cle{C horizon}: the \emph{regolith}, made of partially weathered fragments of the parent rock. The original rock structures may still be visible, but the rock is softened, discoloured and crumbly. There is practically no organic matter here. In Cameroon, the C horizon over granite or gneiss of the Congo craton basement often appears as pale, coarse, sandy material called \emph{saprolite}. Over volcanic rocks (basalt, trachyte) of the Cameroon Volcanic Line it may be a reddish, clay-rich saprolite.

\textbf{The R horizon --- the bedrock}\par
At the base lies the \cle{R horizon} (R for \emph{rock}): the unweathered \emph{bedrock}, hard and continuous. It is not strictly soil, but it is the parent material from which the whole profile was derived.

\begin{propriete}[Eluviation and illuviation]
In a mature, freely drained soil, \cle{eluviation} (the washing \emph{out} of clay, soluble minerals, and sesquioxides by percolating water) dominates in the upper mineral horizons (A and E), while \cle{illuviation} (the washing \emph{in} and accumulation of those same materials) dominates in the B horizon. The B horizon is therefore enriched in clay, iron and aluminium oxides at the expense of the horizons above it. (This property must be stated precisely in the examination; no proof is required, but you must be able to link each process to the correct horizon.)
\end{propriete}

\begin{popfigure}
\begin{tikzpicture}[x=1cm,y=1cm]
% O horizon
\fill[popdark!85!black] (0,8.4) rectangle (6,9.2);
\draw[popgreen,thick] (0.4,9.2) -- (0.4,9.5) (1.2,9.2) -- (1.2,9.55) (2.1,9.2) -- (2.1,9.5) (3.0,9.2) -- (3.0,9.6) (4.0,9.2) -- (4.0,9.5) (5.0,9.2) -- (5.0,9.55) (5.7,9.2) -- (5.7,9.5);
% A horizon
\fill[popdark!55!popgold] (0,6.8) rectangle (6,8.4);
% E horizon
\fill[gray!25] (0,5.6) rectangle (6,6.8);
% B horizon
\fill[poporange!70] (0,3.2) rectangle (6,5.6);
% C horizon
\fill[popgold!35] (0,1.2) rectangle (6,3.2);
\foreach \x in {0.4,1.3,2.4,3.2,4.1,5.2} {\draw[gray] (\x,2.2) circle (0.09); \draw[gray] (\x+0.5,1.6) circle (0.07);}
% R horizon
\fill[gray!60] (0,0) rectangle (6,1.2);
\draw[black!70] (0,0.5) -- (6,0.35); \draw[black!70] (0,0.9) -- (6,0.75);
% frame
\draw[thick] (0,0) rectangle (6,9.2);
% labels
\node[white] at (0.55,8.8) {\textbf{O}};
\node at (0.55,7.6) {\textbf{A}};
\node at (0.55,6.2) {\textbf{E}};
\node at (0.55,4.4) {\textbf{B}};
\node at (0.55,2.2) {\textbf{C}};
\node[white] at (0.55,0.6) {\textbf{R}};
% descriptions
\node[anchor=west,text width=6.4cm,align=left] at (6.25,8.8) {\small Litter and decomposing humus (organic matter)};
\node[anchor=west,text width=6.4cm,align=left] at (6.25,7.6) {\small Topsoil: dark, humus + minerals; \textbf{eluviation} (leaching); roots, organisms};
\node[anchor=west,text width=6.4cm,align=left] at (6.25,6.2) {\small Strongly leached, pale; quartz remains (forest/podzolic soils)};
\node[anchor=west,text width=6.4cm,align=left] at (6.25,4.4) {\small Subsoil: \textbf{illuviation}; clay, Fe and Al oxides accumulate; reddish};
\node[anchor=west,text width=6.4cm,align=left] at (6.25,2.2) {\small Partially weathered parent rock (regolith, saprolite)};
\node[anchor=west,text width=6.4cm,align=left] at (6.25,0.6) {\small Unweathered bedrock};
% leaching arrow
\draw[-latex,very thick,popblue] (5.5,8.2) -- (5.5,3.6) node[midway,right,popblue,text width=1.6cm,align=center] {\scriptsize percolating water};
\end{tikzpicture}
\legende{Idealised soil profile showing the master horizons O, A, E, B, C and R, with their typical colours, compositions and dominant processes.}
\end{popfigure}

\begin{center}
\begin{tabularx}{\linewidth}{|l|X|X|X|X|}\hline
\rowcolor{popblueL}
\textbf{Horizon} & \textbf{Common name} & \textbf{Key features} & \textbf{Dominant process} & \textbf{Typical in Cameroon} \\ \hline
O & Organic layer & Litter, humus; dark brown/black & Organic accumulation & Thin in humid forest (rapid decay) \\ \hline
A & Topsoil & Mineral + humus; dark; crumbly; roots & Eluviation (leaching) & Fertile layer, 10--30 cm \\ \hline
E & Leached layer & Pale, sandy/silty; quartz-rich; low Fe/Al & Intense eluviation & Rare/absent in tropics; podzols \\ \hline
B & Subsoil & Clay, Fe/Al oxides; red/yellow; compact & Illuviation (accumulation) & Thick lateritic B (metres) in South \\ \hline
C & Regolith / Saprolite & Weathered rock fragments; rock structure visible & Weathering & Deep saprolite on basement/volcanics \\ \hline
R & Bedrock & Fresh, unweathered rock & --- & Granite/gneiss (Congo craton), basalt (CVL) \\ \hline
\end{tabularx}
\end{center}

\begin{aretenir}[Key points on horizons]
\begin{itemize}
\item \textbf{O}: organic litter and humus at the surface; subdivisions Oi, Oe, Oa.
\item \textbf{A}: dark topsoil, humus mixed with minerals; zone of eluviation; the fertile layer; Ap if cultivated.
\item \textbf{E}: pale, strongly leached layer (not always present); quartz residue; hallmark of podzols.
\item \textbf{B}: subsoil, zone of illuviation; accumulates clay and Fe/Al oxides; reddish in tropics; Bt, Bw, Bs, Bk.
\item \textbf{C}: partially weathered parent material (regolith/saprolite); retains rock structure.
\item \textbf{R}: unweathered bedrock; parent material.
\end{itemize}
\end{aretenir}

\courssub{Comparing Soil Profiles: Humid Forest Zone versus Sahelian North}

The depth and development of a soil profile depend directly on the factors of soil formation studied earlier --- especially \emph{climate} and \emph{time}. Cameroon offers a spectacular natural laboratory for this comparison.

In the \textbf{humid forest zone} of the South (around Yaoundé, Kribi or Ebolowa), high temperatures (\SIrange{22}{26}{\ensuremath{{}^{\circ}\mathrm{C}}}) and abundant rainfall (often above \SI{1500}{mm} per year, up to \SI{4000}{mm} at Debundscha) drive intense chemical weathering and leaching over very long periods. The result is a \emph{deep lateritic profile}, commonly several metres to tens of metres thick: a thin O and A horizon, a thick red or yellow B horizon rich in iron and aluminium oxides (ferralitic/Bs horizon), and a deep saprolitic C horizon over the crystalline basement. These are \emph{Ferralsols} (FAO) or \emph{Oxisols} (USDA).

In the \textbf{Sahelian north} (around Maroua and the Far North), rainfall is low (\SIrange{500}{1000}{mm} per year) and seasonal, chemical weathering is weak, and wind and sheet erosion remove fine material. Soils are \emph{thin and immature}: a shallow, pale A horizon poor in humus rests almost directly on a stony C horizon and bedrock. The B horizon is weakly developed (Bw or absent), and organic matter is scarce because vegetation cover is sparse. These are often \emph{Leptosols}, \emph{Regosols} or \emph{Arenosols}.

\begin{popfigure}
\begin{tikzpicture}[x=1cm,y=1cm]
% ---- Left: lateritic profile ----
\begin{scope}
\fill[popdark!85!black] (0,7.6) rectangle (4,8.0);
\fill[popdark!55!popgold] (0,6.6) rectangle (4,7.6);
\fill[poporange!70] (0,3.0) rectangle (4,6.6);
\fill[popgold!35] (0,1.0) rectangle (4,3.0);
\fill[gray!60] (0,0) rectangle (4,1.0);
\draw[thick] (0,0) rectangle (4,8.0);
\draw[popgreen,thick] (0.6,8.0)--(0.6,8.3) (1.6,8.0)--(1.6,8.35) (2.6,8.0)--(2.6,8.3) (3.4,8.0)--(3.4,8.35);
\node at (0.45,7.8) {\scriptsize\textbf{O}};
\node at (0.45,7.1) {\scriptsize\textbf{A}};
\node at (0.45,4.8) {\scriptsize\textbf{B}};
\node at (0.45,2.0) {\scriptsize\textbf{C}};
\node[white] at (0.45,0.5) {\scriptsize\textbf{R}};
\node[align=center] at (2,-0.7) {\small\textbf{Humid forest zone}\\ \small (deep lateritic profile, South Cameroon)};
\node[anchor=west,text width=2.6cm,align=left] at (4.15,4.8) {\scriptsize thick red B horizon (Fe/Al oxides)};
\node[anchor=west,text width=2.6cm,align=left] at (4.15,2.0) {\scriptsize deep saprolite};
\end{scope}
% ---- Right: thin arid profile ----
\begin{scope}[xshift=8.2cm]
\fill[popgold!50] (0,7.2) rectangle (4,8.0);
\fill[popgold!30] (0,6.2) rectangle (4,7.2);
\foreach \x in {0.5,1.4,2.2,3.1,3.7} {\draw[gray] (\x,6.7) circle (0.08);}
\fill[gray!60] (0,5.2) rectangle (4,6.2);
\draw[thick] (0,5.2) rectangle (4,8.0);
\draw[popgreen!70!black,thick] (1.0,8.0)--(1.0,8.25) (3.0,8.0)--(3.0,8.2);
\node at (0.45,7.6) {\scriptsize\textbf{A}};
\node at (0.45,6.7) {\scriptsize\textbf{C}};
\node[white] at (0.45,5.7) {\scriptsize\textbf{R}};
\node[align=center] at (2,4.5) {\small\textbf{Sahelian zone}\\ \small (thin immature soil, Far North)};
\node[anchor=west,text width=2.6cm,align=left] at (4.15,7.6) {\scriptsize thin pale A, little humus};
\node[anchor=west,text width=2.6cm,align=left] at (4.15,6.7) {\scriptsize stony regolith};
\end{scope}
\end{tikzpicture}
\legende{Comparison of soil profiles in Cameroon: deep lateritic profile of the humid south (left) and thin, immature soil of the arid north (right), drawn to the same horizontal scale.}
\end{popfigure}

\begin{savaistu}[Laterite and road building in Cameroon]
The red lateritic material of the B horizon has been used for generations in Cameroon to surface rural roads and to make sun-dried or compressed earth blocks for houses. Laterite hardens on exposure to air because its iron oxides cement the particles together --- a property known since antiquity (the word \emph{laterite} comes from the Latin \emph{later}, ``brick''). But when the protective A horizon is stripped off, the exposed subsoil erodes quickly, which is one reason why road cuttings degrade so fast in the rainy season. The term \emph{laterite} was coined by Francis Buchanan in 1807 in India; in Cameroon it is a major construction material.
\end{savaistu}

\courssub{Other Profile Types in Cameroon}

\begin{itemize}
\item \textbf{Volcanic soils (Andosols)} on recent ashes and lava of Mount Cameroon and the Cameroon Volcanic Line: often show an \textbf{A--C--R} or \textbf{A--Bw--C--R} sequence with a thick, dark, humus-rich A horizon (high organic matter, low bulk density, high water retention) over weakly weathered volcanic parent material. Little leaching due to youth.
\item \textbf{Alluvial soils (Fluvisols)} in the Wouri estuary, Logone and Benue floodplains: stratified layers (C1, C2, ...) of sediment with buried A horizons; no well-developed B horizon due to frequent deposition.
\item \textbf{Hydromorphic soils (Gleysols)} in valley bottoms (bas-fonds) and around Lake Chad: grey/blue-green colours (gleying) due to reduction of iron in waterlogged conditions; mottled Bg horizons.
\item \textbf{Vertisols} in the Logone plain and parts of the Benue trough: high smectite clay content causes deep cracking when dry and slickensides in B horizon; profile shows A--Bw--C with wedge-shaped structure.
\end{itemize}

\courssub{Reading and Labelling a Soil Profile: a Practical Skill}

In GCE Paper 2 you may be given a photograph or diagram of a soil profile and asked to label the horizons and justify your choices. The following method guarantees full marks.

\begin{methode}[How to describe a soil profile in an exam]
\begin{enumerate}
\item \textbf{Work from top to bottom}, horizon by horizon, never jumping around.
\item For each horizon, \textbf{state four things}: its \emph{depth} (or relative thickness), its \emph{colour} (use Munsell notation if known, e.g. 5YR 4/6), its \emph{texture} (sandy, clayey, stony, crumbly), and its \emph{composition} (humus, roots, oxides, rock fragments, concretions).
\item \textbf{Name the horizon} with its letter and its common name (e.g.\ ``B horizon --- subsoil'') and any subordinate distinction (e.g. Bt, Bs).
\item \textbf{Name the dominant process} where relevant: eluviation/leaching in A and E, illuviation/accumulation in B, weathering in C, organic accumulation in O.
\item Describe \textbf{horizon boundaries} (abrupt, clear, gradual, diffuse; smooth, wavy, irregular, broken).
\item Conclude with \textbf{one sentence on the soil type and environment} (e.g.\ deep lateritic soil/Ferralsol of a humid tropical climate on basement complex).
\end{enumerate}
\end{methode}

\begin{exemplebox}[Identifying a horizon from a description]
\emph{``A pale grey, sandy layer, almost free of iron compounds, lying between a dark humus-rich layer above and a reddish clay-rich layer below.''} This is the \textbf{E horizon}: its light colour shows that iron and clay have been washed out (eluviation), and its position between A and B confirms it.
\end{exemplebox}

\begin{exemplebox}[Identifying a subordinate horizon]
\emph{``A subsoil horizon with distinct clay skins (argillans) coating ped faces and lining pores, overlying a saprolite.''} This is a \textbf{Bt horizon} (B with clay illuviation), typical of an \emph{Alfisol} or \emph{Luvisol}, less common in the deeply leached South but found in transitional zones.
\end{exemplebox}

\begin{exoresolu}[Worked exercise: labelling a profile]
A soil pit dug near Bafoussam reveals, from the surface downwards: (1) a 3 cm layer of dead leaves and dark decomposing matter; (2) a 25 cm dark brown crumbly layer full of roots; (3) a 90 cm red, compact, clay-rich layer with no roots; (4) a pale yellow layer of soft, crumbly rock fragments; (5) hard fresh volcanic rock. Identify each layer and name the dominant process in layers (2) and (3).
\tcblower
\textbf{Solution.}
\begin{itemize}
\item (1) \textbf{O horizon} --- organic litter and humus (Oe/Oa).
\item (2) \textbf{A horizon} (topsoil): dark because humus is mixed with minerals; roots confirm biological activity. Dominant process: \textbf{eluviation} (leaching by percolating rainwater).
\item (3) \textbf{B horizon} (subsoil): red colour and clay enrichment show \textbf{illuviation} --- accumulation of iron oxides and clay washed down from above. Absence of roots and compactness are typical. Subordinate: \textbf{Bs} (sesquioxide accumulation) or \textbf{Bt} if clay skins visible.
\item (4) \textbf{C horizon}: partially weathered parent material (regolith/saprolite) derived from basalt/trachyte.
\item (5) \textbf{R horizon}: unweathered bedrock (volcanic rock of the Cameroon Volcanic Line).
\end{itemize}
Conclusion: this is a \emph{deep, mature profile of a humid tropical environment (Ferralsol/Andosol transition)}, developed on volcanic parent rock.
\end{exoresolu}

\begin{exercice}[Test yourself]
\begin{enumerate}
\item Define \emph{soil profile} and \emph{soil horizon}.
\item Draw and label a complete soil profile, indicating the dominant process in each mineral horizon.
\item Explain why the B horizon of soils around Yaoundé is red while the E horizon of a podzol is pale grey.
\item A pupil claims that humus accumulates at the base of the profile ``because it is heavy''. Correct this statement.
\item Distinguish between a \textbf{Bt} horizon and a \textbf{Bs} horizon. Give an example of where each might be found in Cameroon.
\item Describe the horizon sequence you would expect in a soil developed on recent volcanic ash on the slopes of Mount Cameroon.
\end{enumerate}
\end{exercice}

\begin{corrige}[Exercise 1]
\begin{enumerate}
\item A soil profile is a vertical section through the soil from the surface to the bedrock; a soil horizon is one of the horizontal layers distinguished within it by colour, texture, structure, consistence and composition.
\item The diagram should show, from top to bottom: O (litter/humus, organic accumulation), A (topsoil, eluviation), E (leached, if present, intense eluviation), B (subsoil, illuviation of clay and Fe/Al oxides), C (weathered parent material, weathering), R (bedrock).
\item Around Yaoundé, intense leaching under a humid climate removes soluble minerals but leaves iron oxides, which accumulate in the B horizon and give it a red colour (Bs horizon). In a podzol, even the iron is washed out of the E horizon under very acidic conditions, leaving pale quartz sand.
\item Humus forms from dead plants and animals at the surface and is mixed into the topsoil by organisms; it is therefore found in the O and A horizons at the \emph{top} of the profile, not at the bottom. The base of the profile is parent material and bedrock, which contain no humus.
\item \textbf{Bt}: horizon of clay illuviation (argillans/clay skins visible); typical of Luvisols/Alfisols, found in less intensely leached zones (e.g., transitional zones between forest and savanna, Adamawa plateau). \textbf{Bs}: horizon of sesquioxide (Fe/Al oxide) accumulation; typical of Ferralsols/Oxisols, dominant in the humid forest zone of South Cameroon.
\item On recent volcanic ash (Mount Cameroon): \textbf{O} (thin litter) -- \textbf{A} (thick, dark, humus-rich, low bulk density, andic properties) -- \textbf{Bw} (weakly developed, structure formation) or directly \textbf{C} (fresh ash/lapilli) -- \textbf{R} (basalt). Little leaching due to youth; high organic matter, high water retention.
\end{enumerate}
\end{corrige}

\begin{attention}[Exam tip]
Learn \textbf{one} labelled profile diagram perfectly --- the six-horizon profile with colours and processes shown in this section. Profile-labelling and profile-description questions appear very frequently in GCE Paper 2, and a clear, correctly labelled diagram with brief annotations often earns most of the marks on its own. Practise drawing it from memory in under three minutes. Remember to include subordinate distinctions (Ap, Bt, Bs, Bw) if the description warrants it.
\end{attention}

\begin{experience}[Field / Practical Activity: Soil Profile Description]
\textbf{Objective:} Describe a soil profile in the field or from a photograph/prepared pit.
\textbf{Materials:} Spade/auger, Munsell colour chart, tape measure, knife, hand lens, water bottle, HCl (10\%) for carbonate test, notebook.
\textbf{Procedure:}
\begin{enumerate}
\item Dig or locate a fresh soil pit at least 1 m deep (or to bedrock).
\item Clean one face vertically with a knife to reveal structure and colours.
\item Measure and record the depth of each horizon boundary.
\item For each horizon, record: (a) Master horizon symbol + subordinate (e.g., Ap, Bt1, Bt2, BC), (b) Depth range (cm), (c) Moist colour (Munsell), (d) Texture (field estimate: sand/silt/clay), (e) Structure (grade, size, type), (f) Consistence (wet, moist, dry), (g) Roots (abundance, size), (h) Pores, (i) Concretions/nodules, (j) Reaction to HCl, (k) Boundary distinctness and topography.
\item Sketch the profile to scale, label horizons, add a scale bar and north arrow.
\item Classify the soil type (e.g., Ferralsol, Andosol, Fluvisol) and relate to landscape position and parent material.
\end{enumerate}
\textbf{Safety:} Do not enter pits deeper than 1 m without shoring; wear gloves; wash hands after handling soil.
\end{experience}

Now that we can recognise and describe the layers of a soil, the next step is to examine what gives each soil its particular behaviour: its \emph{physical characteristics} --- texture, structure, porosity and water retention --- which we study in the following section.

\coursec{Physical Characteristics of Soil}

Having examined the vertical arrangement of horizons in a soil profile, we now turn to the intrinsic physical properties that determine how a soil behaves — how it holds water, how easily roots penetrate it, how it is worked by farmers, and ultimately how productive it can be. These properties are the \emph{physical characteristics of soil}: texture, structure, colour, porosity, aeration, water-holding capacity and depth. In Cameroon, from the volcanic slopes of Mount Cameroon to the ferralitic plateaus of the Adamawa, these characteristics explain why certain soils support thriving cocoa or coffee plantations while others require careful management for maize or rice.

\courssub{Soil Texture: The Relative Proportions of Sand, Silt and Clay}

\begin{definition}[Soil Texture]
Soil texture is the \textbf{relative proportion by weight} of the three mineral particle-size fractions in a soil: \cle{sand} ($2.0$–$0.05\ \mathrm{mm}$), \cle{silt} ($0.05$–$0.002\ \mathrm{mm}$) and \cle{clay} ($<0.002\ \mathrm{mm}$). Texture is a \emph{permanent} property — it does not change with cultivation or weathering on human timescales.
\end{definition}

The size limits above follow the international (ISSS/USDA) classification used in GCE examinations. Note that \emph{gravel} ($>2\ \mathrm{mm}$) and \emph{stones} are not part of the fine earth fraction and are removed before texture analysis.

\begin{propriete}[The Textural Triangle]
Any soil can be placed on the \textbf{textural triangle} (also called the texture diagram) by reading its percentages of sand, silt and clay. The three axes represent the three fractions; the sum is always $100\%$. The triangle is divided into \textbf{12 textural classes} (sand, loamy sand, sandy loam, loam, silt loam, silt, sandy clay loam, clay loam, silty clay loam, sandy clay, silty clay, clay). For agricultural purposes, we often group them into four broad categories:
\begin{itemize}
\item \textbf{Sandy soils:} $>70\%$ sand, $<15\%$ clay.
\item \textbf{Silty soils:} $>80\%$ silt, $<12\%$ clay.
\item \textbf{Clayey soils:} $>35\%$ clay (often $>40\%$).
\item \textbf{Loamy soils:} roughly balanced mixtures — the most favourable for agriculture.
\end{itemize}
\end{propriete}

\begin{popfigure}
\begin{tikzpicture}[scale=0.7]
% Triangle vertices
\coordinate (S) at (0,0);     % Sand (100% sand, 0% silt, 0% clay)
\coordinate (Si) at (10,0);   % Silt (0% sand, 100% silt, 0% clay)
\coordinate (C) at (5,8.66);  % Clay (0% sand, 0% silt, 100% clay)
% Draw triangle outline
\draw[popdark, thick] (S) -- (Si) -- (C) -- cycle;
% Grid lines for constant sand % (parallel to Si-C side)
\foreach \p in {10,20,...,90} {
    \pgfmathsetmacro{\f}{\p/100}
    \coordinate (A) at ($(S)!\f!(Si)$); % point on base S-Si
    \coordinate (B) at ($(S)!\f!(C)$);  % point on left side S-C
    \draw[popline, very thin] (A) -- (B);
    % Label on base (sand axis)
    \node[popmuted, font=\tiny, anchor=north] at ($(S)!0.5!(A)$) {\p\%};
}
% Grid lines for constant silt % (parallel to S-C side)
\foreach \p in {10,20,...,90} {
    \pgfmathsetmacro{\f}{\p/100}
    \coordinate (A) at ($(Si)!\f!(S)$); % point on base Si-S
    \coordinate (B) at ($(Si)!\f!(C)$);  % point on right side Si-C
    \draw[popline, very thin] (A) -- (B);
    % Label on right side (silt axis)
    \node[popmuted, font=\tiny, anchor=west] at ($(Si)!0.5!(B)$) {\p\%};
}
% Grid lines for constant clay % (parallel to S-Si base)
\foreach \p in {10,20,...,90} {
    \pgfmathsetmacro{\f}{\p/100}
    \coordinate (A) at ($(C)!\f!(S)$);  % point on left side C-S
    \coordinate (B) at ($(C)!\f!(Si)$); % point on right side C-Si
    \draw[popline, very thin] (A) -- (B);
    % Label on left side (clay axis)
    \node[popmuted, font=\tiny, anchor=east] at ($(C)!0.5!(A)$) {\p\%};
}
% Vertex labels
\node[popdark, font=\bfseries] at (S) [below left] {Sand};
\node[popdark, font=\bfseries] at (Si) [below right] {Silt};
\node[popdark, font=\bfseries] at (C) [above] {Clay};
% Axis labels (corrected)
\node[popdark, rotate=60, font=\small] at (2.5,4.33) {Clay \%};   % left side
\node[popdark, rotate=-60, font=\small] at (7.5,4.33) {Silt \%};  % right side
\node[popdark, font=\small] at (5,-0.5) {Sand \%};                 % base
% Example point: 40% sand, 40% silt, 20% clay -> barycentric (0.4,0.4,0.2)
\coordinate (P) at (5, 1.732);
\fill[poporange] (P) circle (3pt);
\node[poporange, font=\small\bfseries, anchor=south] at (P) {Loam (40/40/20)};
\end{tikzpicture}
\legende{The soil textural triangle. Follow the grid lines parallel to each side to read the percentage of each fraction. The orange dot marks the ideal loam (40\% sand, 40\% silt, 20\% clay).}
\end{popfigure}

\begin{aretenir}[Reading the Textural Triangle — Exam Tip]
To locate a soil on the triangle:
\begin{enumerate}
\item Find the \% sand on the \textbf{base} (0–100\% from left to right). Draw a line \textbf{parallel to the right side} (silt–clay side).
\item Find the \% silt on the \textbf{right side} (0–100\% from bottom to top). Draw a line \textbf{parallel to the left side} (sand–clay side).
\item Find the \% clay on the \textbf{left side} (0–100\% from bottom to top). Draw a line \textbf{parallel to the base} (sand–silt side).
\item The three lines intersect at a single point — the soil's texture class.
\end{enumerate}
Practise with several combinations until you can do it quickly.
\end{aretenir}

\begin{propriete}[Loam — The Ideal Agricultural Texture]
A \textbf{loam} contains approximately \textbf{40\% sand, 40\% silt and 20\% clay}. This balance gives:
\begin{itemize}
\item Good \textbf{drainage} and \textbf{aeration} (from sand).
\item Good \textbf{water-holding capacity} and \textbf{nutrient retention} (from clay).
\item Good \textbf{workability} and \textbf{root penetration} (from silt).
\end{itemize}
Most high-yielding crops (maize, cocoa, coffee, oil palm) perform best on loamy soils. In Cameroon, the volcanic ash-derived soils on the flanks of Mount Cameroon and the recent alluvial soils of the Wouri and Sanaga valleys often have a loamy texture.
\end{propriete}

\begin{methode}[The Jar Settling Test — Estimating Texture in the Field]
This simple test gives a visual estimate of the sand, silt and clay percentages.
\begin{enumerate}
\item Take a representative sample of fine earth (remove gravel, roots, organic debris).
\item Fill a straight-sided glass jar about 1/3 full with the soil.
\item Add water to about 3/4 full, add a teaspoon of detergent (dispersing agent) and shake vigorously for 5–10 minutes to break all aggregates.
\item Set the jar on a level surface and time the settling:
\begin{itemize}
\item \textbf{Sand} settles in \textbf{1–2 minutes}. Mark the top of the sand layer.
\item \textbf{Silt} settles in \textbf{1–2 hours}. Mark the top of the silt layer.
\item \textbf{Clay} may take \textbf{several days} to settle completely. The water clears when clay has settled.
\end{itemize}
\item Measure the thickness of each layer (and the total thickness). Compute percentages:
\[
\% \text{sand} = \frac{\text{sand thickness}}{\text{total thickness}} \times 100,\quad
\% \text{silt} = \frac{\text{silt thickness}}{\text{total thickness}} \times 100,\quad
\% \text{clay} = 100 - (\% \text{sand} + \% \text{silt}).
\]
\item Plot the percentages on the textural triangle to name the texture class.
\end{enumerate}
\end{methode}

\begin{exoresolu}[Determining Texture from a Jar Test]
A student performs the jar settling test on a soil sample from a farm near Buea. After settling, the layers measure: sand = 2.8 cm, silt = 1.6 cm, clay = 0.6 cm. Total thickness = 5.0 cm.
\tcblower
\textbf{Solution:}
\[
\% \text{sand} = \frac{2.8}{5.0} \times 100 = 56\%
\]
\[
\% \text{silt} = \frac{1.6}{5.0} \times 100 = 32\%
\]
\[
\% \text{clay} = 100 - (56 + 32) = 12\%
\]
Plot (56\% sand, 32\% silt, 12\% clay) on the textural triangle. The point falls in the \textbf{sandy loam} class (borderline loam). This soil will drain well but may need organic matter to improve water retention for dry-season crops.

\end{exoresolu}

\courssub{Soil Structure: The Arrangement of Particles into Aggregates}

\begin{definition}[Soil Structure]
Soil structure is the \textbf{arrangement of primary soil particles (sand, silt, clay) into secondary units called \cle{peds} or \cle{aggregates}}. Unlike texture, structure is \emph{dynamic} — it can be improved by organic matter, root action, earthworms, and careful tillage, or destroyed by compaction, excessive cultivation, or sodium dispersion.
\end{definition}

\begin{attention}[Texture vs. Structure — Do Not Confuse Them!]
\begin{itemize}
\item \textbf{Texture} = \emph{what the soil is made of} (proportions of sand, silt, clay). It is \textbf{inherited} from the parent material and changes only over geological time.
\item \textbf{Structure} = \emph{how the particles are organised} (aggregates, pore spaces). It is \textbf{manageable} — farmers can improve or degrade it within a single season.
\end{itemize}
\textbf{Exam trap:} A question may ask "How does texture affect drainage?" (answer: clay holds water, sand drains fast) versus "How does structure affect drainage?" (answer: crumb structure creates large pores for rapid drainage; platy structure impedes drainage). Always identify which property is being discussed.
\end{attention}

The main structural types (shape of peds) are:

\begin{center}
\begin{tabularx}{\linewidth}{|l|X|X|}\hline
\textbf{Type} & \textbf{Description} & \textbf{Typical Horizon / Significance} \\ \hline
\rowcolor{popblueL}
\textbf{Crumb / Granular} & Small, rounded, porous aggregates (1–10 mm). & \textbf{Topsoil (A horizon)} rich in organic matter. Excellent aeration, drainage, root growth, easy tillage. \\ \hline
\textbf{Blocky} (angular or subangular) & Roughly cube-shaped, 10–50 mm. & \textbf{Subsoil (B horizon)}. Moderate drainage; can be improved by organic matter. \\ \hline
\textbf{Platy} & Flat, plate-like, horizontal. & Often in \textbf{compacted layers (E or Bt)} or surface crusts. \textbf{Impedes} water movement and root penetration. \\ \hline
\textbf{Prismatic / Columnar} & Vertical columns, 50–100+ mm, often with flat or rounded tops. & \textbf{Clay-rich subsoils (Bt)} in humid tropics; columnar tops may have salt caps in arid zones. Can restrict roots if dense. \\ \hline
\textbf{Single grain} & No aggregation — loose sand. & \textbf{Sandy soils (A or C)}. Very rapid drainage, low water holding. \\ \hline
\textbf{Massive} & No visible structure — coherent solid mass. & \textbf{Compacted or clayey subsoils}. Very poor aeration and drainage. \\ \hline
\end{tabularx}
\end{center}

\begin{savaistu}[Cameroon Context: Structure in Volcanic and Ferralitic Soils]
On the slopes of Mount Cameroon, young volcanic ash soils often have a \textbf{crumb structure} due to high organic matter and allophane minerals — they are friable and highly productive. In contrast, the deep ferralitic soils (Oxisols) of the southern plateau have a \textbf{weak subangular blocky to granular structure} stabilised by iron oxides; they are permeable but low in nutrient reserves. In the irrigated rice schemes of the Logone floodplain, \textbf{platy structure} can develop in the puddled layer, which is actually desirable to reduce percolation losses.
\end{savaistu}

\courssub{Soil Colour: A Quick Diagnostic Indicator}

Soil colour is primarily due to \textbf{organic matter} (dark) and \textbf{iron oxides} (red, yellow, brown). It is described using the \textbf{Munsell colour notation} (e.g., 10YR 3/2 = very dark greyish brown), but in the field we use simple terms:

\begin{propriete}[Colour as an Indicator of Soil Conditions]
\begin{itemize}
\item \textbf{Dark brown / black:} High \textbf{humus} content (organic matter). Typical of well-drained topsoils (A horizons) under forest or grassland. \emph{Example:} Forest soils of the Congo Basin margin.
\item \textbf{Red / reddish brown:} \textbf{Hematite} ($\mathrm{Fe_2O_3}$) — well-drained, oxidising conditions, often old, highly weathered soils (laterites). \emph{Example:} Ferralitic soils of the Adamawa and South Cameroon plateaus.
\item \textbf{Yellow / yellowish brown:} \textbf{Goethite} ($\mathrm{FeO\cdot OH}$) — slightly less drained than red soils, still oxidising.
\item \textbf{Grey / blue-grey / greenish:} \textbf{Gley colours} — reduced iron ($\mathrm{Fe^{2+}}$) due to \textbf{waterlogging} (anaerobic conditions). \emph{Example:} Hydromorphic soils in the Wouri estuary, Nyong valley, or inland valleys (bas-fonds) used for rice.
\item \textbf{Mottles (spots of contrasting colour):} Indicate \textbf{fluctuating water table} — alternating oxidation (red/yellow) and reduction (grey).
\end{itemize}
\end{propriete}

\begin{experience}[Observing Soil Colour in a Profile]
Dig a pit to 1 m depth in a typical Cameroonian landscape (e.g., a cocoa farm on ferralitic soil). Clean a vertical face and observe:
\begin{enumerate}
\item Top 20 cm: dark brown (10YR 3/3) — humus-rich A horizon.
\item 20–60 cm: strong brown (7.5YR 5/8) — B horizon with iron oxides.
\item 60–100 cm: red (2.5YR 4/8) — deeper B horizon, more hematite.
\item If you dig in a valley bottom: grey matrix (5Y 5/1) with yellow mottles (10YR 6/8) — gleyed horizon.
\end{enumerate}
Record the Munsell notation if possible; otherwise use the descriptive terms above.
\end{experience}

\courssub{Soil Porosity, Aeration and Water-Holding Capacity}

\begin{definition}[Porosity]
\cle{Porosity} is the \textbf{percentage of soil volume occupied by pore spaces} (voids between particles and aggregates). Total porosity = (1 – bulk density / particle density) × 100\%. Pores are classified by size:
\begin{itemize}
\item \textbf{Macropores} ($>0.06\ \mathrm{mm}$): drain by gravity, conduct air and water rapidly.
\item \textbf{Micropores} ($<0.06\ \mathrm{mm}$): hold water against gravity (capillary water).
\end{itemize}
\end{definition}

Texture and structure together control porosity:
\begin{itemize}
\item \textbf{Sandy soils:} Large particles → large pores (macropores) → \textbf{high total porosity} (40–50\%) but \textbf{low water-holding capacity} because water drains quickly through macropores.
\item \textbf{Clayey soils:} Small particles → many micropores → \textbf{high total porosity} (50–60\%) and \textbf{high water-holding capacity}, but \textbf{poor aeration} and \textbf{slow drainage} because micropores retain water and restrict air exchange.
\item \textbf{Loamy soils:} Balanced pore-size distribution → \textbf{good aeration} (enough macropores) and \textbf{good available water} (enough micropores).
\end{itemize}

\begin{popfigure}
\begin{tikzpicture}
\begin{axis}[
    ybar,
    bar width=12pt,
    width=0.9\linewidth,
    height=8cm,
    ylabel={Relative rating (arbitrary units)},
    xtick={1,2,3},
    xticklabels={Sandy, Loamy, Clayey},
    ymin=0, ymax=10,
    legend style={at={(0.5,-0.15)}, anchor=north, legend columns=2},
    ymajorgrids=true,
    grid style=dashed,
]
\addplot[fill=poporange] coordinates {(1,8) (2,5) (3,2)};
\addplot[fill=popblue] coordinates {(1,2) (2,7) (3,9)};
\addplot[fill=popgreen] coordinates {(1,3) (2,8) (3,4)};
\legend{Drainage rate, Water-holding capacity, Aeration (macroporosity)}
\end{axis}
\end{tikzpicture}
\legende{Comparison of drainage rate, water-holding capacity and aeration (macroporosity) for sandy, loamy and clayey soils. Loam offers the best compromise.}
\end{popfigure}

\begin{aretenir}[Key Relationships]
\begin{itemize}
\item \textbf{Field capacity (FC):} Water held after free drainage ceases (≈ -33 kPa). Clay > Loam > Sand.
\item \textbf{Permanent wilting point (PWP):} Water held too tightly for roots to extract (≈ -1500 kPa). Clay > Loam > Sand.
\item \textbf{Available water capacity (AWC) = FC – PWP.} Loam often has the \textbf{highest AWC} because clay holds much water but at tensions too high for plants (small pores), while sand holds little water at any tension.
\item \textbf{Aeration:} Roots need $\mathrm{O_2}$; $\mathrm{CO_2}$ must diffuse out. Macropores ($>0.06\ \mathrm{mm}$) are essential. Clayey soils easily become anaerobic when wet.
\end{itemize}
\end{aretenir}

\courssub{Soil Depth and Its Agricultural Significance}

\cle{Soil depth} (or \cle{solum depth}) is the thickness of the weathered material (A + B horizons) above the unweathered parent material (C horizon) or bedrock. It is classified as:
\begin{itemize}
\item \textbf{Very shallow} $< 25\ \mathrm{cm}$
\item \textbf{Shallow} $25$–$50\ \mathrm{cm}$
\item \textbf{Moderately deep} $50$–$100\ \mathrm{cm}$
\item \textbf{Deep} $100$–$150\ \mathrm{cm}$
\item \textbf{Very deep} $> 150\ \mathrm{cm}$
\end{itemize}

\begin{propriete}[Importance of Soil Depth]
\begin{itemize}
\item \textbf{Root development:} Deep soils allow deep rooting → access to more water and nutrients → drought resilience. Shallow soils restrict roots, causing moisture stress.
\item \textbf{Water storage:} Total available water = AWC × rooting depth. A deep loam can store $200+$ mm of available water; a shallow sand may store $<50$ mm.
\item \textbf{Management:} Deep soils tolerate more intensive cultivation; shallow soils require conservation tillage, agroforestry, or permanent cover.
\end{itemize}
In Cameroon, the deep volcanic soils on Mount Cameroon (often $>2\ \mathrm{m}$) support perennial tree crops (banana, oil palm) with high yields, while the shallow lithosols on steep slopes of the Mandara Mountains are best left under permanent vegetation or used for grazing.
\end{propriete}

\courssub{Simple Field Tests for Physical Characteristics}

Besides the jar settling test (Method box above), every geology student should master these quick field assessments:

\begin{methode}[Feel / Ribbon Test for Texture]
\begin{enumerate}
\item Take a handful of moist soil (not wet, not dry). Knead it to break aggregates.
\item \textbf{Gritty feel} → sand dominant.
\item \textbf{Silky / floury feel} → silt dominant.
\item \textbf{Sticky, plastic, forms long ribbons} → clay dominant.
\item \textbf{Ribbon test:} Squeeze soil between thumb and forefinger to form a ribbon.
\begin{itemize}
\item No ribbon (crumbles) → sand / loamy sand.
\item Short ribbon ($<2.5\ \mathrm{cm}$) → sandy loam / silt loam.
\item Medium ribbon ($2.5$–$5\ \mathrm{cm}$) → loam / clay loam.
\item Long ribbon ($>5\ \mathrm{cm}$) → clay / sandy clay / silty clay.
\end{itemize}
\end{enumerate}
\end{methode}

\begin{methode}[Visual Assessment of Structure and Colour]
\begin{enumerate}
\item Extract a spade-sized block of soil from the topsoil and subsoil.
\item \textbf{Structure:} Drop the block gently; observe how it breaks. Crumb = good; large hard blocks = poor; platy = compacted.
\item \textbf{Colour:} Note the dominant colour and any mottles. Use a Munsell chart if available, otherwise describe (dark, red, grey, mottled).
\item \textbf{Porosity:} Look for visible pores, root channels, earthworm burrows.
\end{enumerate}
\end{methode}

\begin{exoresolu}[Interpreting Field Observations]
A student examines a soil profile in a maize field near Ngaoundéré. Observations:
\begin{itemize}
\item Topsoil (0–25 cm): dark brown, crumb structure, many roots, friable.
\item Subsoil (25–80 cm): reddish brown, weak subangular blocky, few roots, firm.
\item Below 80 cm: weathered granite (C horizon).
\item Jar test on topsoil: 45\% sand, 35\% silt, 20\% clay.
\end{itemize}
\textbf{Questions:}
\begin{enumerate}
\item Name the texture class of the topsoil.
\item Assess the structure of both horizons for maize root growth.
\item What does the colour change indicate?
\end{enumerate}
\tcblower
\textbf{Solution:}
\begin{enumerate}
\item 45/35/20 plots in the \textbf{loam} class (near the centre of the triangle). Excellent texture for maize.
\item Topsoil: \textbf{crumb structure} — ideal for seedling emergence, aeration, water infiltration. Subsoil: \textbf{weak subangular blocky} — moderate; roots can penetrate but may be limited by firmness. Overall good rooting depth (80 cm).
\item Dark topsoil = organic matter accumulation. Reddish subsoil = iron oxide (hematite) accumulation — well-drained, oxidising conditions typical of ferralitic soils. No mottles → no waterlogging.
\end{enumerate}

\end{exoresolu}

\courssub{Common Mistakes and Exam Tips}

\begin{attention}[Clay Soils Are Not Always "Best" — A Classic Trap]
Many students think "clay holds nutrients and water, so clay soil is the most fertile". \textbf{Wrong.} Clayey soils have high \textbf{potential} fertility (high CEC, high water storage) but often have \textbf{poor actual productivity} because:
\begin{itemize}
\item Poor drainage → waterlogging, denitrification, root rot.
\item Hard when dry, sticky when wet → difficult tillage (narrow workability window).
\item Slow warming in spring (in temperate zones) or after rains.
\item Prone to compaction and crusting.
\end{itemize}
\textbf{Correct statement:} "Loamy soils are generally the most productive because they balance water holding, drainage, aeration and workability." In exams, if asked "Which texture is best for agriculture?", answer \textbf{loam} and explain why.
\end{attention}

\begin{aretenir}[Exam Checklist for Physical Characteristics]
\begin{itemize}
\item Define texture and structure \textbf{distinctly}.
\item Know the size limits: sand $2$–$0.05\ \mathrm{mm}$, silt $0.05$–$0.002\ \mathrm{mm}$, clay $<0.002\ \mathrm{mm}$.
\item Be able to \textbf{plot and read} any point on the textural triangle (follow the three grid directions).
\item Know the \textbf{four broad classes} and the \textbf{ideal loam composition} (40/40/20).
\item Link colour to processes: dark = humus; red/yellow = Fe oxides (drained); grey/blue = gley (waterlogged).
\item Explain how texture and structure control \textbf{porosity, aeration, drainage, water-holding capacity}.
\item Describe the \textbf{jar settling test} and \textbf{feel/ribbon test} step by step.
\item Relate \textbf{soil depth} to rooting volume and drought resilience.
\end{itemize}
\end{aretenir}

\begin{exercice}[Practice Questions]
\begin{enumerate}
\item A soil sample gives the following jar test layers: sand 3.2 cm, silt 1.0 cm, clay 0.8 cm. Determine the texture class.
\item Explain why a soil with 60\% clay may have lower \emph{available} water capacity than a loam with 20\% clay.
\item A farmer in the Logone floodplain observes grey soil with yellow mottles at 40 cm depth. What does this indicate about drainage and suitable crops?
\item Distinguish between soil texture and soil structure. Give one management practice that improves structure but does not change texture.
\item Using the textural triangle, locate the point for 20\% sand, 70\% silt, 10\% clay. Name the texture class.
\end{enumerate}
\end{exercice}

\begin{corrige}[Exercise 1]
Total = 5.0 cm. \% sand = 64\%, \% silt = 20\%, \% clay = 16\%. This plots in the \textbf{sandy loam} class (close to loam boundary).
\end{corrige}

\begin{corrige}[Exercise 2]
Clay holds much water in micropores at tensions $>1500\ \mathrm{kPa}$ (permanent wilting point), so a large fraction is \emph{unavailable} to plants. Loam has a more favourable pore-size distribution, giving a higher \emph{available} water capacity (field capacity – wilting point).
\end{corrige}

\begin{corrige}[Exercise 3]
Grey matrix + yellow mottles = \textbf{gleying} due to a fluctuating water table. The soil is \textbf{poorly drained} (hydromorphic). Suitable crops: \textbf{rice} (tolerates flooding), or dry-season vegetables on raised beds. Maize would suffer from waterlogging.
\end{corrige}

\begin{corrige}[Exercise 4]
Texture = particle size distribution (inherited, permanent). Structure = arrangement of particles into aggregates (dynamic, manageable). Adding organic matter (compost, manure) or growing cover crops improves structure (crumb formation) without changing texture.
\end{corrige}

\begin{corrige}[Exercise 5]
20\% sand, 70\% silt, 10\% clay → \textbf{silt loam} (high silt, low clay).
\end{corrige}

\coursec{Chemical and Biological Characteristics of Soil}

In the previous section, we explored the \emph{physical} architecture of the soil — its texture, structure, porosity, and colour. We saw that a soil's physical properties determine how easily roots can penetrate, how water moves, and how air circulates. However, a soil can have a perfect crumb structure and ideal texture yet remain unproductive if its \emph{chemical} and \emph{biological} engines are not running. This section opens the "chemical laboratory" and the "living factory" beneath our feet. We will see how mineral nutrients are stored, released, and cycled, why the pH of the soil solution acts as a master switch for fertility, and how the invisible army of soil organisms drives the entire system. Understanding these characteristics is the key to interpreting soil analysis reports and making sound management decisions for Cameroonian farms, from the volcanic slopes of Mount Cameroon to the ferrallitic plains of the South.

\courssub{The Composition of an Ideal Soil}

Before diving into chemical reactions, we must visualise the soil as a dynamic four-phase system. An "ideal" agricultural soil — often called a \cle{loam} — is not just solid particles; it is a balanced mixture of solids, liquids, and gases.

\begin{popfigure}
\begin{tikzpicture}
    % Pie chart for ideal soil composition
    \def\anglemin{162} % Minerals 45%
    \def\angleorg{18}  % Organic 5%
    \def\anglewat{90}  % Water 25%
    \def\angleair{90}  % Air 25%
    
    % Minerals (Brown)
    \fill[poporange!70!popdark] (0,0) -- (90:3) arc (90:90-\anglemin:3) -- cycle;
    % Organic Matter (Dark Brown)
    \fill[popdark!80!black] (0,0) -- (90-\anglemin:3) arc (90-\anglemin:90-\anglemin-\angleorg:3) -- cycle;
    % Water (Blue)
    \fill[popblue] (0,0) -- (90-\anglemin-\angleorg:3) arc (90-\anglemin-\angleorg:90-\anglemin-\angleorg-\anglewat:3) -- cycle;
    % Air (Light Grey/White pattern)
    \fill[popmuted!30] (0,0) -- (90-\anglemin-\angleorg-\anglewat:3) arc (90-\anglemin-\angleorg-\anglewat:90-\anglemin-\angleorg-\anglewat-\angleair:3) -- cycle;
    
    % Labels with lines
    % Minerals label
    \node[poporange!70!popdark, font=\bfseries] at (1.5, -1.5) {Mineral Matter $\approx 45\%$};
    \draw[poporange!70!popdark, thick] (1.5, -1.3) -- (0.5, -0.2);
    
    % Organic label
    \node[popdark, font=\bfseries] at (-3.5, 1.5) {Organic Matter $\approx 5\%$};
    \draw[popdark, thick] (-2.5, 1.2) -- (-0.5, 0.5);
    
    % Water label
    \node[popblue, font=\bfseries] at (-3.5, -1.5) {Soil Water $\approx 25\%$};
    \draw[popblue, thick] (-2.5, -1.2) -- (-0.5, -0.5);
    
    % Air label
    \node[poppurple, font=\bfseries] at (3.5, 1.5) {Soil Air $\approx 25\%$};
    \draw[poppurple, thick] (2.5, 1.2) -- (0.5, 0.5);
    
    % Center label
    \node[font=\small\bfseries, text width=2.5cm, align=center] at (0,0) {Ideal Soil\\Volume Basis};
\end{tikzpicture}
\legende{Volume composition of an ideal agricultural soil (loam). Note that the solid phase (minerals + organic) occupies roughly half the volume, while the pore space is ideally shared equally between water and air.}
\end{popfigure}

\begin{definition}[Soil Composition (Volume Basis)]
In an ideal, well-structured agricultural soil (loam), the approximate volume percentages are:
\begin{itemize}
    \item \textbf{Mineral Matter ($\sim 45\%$):} Derived from the weathering of parent rock (primary minerals like quartz, feldspars; secondary minerals like clay minerals, oxides). Provides the skeleton and the reservoir of most nutrient elements (K, Ca, Mg, P, Fe, etc.).
    \item \textbf{Organic Matter ($\sim 5\%$):} Includes fresh residues, decomposing material, and \cle{humus} (see below). Though small in volume, it is the chemical and biological heart of the soil.
    \item \textbf{Soil Water ($\sim 25\%$):} The \cle{soil solution}, carrying dissolved ions (nutrients) to roots. Volume varies wildly with rainfall and drainage.
    \item \textbf{Soil Air ($\sim 25\%$):} Occupies the remaining pore space. Crucially different from atmospheric air: higher $\mathrm{CO_2}$ (from root/respiration), lower $\mathrm{O_2}$, near 100\% humidity. Essential for root respiration and aerobic microbial activity.
\end{itemize}
\end{definition}

\begin{attention}[Mass vs. Volume Percentages]
Examination questions often trick students by asking for percentages \emph{by weight} (mass). Because mineral particles are dense ($\sim 2.65\ \mathrm{g/cm^3}$) and organic matter is light ($\sim 1.3\ \mathrm{g/cm^3}$), the \textbf{mass percentage of minerals is much higher ($\sim 90\text{--}99\%$)} while organic matter drops to $\sim 1\text{--}5\%$. Always check the basis: \emph{volume} for porosity/aeration discussions, \emph{mass} for chemical analysis/fertiliser calculations.
\end{attention}

\courssub{Humus: The Chemical and Biological Engine}

Of that tiny $5\%$ organic matter volume, the most important fraction is \cle{humus}.

\begin{definition}[Humus]
\cle{Humus} is the dark, amorphous, colloidal, relatively stable end-product of the microbial decomposition of plant and animal residues. It has lost all visible cellular structure. It is not a single chemical compound but a complex mixture of high-molecular-weight organic polymers (humic acids, fulvic acids, humin).
\end{definition}

\begin{propriete}[Key Roles of Humus in Soil Fertility]
Humus is the "life insurance" of the soil. Its functions are disproportionate to its quantity:
\begin{enumerate}
    \item \textbf{Cation Exchange Capacity (CEC):} Humus colloids carry a high density of negative charges (from $\mathrm{COOH}$ and $\mathrm{OH}$ groups), giving them a CEC of $150\text{--}300\ \mathrm{cmol_c/kg}$ — far higher than most clay minerals. It acts as the primary nutrient "bank" in highly weathered tropical soils.
    \item \textbf{Water Retention:} Humus can hold up to $20\times$ its weight in water, acting as a sponge in the soil profile.
    \item \textbf{Soil Structure:} Humus binds mineral particles into stable aggregates (crumbs) via polysaccharides and fungal hyphae, creating the pore space for air and water.
    \item \textbf{pH Buffering:} The weak acidic functional groups resist pH changes, protecting roots and microbes from sudden acid/alkaline shocks.
    \item \textbf{Nutrient Source:} Slow mineralisation releases $\mathrm{N, P, S}$ and micronutrients in sync with plant demand (unlike soluble fertilisers which can leach).
    \item \textbf{Chelation:} Humic/fulvic acids form soluble complexes with Fe, Zn, Cu, Mn, preventing their precipitation as unavailable hydroxides in alkaline soils or toxicity in acid soils.
\end{enumerate}
\end{propriete}

\begin{savaistu}[Humus in Cameroon's Agro-Ecological Zones]
In the \textbf{Humid Forest Zone (South, East, Centre)}, high temperatures and rainfall drive rapid decomposition. Humus turnover is fast; without constant organic input (fallow, mulch), humus levels plummet, leaving the soil dependent on the low CEC of kaolinite/oxides. In the \textbf{Western Highlights (West, Northwest)} on volcanic ash, the \emph{andic} properties (allophane, imogolite) protect organic matter, leading to high humus accumulation (mollic/umbric epipedons) and exceptional fertility. In the \textbf{Sudano-Sahelian North}, low rainfall limits biomass production, so humus is naturally low ($<1\%$), making soils fragile and prone to crusting.
\end{savaistu}

\courssub{Soil pH: The Master Variable}

If humus is the engine, \cle{soil pH} is the dashboard warning light — it tells you instantly whether the chemical environment is friendly or hostile to plant roots and microbes.

\begin{definition}[Soil pH]
Soil pH is the negative logarithm of the hydrogen ion activity in the soil solution: $\mathrm{pH} = -\log_{10} a_{\mathrm{H}^+}$. It measures the \emph{intensity} of acidity or alkalinity, not the total quantity (which is "reserve acidity").
\end{definition}

\begin{popfigure}
\begin{tikzpicture}[scale=0.9]
    % pH Scale bar
    \draw[popline, line width=4pt] (0,0) -- (14,0);
    % Ticks and numbers
    \foreach \x/\lbl in {0/0, 1/1, 2/2, 3/3, 4/4, 5/5, 6/6, 7/7, 8/8, 9/9, 10/10, 11/11, 12/12, 13/13, 14/14} {
        \draw (\x, -0.15) -- (\x, 0.15);
        \node[below, font=\small] at (\x, -0.3) {\lbl};
    }
    % Colour zones
    \fill[poppink!60, opacity=0.6] (0, -0.15) rectangle (4.5, 0.8); % Strongly Acid
    \fill[poporange!60, opacity=0.6] (4.5, -0.15) rectangle (5.5, 0.8); % Acid
    \fill[popgreen!60, opacity=0.6] (5.5, -0.15) rectangle (7.5, 0.8); % Optimal
    \fill[popblue!60, opacity=0.6] (7.5, -0.15) rectangle (8.5, 0.8); % Alkaline
    \fill[poppurple!60, opacity=0.6] (8.5, -0.15) rectangle (14, 0.8); % Strongly Alkaline
    
    % Zone Labels
    \node[rotate=90, font=\bfseries\small, popdark] at (2.25, 1.2) {Strongly Acid};
    \node[rotate=90, font=\bfseries\small, popdark] at (5, 1.2) {Mod. Acid};
    \node[rotate=90, font=\bfseries\small, popdark] at (6.5, 1.2) {Optimal Range};
    \node[rotate=90, font=\bfseries\small, popdark] at (8, 1.2) {Mod. Alkaline};
    \node[rotate=90, font=\bfseries\small, popdark] at (11.25, 1.2) {Strongly Alkaline};
    
    % Crop Ranges (Horizontal bars above)
    % Maize: 5.5 - 7.0
    \draw[popgreen, line width=3pt] (5.5, 1.6) -- (7.0, 1.6);
    \node[above, font=\small, popgreen] at (6.25, 1.8) {Maize (5.5--7.0)};
    % Cocoa: 5.0 - 6.5
    \draw[poporange, line width=3pt] (5.0, 2.2) -- (6.5, 2.2);
    \node[above, font=\small, poporange] at (5.75, 2.4) {Cocoa (5.0--6.5)};
    % Tea: 4.5 - 5.5 (Acid loving)
    \draw[poppurple, line width=3pt] (4.5, 2.8) -- (5.5, 2.8);
    \node[above, font=\small, poppurple] at (5.0, 3.0) {Tea (4.5--5.5)};
    % Cassava: 5.0 - 7.5 (Tolerant)
    \draw[popblue, line width=3pt] (5.0, 3.4) -- (7.5, 3.4);
    \node[above, font=\small, popblue] at (6.25, 3.6) {Cassava (5.0--7.5)};
    % Oil Palm: 4.5 - 6.5
    \draw[poppink, line width=3pt] (4.5, 4.0) -- (6.5, 4.0);
    \node[above, font=\small, poppink] at (5.5, 4.2) {Oil Palm (4.5--6.5)};
    
    % Neutral point marker
    \draw[thick, popdark] (7, -0.15) -- (7, 0.15);
    \node[below, font=\bfseries\small] at (7, -0.5) {Neutral (7.0)};
\end{tikzpicture}
\legende{The soil pH scale (0--14) showing acidity/alkalinity zones and the preferred pH ranges for major Cameroonian crops. Most staple crops thrive in the slightly acidic to neutral window (pH 5.5--7.0). Tea is a notable acidophile.}
\end{popfigure}

\begin{propriete}[The pH Preference Rule]
\textbf{Most agricultural crops grow best in slightly acidic to neutral soils (pH 5.5--7.0).}
In this range:
\begin{itemize}
    \item Most essential nutrients are maximally available (see "Nutrient Availability vs pH" concept).
    \item Microbial activity (nitrification, nitrogen fixation, decomposition) is optimal.
    \item Toxic elements ($\mathrm{Al^{3+}}, \mathrm{Mn^{2+}}$) are precipitated/immobilised.
    \item Root growth is uninhibited.
\end{itemize}
\end{propriete}

\paragraph{Why are Cameroonian soils often acidic?}
In the humid tropics (South, East, Littoral, Southwest), annual rainfall ($>1500\ \mathrm{mm}$) far exceeds evapotranspiration. This creates a net downward flux of water (leaching). Basic cations ($\mathrm{Ca^{2+}}, \mathrm{Mg^{2+}}, \mathrm{K^+}, \mathrm{Na^+}$) are washed out of the profile, replaced on exchange sites by $\mathrm{H^+}$ and $\mathrm{Al^{3+}}$ from carbonic acid ($\mathrm{H_2CO_3}$) and organic acids. The result: \cle{ferrallitic soils} (Oxisols/Ultisols) with pH often $4.0\text{--}5.0$ in the topsoil.
Conversely, in the \textbf{Sahelian North} (low rainfall, high evaporation), capillary rise brings bases up; soils are neutral to alkaline (Vertisols, pH $7.0\text{--}8.5$). On the \textbf{Cameroon Volcanic Line} (Mount Cameroon, Bamboutos, Adamawa), fresh volcanic ash (basalt, trachyte) weathers to release bases rapidly, maintaining pH $6.0\text{--}7.5$ (Andosols, Vertisols) — the "breadbaskets" of the country.

\begin{methode}[Measuring Soil pH in the Field / Lab (GCE Practical)]
\textbf{Apparatus:} Soil sample, distilled water (or $\mathrm{0.01\ M\ CaCl_2}$ for lab standard), beaker, glass rod, pH meter (calibrated) \emph{or} BDH/Universal indicator solution + colour chart \emph{or} pH paper.
\textbf{Procedure (Water Method, 1:2.5 ratio):}
\begin{enumerate}
    \item Air-dry and sieve soil ($<2\ \mathrm{mm}$).
    \item Weigh $20\ \mathrm{g}$ soil into a beaker.
    \item Add $50\ \mathrm{ml}$ distilled water ($\mathrm{soil:water} = 1:2.5$).
    \item Stir vigorously for 30 seconds; let stand 30 minutes (stir occasionally).
    \item \textbf{If pH meter:} Rinse electrode, immerse in suspension, read stable value. Rinse between samples.
    \item \textbf{If indicator:} Add few drops of indicator to a small aliquot of supernatant, compare colour to chart.
    \item \textbf{If pH paper:} Dip paper into supernatant, compare.
\end{enumerate}
\textbf{Interpretation for Crop Choice:}
\begin{itemize}
    \item $\mathrm{pH} < 4.5$: Very strongly acidic. $\mathrm{Al}$ toxicity likely. Liming essential before most crops. Suitable: Tea, Pineapple.
    \item $\mathrm{pH}\ 4.5\text{--}5.5$: Strongly acidic. Good for Cocoa, Oil Palm, Rubber, Cassava. Liming needed for Maize, Legumes, Vegetables.
    \item $\mathrm{pH}\ 5.5\text{--}7.0$: Optimal for Maize, Beans, Groundnut, Tomato, Cabbage, Plantain.
    \item $\mathrm{pH} > 7.5$: Alkaline. $\mathrm{P, Fe, Zn, Mn, Cu}$ availability drops. Suitable: Cotton, Sorghum, Millet (with management).
\end{itemize}
\textbf{Note:} $\mathrm{pH_{CaCl_2}}$ is typically $0.5\text{--}1.0$ units lower than $\mathrm{pH_{water}}$ because $\mathrm{Ca^{2+}}$ displaces $\mathrm{H^+}/\mathrm{Al^{3+}}$ from exchange sites. Always state the method used.
\end{methode}

\courssub{Essential Plant Nutrients and Their Behaviour}

Plants require 17 essential elements. Carbon, Hydrogen, Oxygen come from air/water. The other 14 are mineral nutrients absorbed from the soil solution.

\begin{center}
\begin{tabularx}{\linewidth}{|l|X|X|}\hline
\rowcolor{popblueL}
\textbf{Group} & \textbf{Elements} & \textbf{Primary Role / Mobility in Plant} \\ \hline
\textbf{Primary Macros} & Nitrogen (N), Phosphorus (P), Potassium (K) & N: Proteins, chlorophyll (mobile). P: Energy (ATP), DNA (mobile). K: Osmoregulation, enzyme activation (mobile). \\ \hline
\textbf{Secondary Macros} & Calcium (Ca), Magnesium (Mg), Sulphur (S) & Ca: Cell walls, membranes (immobile). Mg: Chlorophyll centre (mobile). S: Proteins (immobile). \\ \hline
\textbf{Micronutrients (Trace)} & Fe, Mn, Zn, Cu, B, Mo, Cl, Ni & Enzyme cofactors, redox, structure. Most immobile (except Mo, Cl). \\ \hline
\end{tabularx}
\end{center}

\begin{attention}[The "Fertiliser Fallacy" — Locked Nutrients in Acid Soils]
A classic error: \emph{"My maize is yellow, I will apply more NPK fertiliser."} On a ferrallitic soil at $\mathrm{pH}\ 4.2$, added $\mathrm{P}$ instantly precipitates as $\mathrm{Fe/Al}$-phosphates (unavailable). Added $\mathrm{K}$, $\mathrm{Ca}$, $\mathrm{Mg}$ leach rapidly because the CEC is saturated with $\mathrm{Al^{3+}}/\mathrm{H^+}$. \textbf{First correct the pH (liming) to unlock the soil's own reserves and retain applied fertiliser.} Liming raises pH $\to$ precipitates $\mathrm{Al^{3+}}$ $\to$ frees CEC sites $\to$ retains $\mathrm{K^+}, \mathrm{Ca^{2+}}, \mathrm{Mg^{2+}}, \mathrm{NH_4^+}$ $\to$ increases P availability. \emph{Then} apply fertiliser.
\end{attention}

\courssub{Cation Exchange Capacity (CEC): The Nutrient Bank}

We saw in the Physical Characteristics section that clay particles are colloidal. Humus is also colloidal. Colloids carry negative charges (isomorphic substitution in clays; dissociated functional groups in humus). They attract and hold cations ($\mathrm{Ca^{2+}}, \mathrm{Mg^{2+}}, \mathrm{K^+}, \mathrm{NH_4^+}, \mathrm{H^+}, \mathrm{Al^{3+}}$) on their surfaces. This is \cle{Cation Exchange Capacity (CEC)}.

\begin{definition}[Cation Exchange Capacity (CEC)]
The \cle{Cation Exchange Capacity (CEC)} is the total quantity of exchangeable cations that a soil can adsorb per unit mass (or volume) at a given pH, usually expressed in \textbf{centimoles of charge per kilogram} ($\mathrm{cmol_c/kg}$, formerly $\mathrm{meq/100g}$).
\[
\mathrm{CEC} = \sum \text{Exchangeable Bases } (\mathrm{Ca^{2+}, Mg^{2+}, K^+, Na^+}) + \text{Exchangeable Acidity } (\mathrm{H^+, Al^{3+}})
\]
\end{definition}

\begin{propriete}[CEC of Common Soil Colloids (at pH 7)]
\begin{center}
\begin{tabular}{|l|c|}\hline
\rowcolor{popblueL}
\textbf{Colloid} & \textbf{CEC ($\mathrm{cmol_c/kg}$)} \\ \hline
Humus (mature) & $150\text{--}300$ \\ \hline
Vermiculite / Smectite (Montmorillonite) & $100\text{--}150$ \\ \hline
Illite & $20\text{--}40$ \\ \hline
Kaolinite & $3\text{--}15$ \\ \hline
Sesquioxides (Gibbsite, Goethite) & $< 5$ (variable charge) \\ \hline
Allophane / Imogolite (Volcanic) & $50\text{--}150$ (pH dependent) \\ \hline
\end{tabular}
\end{center}
\end{propriete}

\begin{aretenir}[CEC in the Cameroonian Context]
\begin{itemize}
    \item \textbf{Ferrallitic Soils (South/East):} Dominated by Kaolinite + Oxides. \textbf{Very low CEC ($< 5\ \mathrm{cmol_c/kg}$)}. Nutrients leach instantly. Fertility resides \emph{entirely} in the top $10\text{--}20\ \mathrm{cm}$ humus layer (biological cycling). "The forest feeds the soil, the soil feeds the forest."
    \item \textbf{Volcanic Soils (West/NW/SW slopes):} Allophane/Imogolite + Humus. \textbf{High CEC ($> 30\ \mathrm{cmol_c/kg}$)}. High nutrient retention. Resilient to leaching.
    \item \textbf{Vertisols (North/Logone):} Smectite. \textbf{High CEC ($> 40\ \mathrm{cmol_c/kg}$)}. But high $\mathrm{Mg/Ca}$ ratio and shrink-swell limit rooting.
    \item \textbf{Alluvial Soils (Wouri, Logone, Benue valleys):} Mixed mineralogy. Moderate to High CEC. Renewed by flooding.
\end{itemize}
\end{aretenir}

\begin{exoresolu}[Calculating Base Saturation and Liming Requirement]
\textbf{Statement:} A soil sample from a maize field in Mbankomo (South Region) has the following exchangeable cations (extracted by $\mathrm{NH_4OAc}$ pH 7):
$\mathrm{Ca^{2+}} = 2.0$, $\mathrm{Mg^{2+}} = 0.8$, $\mathrm{K^+} = 0.3$, $\mathrm{Na^+} = 0.1$, $\mathrm{Al^{3+}} = 1.5$, $\mathrm{H^+} = 0.4$ (all in $\mathrm{cmol_c/kg}$).
\begin{enumerate}
    \item Calculate the CEC and the Base Saturation (\%BS).
    \item The target \%BS for maize is $70\%$. Assuming the CEC remains constant, calculate the lime requirement ($\mathrm{CaCO_3}$ in $\mathrm{t/ha}$) to raise \%BS to $70\%$, assuming $1\ \mathrm{cmol_c/kg}$ exchangeable acidity neutralised $\approx 0.5\ \mathrm{t/ha}\ \mathrm{CaCO_3}$ (for $20\ \mathrm{cm}$ depth).
\end{enumerate}
\textbf{Detailed Solution:}
\begin{enumerate}
    \item \textbf{CEC} = Sum of all cations $= 2.0 + 0.8 + 0.3 + 0.1 + 1.5 + 0.4 = \mathbf{5.1\ cmol_c/kg}$.
    \item \textbf{Sum of Bases (SB)} $= 2.0 + 0.8 + 0.3 + 0.1 = \mathbf{3.2\ cmol_c/kg}$.
    \item \textbf{\% Base Saturation} $= \frac{SB}{CEC} \times 100 = \frac{3.2}{5.1} \times 100 = \mathbf{62.7\%}$.
    \item \textbf{Target SB} for $70\%$ BS $= 0.70 \times 5.1 = 3.57\ \mathrm{cmol_c/kg}$.
    \item \textbf{Deficit} $= 3.57 - 3.2 = 0.37\ \mathrm{cmol_c/kg}$.
    \item \textbf{Lime Requirement} $= 0.37 \times 0.5\ \mathrm{t/ha} = \mathbf{0.185\ t/ha} \approx \mathbf{185\ kg/ha}\ \mathrm{CaCO_3}$.
    \item \textbf{Interpretation:} The soil is moderately acidic (low CEC, significant $\mathrm{Al^{3+}}$). The lime requirement is relatively low because CEC is low. However, in practice, one would apply $1\text{--}2\ \mathrm{t/ha}$ to ensure pH rises sufficiently to precipitate remaining $\mathrm{Al^{3+}}$ and provide a buffer.
\end{enumerate}
\end{exoresolu}

\courssub{Soil Biology: The Living Factory}

Soil is not a dead medium; it is a living ecosystem. The \cle{soil biomass} in a healthy topsoil can exceed $5\text{--}10\ \mathrm{t/ha}$ (equivalent to a herd of cattle per hectare underground!).

\begin{center}
\begin{tabularx}{\linewidth}{|l|X|X|}\hline
\rowcolor{popblueL}
\textbf{Group} & \textbf{Key Representatives} & \textbf{Ecological / Agronomic Function} \\ \hline
\textbf{Bacteria} & \emph{Azotobacter, Rhizobium, Nitrosomonas, Nitrobacter, Bacillus} & \textbf{Decomposers} (rapid, simple C). \textbf{N-fixers} (free-living \& symbiotic). \textbf{Nitrifiers} ($\mathrm{NH_4^+ \to NO_2^- \to NO_3^-}$). Require neutral pH, good aeration. \\ \hline
\textbf{Fungi} & \emph{Aspergillus, Penicillium, Mycorrhizae (Glomus)} & \textbf{Decomposers} (slow, complex C: lignin, cellulose). \textbf{Mycorrhizae}: extend root surface $\times 100$ for P/Zn/Cu uptake. Tolerate acid pH. Produce glomalin (soil glue). \\ \hline
\textbf{Actinomycetes} & \emph{Streptomyces} & Decompose tough organics (chitin, cellulose). Produce antibiotics (suppress pathogens). "Earthy" smell (geosmin). \\ \hline
\textbf{Algae / Cyanobacteria} & \emph{Nostoc, Anabaena} & Photosynthesise on wet surface. \textbf{N-fixation} (paddy rice). Soil crust formation. \\ \hline
\textbf{Protozoa / Nematodes} & Flagellates, Amoebae; bacterial/fungal feeders & \textbf{Grazers}: regulate microbial populations, mineralise nutrients (excrete $\mathrm{NH_4^+}$). Some plant-parasitic nematodes are pests. \\ \hline
\textbf{Earthworms} & \emph{Lumbricus, Eudrilus (African nightcrawler)} & \textbf{Ecosystem Engineers}: ingest soil + OM $\to$ casts (stable aggregates, high nutrients). Burrows $\to$ macropores (drainage, aeration, root channels). "Nature's plough". \\ \hline
\textbf{Termites} & \emph{Macrotermes, Cubitermes} & Dominant in tropics/savanna. Mound builders. Decompose wood/litter. Move huge soil volumes (bioturbation). Can damage crops/structures. \\ \hline
\end{tabularx}
\end{center}

\begin{propriete}[Nutrient Cycling: The Nitrogen Cycle in Soil]
This is the most critical cycle for productivity.
\[
\ce{Organic N ->[Ammonification/Proteolysis] NH4+ ->[Nitrification] NO3-}
\]
\begin{itemize}
    \item \textbf{Mineralisation (Ammonification):} Microbes release $\mathrm{NH_4^+}$. Fast in warm, moist, neutral, aerated soils.
    \item \textbf{Nitrification:} $\mathrm{NH_4^+ \to NO_3^-}$. Acid-sensitive (stops at $\mathrm{pH} < 5.5$). $\mathrm{NO_3^-}$ is mobile (leaches); $\mathrm{NH_4^+}$ is held on CEC. Carried out by \emph{Nitrosomonas} and \emph{Nitrobacter}, requires $\mathrm{O_2}$.
    \item \textbf{Immobilisation:} Microbes compete with plants for $\mathrm{NH_4^+}/\mathrm{NO_3^-}$ if C:N ratio of residue $> 25\text{--}30$ (e.g., maize stover, sawdust). \textbf{Warning:} Fresh high-C residues cause temporary N deficiency ("N hunger").
    \item \textbf{Denitrification:} Anaerobic (waterlogged) $\to$ $\mathrm{NO_3^- \to N_2O \to N_2}$ (gas loss). Major loss pathway in flooded rice / compacted soils.
    \item \textbf{Biological N Fixation (BNF):} $\mathrm{N_2 \to NH_3}$ by \emph{Rhizobium} (legumes), \emph{Azotobacter}, cyanobacteria. Requires Mo, Fe, $\mathrm{Ca}$, neutral pH, low $\mathrm{NO_3^-}$.
\end{itemize}
\end{propriete}

\begin{experience}[Observing Soil Life: The "Cotton Strip" or "Tea Bag" Assay (Simple Field Test)]
\textbf{Objective:} Compare decomposition rates (biological activity) in different soil management plots (e.g., Forest vs. Continuous Maize vs. Maize+Mucuna cover crop).
\textbf{Materials:} Standard cotton fabric strips ($10\times 10\ \mathrm{cm}$) or commercial tea bags (green tea = labile C, rooibos = recalcitrant C), mesh bags, tent pegs, balance.
\textbf{Method:}
\begin{enumerate}
    \item Weigh dry strips/bags ($W_0$).
    \item Bury at $10\ \mathrm{cm}$ depth in target plots (3 replicates). Mark with pegs.
    \item Retrieve after 30, 60, 90 days.
    \item Wash gently, dry at $60^\circ\mathrm{C}$, weigh ($W_t$).
    \item Calculate mass loss $\% = \frac{W_0 - W_t}{W_0} \times 100$.
\end{enumerate}
\textbf{Expected Result:} Forest/Cover crop plots show highest mass loss (active microbes, mesofauna). Bare/continuous maize plots show lowest (low OM, high temp, low moisture). Links directly to \emph{humus formation} and \emph{nutrient release rates}.
\end{experience}

\courssub{Soil Water and Soil Air: The Medium of Life}

We introduced them in the pie chart; now we understand their chemical/biological roles.

\begin{propriete}[Soil Water — The Solvent and Transport Medium]
\begin{itemize}
    \item \textbf{Nutrient Transport:} Mass flow (with transpiration stream) brings $\mathrm{NO_3^-}, \mathrm{Ca^{2+}, Mg^{2+}, SO_4^{2-}}$. Diffusion (concentration gradient) brings $\mathrm{P, K, NH_4^+}$ — \emph{requires continuous water films}.
    \item \textbf{Chemical Reactions:} Weathering, dissolution/precipitation, ion exchange all occur in aqueous phase.
    \item \textbf{Plant Uptake:} Roots absorb ions only in solution.
    \item \textbf{Management:} Field Capacity (FC) $\to$ optimal air/water balance. Permanent Wilting Point (PWP) $\to$ water held too tightly ($> 1500\ \mathrm{kPa}$). Available Water Capacity (AWC) $= \mathrm{FC} - \mathrm{PWP}$. \textbf{Clay + Humus} maximise AWC.
\end{itemize}
\end{propriete}

\begin{propriete}[Soil Air — The Respiratory Gas Exchange]
\begin{itemize}
    \item \textbf{Composition:} $\mathrm{O_2}\ (15\text{--}20\%),\ \mathrm{CO_2}\ (0.5\text{--}5\%),\ \mathrm{N_2}\ (\text{rest})$. $\mathrm{CO_2}$ is $10\text{--}100\times$ atmospheric due to root/microbial respiration.
    \item \textbf{Root Respiration:} $\ce{C6H12O6 + 6O2 -> 6CO2 + 6H2O + Energy}$. Roots die in anaerobic conditions ($\mathrm{O_2} < 10\%$).
    \item \textbf{Microbial Respiration:} Aerobic decomposition $\gg$ Anaerobic (speed, energy yield, no toxic byproducts like $\mathrm{CH_4}, \mathrm{H_2S}, \mathrm{Fe^{2+}, Mn^{2+}}$).
    \item \textbf{Gas Diffusion:} Governed by \textbf{Air-Filled Porosity (AFP)}. AFP $> 10\text{--}15\%$ needed for adequate $\mathrm{O_2}$ diffusion. Compaction / Waterlogging $\to$ AFP $\to 0$ $\to$ Anaerobiosis $\to$ Root rot, Denitrification, Toxicity.
\end{itemize}
\end{propriete}

\courssub{Synthesis: Linking Chemistry, Biology and Physics to Fertility}

A GCE Advanced Level answer on soil fertility must integrate the three pillars. Here is the "mental model" for a top-mark essay or structured question.

\begin{aretenir}[The Fertility Triangle — What to Write in the Exam]
When asked "Explain the factors determining soil fertility" or "Why are volcanic soils more fertile than ferrallitic soils?", structure your answer around the interaction of:
\begin{enumerate}
    \item \textbf{Physical Platform:} Texture $\to$ Structure $\to$ Porosity $\to$ Root penetration, Water holding (AWC), Aeration (AFP). \emph{No roots = no uptake.}
    \item \textbf{Chemical Bank:} CEC (Clay type + Humus) $\to$ Nutrient retention (vs Leaching). pH $\to$ Nutrient availability / Toxicity / Microbial habitat. Base Saturation $\to$ Nutrient balance (Ca:Mg:K).
    \item \textbf{Biological Engine:} OM inputs $\to$ Humus formation $\to$ Nutrient mineralisation (N, P, S) $\to$ Structure stabilisation $\to$ Disease suppression $\to$ Mycorrhizal P uptake.
\end{enumerate}
\textbf{The Cameroonian Contrast:}
\begin{itemize}
    \item \textbf{Ferrallitic (South):} Good Physical (deep, porous) $\times$ Poor Chemical (Low CEC, Acid, Al toxic) $\times$ Biological dependent (Nutrients locked in biomass). \emph{Management:} Maintain OM cover, Agroforestry, Targeted liming + P rock phosphate.
    \item \textbf{Volcanic (West/NW):} Good Physical (Andic, high porosity) $\checkmark$ Good Chemical (High CEC, pH 6-7, High bases) $\checkmark$ Good Biological (High OM, protected). \emph{Management:} Sustain OM, prevent erosion (steep slopes), balance K/Mg/Ca.
    \item \textbf{Vertisols (North):} Poor Physical (Hard when dry, sticky wet, cracks cut roots) $\times$ Good Chemical (High CEC, High bases) $\checkmark$ Moderate Biological (Seasonal). \emph{Management:} Timely tillage (moist window), drainage, Ca/gypsum for structure.
\end{itemize}
\end{aretenir}

\begin{exoresolu}[Integrated Interpretation of a Soil Analysis Report]
\textbf{Statement:} A farmer in Dschang (West Region, volcanic soil) and a farmer in Sangmélima (South Region, ferrallitic soil) send topsoil ($0\text{--}20\ \mathrm{cm}$) samples to the lab. Results:
\begin{center}
\begin{tabular}{|l|c|c|}\hline
\rowcolor{popblueL}
\textbf{Parameter} & \textbf{Dschang (Volcanic)} & \textbf{Sangmélima (Ferrallitic)} \\ \hline
pH ($\mathrm{H_2O}$) & 6.2 & 4.3 \\ \hline
Organic C (\%) & 4.5 & 1.2 \\ \hline
Total N (\%) & 0.35 & 0.08 \\ \hline
Available P (Bray-1, mg/kg) & 25 & 3 \\ \hline
Exchangeable $\mathrm{Ca}$ ($\mathrm{cmol_c/kg}$) & 12.0 & 0.5 \\ \hline
Exchangeable $\mathrm{Mg}$ ($\mathrm{cmol_c/kg}$) & 4.0 & 0.2 \\ \hline
Exchangeable $\mathrm{K}$ ($\mathrm{cmol_c/kg}$) & 0.8 & 0.1 \\ \hline
Exchangeable $\mathrm{Al}$ ($\mathrm{cmol_c/kg}$) & 0.0 & 1.8 \\ \hline
CEC ($\mathrm{cmol_c/kg}$) & 28.0 & 4.5 \\ \hline
Base Saturation (\%) & 60 & 18 \\ \hline
\end{tabular}
\end{center}
\textbf{Questions:}
\begin{enumerate}
    \item Compare the chemical fertility status of the two soils.
    \item Explain the role of pH and CEC in the observed Available P difference.
    \item Recommend specific management practices for \emph{each} soil for maize cultivation.
\end{enumerate}
\textbf{Detailed Solution:}
\begin{enumerate}
    \item \textbf{Comparison:}
        \begin{itemize}
            \item \textbf{Dschang:} Near-neutral pH (optimal), High OM/N, High CEC, High bases (Ca, Mg, K), Zero $\mathrm{Al}$, Good BS (60\%). \textbf{High inherent fertility.} Limiting factor might be P fixation by allophane (though Bray-1 shows 25 mg/kg is adequate) or K/Mg balance.
            \item \textbf{Sangmélima:} Strongly acidic, Low OM/N, Very Low CEC, Very Low bases, High $\mathrm{Al}$ (toxic!), Very Low BS (18\%). \textbf{Very low inherent fertility.} Fertility entirely dependent on the thin humus layer. High $\mathrm{Al}$ saturation ($= 1.8/4.5 = 40\%$) explains root stunting.
        \end{itemize}
    \item \textbf{pH/CEC vs P Availability:}
        \begin{itemize}
            \item \textbf{Sangmélima (pH 4.3):} High $\mathrm{Fe/Al}$ oxides (sesquioxides) dominate. At low pH, oxide surfaces are positively charged ($\mathrm{pH} < \mathrm{PZC}$) and/or $\mathrm{Al^{3+}}$ in solution precipitates P as $\mathrm{AlPO_4}$ (variscite) / $\mathrm{FePO_4}$ (strengite). Extremely strong fixation $\to$ Available P = 3 mg/kg (deficient).
            \item \textbf{Dschang (pH 6.2):} Allophane/Imogolite have variable charge. At pH 6.2, they retain anions (P) but less intensely than crystalline oxides at low pH. High OM competes for fixation sites. Organic acids chelate $\mathrm{Al/Fe}$. Result: Moderate fixation, Available P = 25 mg/kg (sufficient).
        \end{itemize}
    \item \textbf{Management Recommendations:}
        \begin{itemize}
            \item \textbf{Sangmélima (Ferrallitic):} \textbf{URGENT: LIMING.} Apply $1\text{--}2\ \mathrm{t/ha}\ \mathrm{CaCO_3}$ (or dolomite for Mg) to raise pH $> 5.5$, precipitate $\mathrm{Al}$, increase CEC (variable charge), boost microbial activity. \textbf{P SOURCE:} Use Reactive Phosphate Rock (RPR) — cheaper, slow release, suited to acid soils; or TSP/SSP but band-placed near seed to reduce fixation volume. \textbf{OM:} \emph{Mucuna}/\emph{Crotalaria} cover crops, compost, mulch to rebuild humus/CEC. \textbf{K:} Split applications (leaching risk). \textbf{Agroforestry:} \emph{Acacia, Leucaena} for nutrient pumping.
            \item \textbf{Dschang (Volcanic):} \textbf{MAINTENANCE.} No lime needed (pH 6.2). Monitor pH drop from $\mathrm{NH_4^+}$ fertilisers/leaching. \textbf{P:} Standard TSP/NPK placement. \textbf{K/Mg Balance:} High Mg (4.0) vs K (0.8) $\to$ Mg:K ratio $= 5:1$ (ideal $\sim 2\text{--}3:1$). Risk of K deficiency / Grass tetany in livestock. Apply $\mathrm{K_2SO_4}$ or $\mathrm{KCl}$ (avoid excess Cl if sensitive crops). \textbf{Erosion Control:} Contour ridges, hedgerows (steep slopes). \textbf{OM:} Return residues, compost.
        \end{itemize}
\end{enumerate}
\end{exoresolu}

\begin{exercice}[Chapter Consolidation Exercise]
\begin{enumerate}
    \item Define Cation Exchange Capacity (CEC) and state its unit. Explain why humus has a higher CEC per unit mass than kaolinite.
    \item A soil has a pH of 4.8. Explain the chemical forms in which Nitrogen, Phosphorus, and Aluminium are likely to exist in the soil solution, and the consequences for a maize crop.
    \item Describe the role of earthworms and mycorrhizal fungi in maintaining soil chemical fertility and physical structure.
    \item A farmer applies $200\ \mathrm{kg/ha}$ of NPK 20-10-10 to a maize field on a ferrallitic soil (pH 4.2, CEC $3\ \mathrm{cmol_c/kg}$) without liming. Predict the fate of the N, P, and K applied. Justify your answer.
    \item Using the concept of Base Saturation, explain why volcanic soils of the Cameroon Volcanic Line generally support higher crop yields than ferrallitic soils of the Southern Plateau, even under similar rainfall.
\end{enumerate}
\end{exercice}

\begin{corrige}[Chapter Consolidation Exercise]
\begin{enumerate}
    \item \textbf{CEC Definition:} Total quantity of exchangeable cations a soil can hold per unit mass at a given pH. Unit: $\mathrm{cmol_c/kg}$ (or $\mathrm{meq/100g}$). \textbf{Why Humus > Kaolinite:} Humus has abundant $\mathrm{-COOH}$ and $\mathrm{-OH}$ groups (pH-dependent charge) giving $150\text{--}300\ \mathrm{cmol_c/kg}$. Kaolinite has low isomorphic substitution and only edge charges ($\mathrm{-OH}$), giving $3\text{--}15\ \mathrm{cmol_c/kg}$.
    \item \textbf{pH 4.8 Scenario:}
        \begin{itemize}
            \item \textbf{N:} Nitrification inhibited ($\mathrm{pH} < 5.5$). N remains as $\mathrm{NH_4^+}$ (held on CEC, less leaching but less mobile for roots). Urea/ammonium fertilisers preferred over nitrate.
            \item \textbf{P:} Fixed as $\mathrm{AlPO_4}$ / $\mathrm{FePO_4}$ (very insoluble). Low availability. Root growth restricted by $\mathrm{Al}$ $\to$ less P foraging.
            \item \textbf{Al:} $\mathrm{Al^{3+}}$ dominant soluble form (toxic). Damages root tips $\to$ stubby roots $\to$ water/nutrient uptake failure. Maize is Al-sensitive.
        \end{itemize}
    \item \textbf{Earthworms:} Ingest mineral+OM $\to$ casts: higher nutrient availability (mineralised), stable aggregates (structure), macropores (aeration/drainage/root channels). \textbf{Mycorrhizae:} Hyphae extend depletion zone $\to$ access immobile P, Zn, Cu. Glomalin $\to$ aggregate stability. Receive plant C $\to$ symbiosis.
    \item \textbf{Fate of NPK on Ferrallitic pH 4.2:}
        \begin{itemize}
            \item \textbf{N (20\% = 40 kg N):} If urea/ammonium $\to$ stays $\mathrm{NH_4^+}$ (held on low CEC, some volatilisation as $\mathrm{NH_3}$ if surface applied). If nitrate $\to$ leaches immediately (no $\mathrm{NO_3^-}$ retention). Low recovery.
            \item \textbf{P (10\% = 20 kg $\mathrm{P_2O_5}$ $\approx$ 8.7 kg P):} Instant fixation by $\mathrm{Al/Fe}$ oxides $\to$ $< 10\%$ recovery first season. Residual builds "P bank" but unavailable.
            \item \textbf{K (10\% = 20 kg $\mathrm{K_2O}$ $\approx$ 16.6 kg K):} Low CEC ($3\ \mathrm{cmol_c/kg}$) saturated with $\mathrm{Al/H}$ $\to$ $\mathrm{K^+}$ cannot be held $\to$ leaches rapidly with high rainfall. Very low recovery.
        \end{itemize}
        \textbf{Conclusion:} Waste of money. Lime first.
    \item \textbf{Base Saturation Argument:}
        \begin{itemize}
            \item \textbf{Volcanic:} Parent material (basalt/trachyte) rich in Ca, Mg, K. Weathering releases bases $\to$ high exchangeable bases $\to$ High BS ($> 50\text{--}80\%$). pH buffered near neutral. Nutrients available, low toxicity.
            \item \textbf{Ferrallitic:} Intense leaching (high rainfall, old landscapes) removes bases $\to$ exchange complex dominated by $\mathrm{H^+}, \mathrm{Al^{3+}}$ $\to$ Low BS ($< 20\%$). Low pH $\to$ Al toxicity, P fixation, low microbial activity.
            \item \textbf{Result:} Volcanic soils have a large "chemical buffer" (reserve of bases on CEC) supplying crops continuously. Ferrallitic soils have an empty "bank" — crops survive only on the "cash flow" of rapid nutrient cycling from surface OM.
        \end{itemize}
\end{enumerate}
\end{corrige}

\begin{experience}[Case Study: The Lake Nyos Disaster — A Soil/Gas Link (Extension)]
\textbf{Context:} Lake Nyos (Northwest Region) is a crater lake on the Cameroon Volcanic Line. In 1986, a limnic eruption released $\sim 1.6\ \mathrm{million\ tonnes}$ of $\mathrm{CO_2}$, killing 1700+ people and livestock.
\textbf{Geochemical Link:} Magmatic $\mathrm{CO_2}$ dissolves in deep lake water under pressure. The lake is meromictic (stratified). The deep water becomes saturated. A trigger (landslide, rain) causes overturn $\to$ pressure release $\to$ explosive degassing.
\textbf{Soil Relevance:} The surrounding volcanic soils (Andosols) are highly permeable. $\mathrm{CO_2}$ diffuses from soil into homes (basements) — a chronic hazard. Soil gas monitoring ($\mathrm{CO_2}, \mathrm{Rn}$) is now part of risk management. This illustrates the \emph{soil air} component connecting deep geology to surface safety.
\end{experience}

\coursec{Soil Usage, Sustainable Management and Protection}

Having examined the physical, chemical and biological characteristics that define a soil's potential, we now turn to the critical interface between soil and society. A soil profile is not merely a scientific curiosity; it is the foundation of Cameroon's economy, feeding millions through agriculture, providing raw materials for construction, and regulating our water and climate. Yet this foundation is fragile. In this section, we analyse how we use soil, why it is effectively a non-renewable resource on human timescales, the specific threats it faces in our national context, and the science-based strategies for its sustainable management and protection.

\courssub{Economic Uses of Soil}

Soil is the primary capital of the primary sector. In Cameroon, where agriculture employs over 60\% of the active population and contributes significantly to GDP, the economic value of soil is direct and immense.

\begin{itemize}
    \item \textbf{Food and Cash Crop Agriculture:} The volcanic soils of the Cameroon Volcanic Line (CVL) — particularly the andosols around Mount Cameroon, the Bamboutos Mountains, and the Western Highlands — support intensive cultivation of \emph{cocoa}, \emph{coffee}, \emph{banana}, \emph{plantain}, and \emph{maize}. In the sedimentary basin (Douala, Kribi) and the Benue trough, alluvial soils support rice, oil palm, and rubber. The ferralitic soils of the South Cameroon Plateau, though acidic and leached, are managed for cocoa and robusta coffee. In the North, the vertisols and tropical ferruginous soils of the Benue and Logone plains are the cotton and sorghum belts.
    \item \textbf{Livestock Grazing:} The savannah zones (Adamawa, North, Far North) rely on natural pastures growing on shallow, often ferruginous or vertic soils. Transhumance patterns are dictated by soil moisture retention and grass nutritional quality, which are soil-dependent properties.
    \item \textbf{Construction Materials:} Lateritic horizons (ironstone gravel, \emph{durcrust}) are excavated for road sub-base and surfacing. Clay-rich horizons (often Bt horizons in ferralitic soils or alluvial clays) are the raw material for \emph{brick making} and \emph{pottery}, a major informal sector activity in towns like Foumban, Baham, and Maroua.
    \item \textbf{Sand and Gravel Extraction:} Alluvial deposits along the Wouri, Sanaga, Benue, and Logone rivers provide sand and gravel for the booming construction industry in Yaoundé, Douala, and secondary cities.
    \item \textbf{Engineering Foundation:} Every building, road, bridge, and dam rests on soil or rock. Understanding the bearing capacity, shrink-swell potential (critical for vertisols), and compressibility of the subsoil is a prerequisite for safe, durable infrastructure.
\end{itemize}

\begin{savaistu}[Cameroon's "Brown Gold"]
Cocoa and coffee, historically Cameroon's main export earners, are essentially "packaged soil fertility". Every tonne of cocoa beans exported represents a net export of nutrients (N, P, K, Mg) that were extracted from the soil. Without replenishment, this is soil mining, not farming.
\end{savaistu}

\courssub{Soil as a Non-Renewable Resource}

A fundamental concept in economic geology and environmental science is the \textbf{timescale of formation} versus the \textbf{timescale of degradation}.

\begin{propriete}[Soil Formation Timescale]
Under favourable humid tropical conditions, forming \SI{1}{cm} of topsoil (A horizon) takes approximately \textbf{100 to 400 years}. In drier or cooler climates, it takes significantly longer. A deep, mature soil profile (1–2 m) represents \textbf{millennia} of pedogenesis.
\end{propriete}

\begin{aretenir}[Human Timescale vs Geological Timescale]
On a human timescale (decades to a few centuries), soil is \textbf{effectively non-renewable}. Once the topsoil is lost to erosion or sealed by urbanisation, it cannot be replaced within the lifetime of a civilisation. This contrasts with forests (decades) or water (annual cycle). Soil conservation is therefore not optional; it is an existential necessity.
\end{aretenir}

\begin{definition}[Soil Erosion]
\textbf{Soil erosion} is the detachment, transport, and deposition of soil particles by agents such as water (rainfall, runoff, streams) or wind, accelerated by human activities that remove protective vegetation cover.
\end{definition}

\courssub{Threats to Soils in Cameroon}

Cameroon's diverse agro-ecological zones face a specific suite of threats, often acting in synergy.

\begin{center}
\begin{tabularx}{\linewidth}{|l|X|X|}
\hline
\rowcolor{popblueL}
\textbf{Threat} & \textbf{Main Affected Zones} & \textbf{Mechanism / Driver} \\
\hline
Water Erosion (Sheet, Rill, Gully) & Western Highlands, Adamawa, South-West slopes (Mount Cameroon), Mandara Mountains & Steep slopes + intense rainfall + deforestation / clean weeding \\
\hline
Wind Erosion & Far North (Sahelian zone), parts of North & Dry harmattan winds + bare soil after harvest / overgrazing \\
\hline
Deforestation & South Cameroon Plateau, East Region, Coastal plain & Logging, slash-and-burn agriculture, fuelwood collection, plantation expansion \\
\hline
Overgrazing & Far North, North, Adamawa (dry season) & Stocking rates $>$ carrying capacity; trampling destroys structure, removes cover \\
\hline
Bush Fires & Savannah zones (Adamawa, North, Far North, parts of East) & Annual burning for hunting, pasture renewal; destroys organic matter, soil fauna \\
\hline
Monoculture / Over-cultivation & Cocoa/coffee belts (South-West, Centre, South, West), Cotton basin (North) & Nutrient mining, no fallow, structure collapse, pest buildup \\
\hline
Chemical Pollution & Industrial zones (Douala, Limbe), intensive plantations (banana, oil palm) & Agrochemical misuse (pesticides, excess fertiliser), industrial effluents, oil spills \\
\hline
Urbanisation / Soil Sealing & Yaoundé, Douala, Bafoussam, Bamenda, Maroua & Permanent covering with impermeable surfaces (concrete, asphalt); total loss of soil functions \\
\hline
\end{tabularx}
\end{center}

\begin{popfigure}
\legende{Cause-and-effect flow chart of the deforestation–poverty trap in the Cameroonian highlands.}
\begin{tikzpicture}[
    node distance=1.2cm and 1.5cm,
    box/.style={draw=popdark, thick, rounded corners, fill=popblueL, text width=2.8cm, align=center, minimum height=1.2cm, font=\small},
    arrow/.style={-stealth, thick, popdark}
]
\node[box, align=center] (defor){Deforestation\\(Logging, Farming, Fuelwood)};
\node[box, right=of defor, align=center] (erosion){Accelerated Erosion\\(Loss of A \& E horizons)};
\node[box, right=of erosion, align=center] (topsoil){Loss of Topsoil\\(\& Nutrients, OM, Structure)};
\node[box, below=of topsoil, align=center] (yield){Declining Yields\\(Food Insecurity)};
\node[box, left=of yield, align=center] (poverty){Rural Poverty\\(Low Income)};
\node[box, left=of poverty, align=center] (clearing){More Land Clearing\\(Marginal/Steep Lands)};

\draw[arrow] (defor) -- (erosion);
\draw[arrow] (erosion) -- (topsoil);
\draw[arrow] (topsoil) -- (yield);
\draw[arrow] (yield) -- (poverty);
\draw[arrow] (poverty) -- (clearing);
\draw[arrow] (clearing) |- (defor); % Feedback loop
\end{tikzpicture}
\end{popfigure}

\courssub{Consequences of Soil Degradation}

The degradation of the soil resource has cascading environmental and socio-economic impacts.

\begin{itemize}
    \item \textbf{Loss of Fertility and Declining Yields:} Removal of the A horizon (rich in OM, nutrients, biota) exposes the B horizon (lower OM, higher clay, possible Al toxicity, lower CEC). Yields drop sharply; farmers respond by clearing more forest (extensification), perpetuating the cycle.
    \item \textbf{Siltation of Rivers and Reservoirs:} Eroded sediment loads rivers (Sanaga, Benue, Wouri, Logone). This reduces reservoir capacity (e.g., Lagdo, Song Loulou, Memve'ele), increases flood risk downstream, raises water treatment costs, and destroys aquatic habitats.
    \item \textbf{Desertification in the Far North:} The combination of rainfall decline, overgrazing, wind erosion, and loss of vegetation cover drives the southward encroachment of desert-like conditions. The "Sahelianisation" of the Far North is a direct manifestation of soil degradation.
    \item \textbf{Landslides and Mass Wasting:} On steep volcanic slopes (Mount Cameroon, Bamboutos, Manengouba), deforestation and uncontrolled cultivation reduce root cohesion and increase pore water pressure, triggering rotational slides and debris flows, threatening lives and infrastructure (e.g., the 2003 Limbe landslides).
    \item \textbf{Carbon Loss and Climate Feedback:} Soil organic carbon (SOC) oxidation releases CO$_2$. Degraded soils sequester less carbon, turning a potential sink into a source.
\end{itemize}

\courssub{Sustainable Soil Management Techniques}

Sustainable management aims to reduce erosion to \emph{geological rates} (tolerable loss $\approx$ formation rate, typically $< \SI{10}{t.ha^{-1}.yr^{-1}}$) while maintaining productivity. The choice of technique depends on the \textbf{slope}, \textbf{climate}, \textbf{soil type}, and \textbf{farming system}.

\begin{definition}[Sustainable Soil Management]
\textbf{Sustainable soil management} is the use of soil resources in a way that maintains their productive capacity, environmental functions (water regulation, carbon storage, biodiversity habitat), and resilience for present and future generations.
\end{definition}

\begin{methode}[Matching Conservation Technique to the Problem]
\begin{enumerate}
    \item \textbf{Diagnose the dominant degradation process:} Water erosion? Wind erosion? Fertility decline? Compaction? Salinisation?
    \item \textbf{Assess the landscape context:} Slope steepness, rainfall intensity, soil texture/depth, land tenure.
    \item \textbf{Select the appropriate measure(s):}
    \begin{itemize}
        \item \textbf{Steep slopes ($> 15\%$):} \textbf{Terracing} (bench, fanya juu) + \textbf{Agroforestry} (permanent tree cover).
        \item \textbf{Gentle to moderate slopes ($5–15\%$):} \textbf{Contour ploughing}, \textbf{Contour bunds} (stone lines, earth bunds), \textbf{Grass strips} (vetiver, elephant grass).
        \item \textbf{Dry, windy zones (Far North):} \textbf{Windbreaks} (rows of \emph{Acacia}, \emph{Neem}, \emph{Casuarina}), \textbf{Mulching}, \textbf{Residue retention}.
        \item \textbf{Moisture conservation (all zones):} \textbf{Mulching} (crop residues, grass), \textbf{Cover crops} (Mucuna, Pueraria, groundnut), \textbf{Minimum tillage}.
        \item \textbf{Fertility restoration:} \textbf{Crop rotation} (cereal–legume), \textbf{Fallowing} (improved fallows with leguminous trees/shrubs), \textbf{Agroforestry} (alley cropping).
    \end{itemize}
    \item \textbf{Combine measures:} Physical structures (terraces) + Agronomic practices (mulching, rotation) + Vegetative barriers (hedgerows) = \textbf{Integrated Watershed Management}.
\end{enumerate}
\end{methode}

\begin{popfigure}
\legende{Cross-section of a cultivated slope: (a) Unprotected up-and-down ploughing leading to gully erosion; (b) Contour ploughing with tied ridges; (c) Bench terracing with riser vegetation.}
\begin{tikzpicture}[scale=0.9, every node/.style={font=\small}]
% --- (a) Unprotected ---
\begin{scope}[xshift=0cm]
\draw[popdark, thick] (0,0) -- (0,4) -- (6,4) -- (6,0) -- cycle; % Box
\node[popdark, font=\bfseries] at (3,3.6) {(a) Unprotected: Up-and-Down Ploughing};
% Slope surface
\draw[poporange, thick] (0.5,3.5) -- (5.5,0.5);
% Furrows (vertical lines = up-down)
\foreach \x in {1,1.8,2.6,3.4,4.2,5} {
    \draw[popmuted] (\x,3.3) -- (\x-0.3,0.7);
}
% Gully forming
\draw[poporange, line width=2pt] (3,3.3) .. controls (3.2,2.5) and (3.1,1.5) .. (3,0.7);
\draw[poporange, line width=1.5pt] (3.2,3.3) .. controls (3.4,2.5) and (3.3,1.5) .. (3.2,0.7);
\node[poporange, align=center] at (4.5,2.5) {Gully\\Erosion};
\draw[->, poporange] (4.5,2.5) -- (3.3,1.8);
% Runoff arrows
\draw[->, popblue, thick] (1,3.3) -- (1,2.5);
\draw[->, popblue, thick] (4,3.3) -- (4,2.5);
\node[popblue] at (0.5,2.5) {Fast Runoff};
\end{scope}

% --- (b) Contour Ploughing ---
\begin{scope}[xshift=7.5cm]
\draw[popdark, thick] (0,0) -- (0,4) -- (6,4) -- (6,0) -- cycle;
\node[popdark, font=\bfseries] at (3,3.6) {(b) Contour Ploughing + Tied Ridges};
% Slope surface
\draw[popgreen, thick] (0.5,3.5) -- (5.5,0.5);
% Furrows (horizontal = contour)
\foreach \y in {3.2, 2.6, 2.0, 1.4, 0.8} {
    \draw[popgreen] (0.7,\y) -- (5.3,\y);
    % Tied ridges (small cross ties)
    \foreach \x in {1.5, 2.5, 3.5, 4.5} {
        \draw[popgreen] (\x,\y) -- (\x,\y-0.15);
    }
}
\node[popgreen, align=center] at (4.5,2.5) {Water\\Infiltrates};
\draw[->, popgreen] (4.5,2.5) -- (3.5,1.8);
\draw[->, popblue, thick] (3,3.5) -- (3,3.0) node[midway, right] {Slow Runoff};
\end{scope}

% --- (c) Bench Terracing ---
\begin{scope}[xshift=15cm]
\draw[popdark, thick] (0,0) -- (0,4) -- (6,4) -- (6,0) -- cycle;
\node[popdark, font=\bfseries] at (3,3.6) {(c) Bench Terracing (Steep Slopes)};
% Terraces
\foreach \y/\h in {3.2/0.5, 2.4/0.5, 1.6/0.5, 0.8/0.5} {
    \fill[popgreenL] (0.5,\y) rectangle (5.5,\y+\h);
    \draw[popgreen, thick] (0.5,\y) -- (5.5,\y); % Tread
    \draw[popgreen, thick] (0.5,\y+\h) -- (5.5,\y+\h); % Next tread
}
% Risers
\foreach \y in {3.2, 2.4, 1.6, 0.8} {
    \draw[popgreen, thick] (0.5,\y) -- (0.5,\y+0.5);
    \draw[popgreen, thick] (5.5,\y) -- (5.5,\y+0.5);
}
% Vegetation on risers
\foreach \y in {3.2, 2.4, 1.6} {
    \foreach \x in {0.6, 1.2, 1.8, 2.4, 3.0, 3.6, 4.2, 4.8, 5.4} {
        \node[popgreen] at (\x,\y+0.25) {\tiny $\triangle$};
    }
}
\node[popgreen, align=center] at (4.5,2.5) {Level Treads,\\Vegetated Risers};
\draw[->, popgreen] (4.5,2.5) -- (3.5,1.8);
\end{scope}
\end{tikzpicture}
\end{popfigure}

\courssub{Maintaining Soil Fertility}

Physical conservation (stopping erosion) must be paired with chemical/biological fertility management.

\begin{itemize}
    \item \textbf{Organic Amendments (Compost, Manure, Green Manure):} The cornerstone of tropical soil fertility. They replenish SOC, feed soil biota, improve structure (aggregation), increase CEC, and buffer pH.
    \begin{itemize}
        \item \emph{Composting:} Controlled decomposition of crop residues, household waste, animal manure. C/N ratio $\approx 30:1$ for efficient decomposition.
        \item \emph{Green Manure / Cover Crops:} \emph{Mucuna pruriens}, \emph{Pueraria phaseoloides}, \emph{Crotalaria}, \emph{Tephrosia}. Fix N, smother weeds, produce biomass for mulch.
    \end{itemize}
    \item \textbf{Mineral Fertilisers (Balanced Use):} Essential to replace nutrients exported in harvests. The "4R" Nutrient Stewardship: \textbf{Right Source}, \textbf{Right Rate}, \textbf{Right Time}, \textbf{Right Place}.
    \begin{itemize}
        \item In Cameroon, NPK formulations (e.g., 20-10-10, 15-15-15) and Urea are common.
        \item \emph{Warning:} Exclusive use of acidifying fertilisers (Urea, Ammonium Sulphate, DAP) on already acidic ferralitic soils accelerates Al/Mn toxicity and base leaching. \textbf{Liming} (CaCO$_3$, Dolomite, Calcium Silicate) is mandatory to raise pH to 5.5–6.5 for most crops.
    \end{itemize}
    \item \textbf{Integrated Soil Fertility Management (ISFM):} The synergy: \emph{Organic matter} improves the \emph{efficiency} of mineral fertiliser (less leaching, better root access) and mineral fertiliser boosts the \emph{biomass production} needed for organic inputs. Neither alone is sufficient at scale.
\end{itemize}

\begin{attention}[GCE Trap: "List vs Explain"]
A common mistake in Paper 2 (Essay) is listing techniques: "Contour ploughing, terracing, mulching, crop rotation..." \textbf{This gains few marks.} You must explain the \textbf{MECHANISM} for each:
\begin{itemize}
    \item \textbf{Contour ploughing:} Furrows act as mini-dams, reducing runoff velocity, increasing infiltration time, trapping sediment.
    \item \textbf{Mulching:} Absorbs raindrop kinetic energy (prevents splash erosion), reduces evaporation, moderates temperature, adds OM on decomposition.
    \item \textbf{Agroforestry (Alley cropping):} Deep tree roots recycle leached nutrients (nutrient pump), prunings provide mulch/N, hedgerows act as physical barriers to runoff.
    \item \textbf{Liming:} Neutralises H$^+$ and Al$^{3+}$, raises pH, increases base saturation, flocculates clay (improves structure).
\end{itemize}
\end{attention}

\courssub{Soil Protection Policies and Institutions in Cameroon}

Technical solutions require an enabling policy environment.

\begin{itemize}
    \item \textbf{Legal Framework:} Law No. 96/12 of 5 August 1996 on Environmental Management (EIA mandatory for large projects). Law No. 74/1 (Land Tenure) — state ownership implies state duty of protection. Forestry Law (1994) — protects watershed forests.
    \item \textbf{MINADER (Ministry of Agriculture and Rural Development):} Extension services (SDR, SDDA), promotion of ISFM, improved seeds, fertiliser subsidy programmes (PAPEA), support for agroforestry and conservation agriculture.
    \item \textbf{MINEPDED (Ministry of Environment, Protection of Nature and Sustainable Development):} National coordination of UNCCD (Desertification), UNFCCC (Climate), CBD (Biodiversity). National Action Programme to Combat Desertification (PAN/LCD). Reforestation campaigns (Operation Green Sahel).
    \item \textbf{MINFOF (Ministry of Forestry and Wildlife):} Management of protected areas (National Parks, Reserves) which protect headwater soils. Community forests.
    \item \textbf{IRAD (Institute of Agricultural Research for Development):} Soil survey, mapping, fertility trials, development of adapted technologies (e.g., \emph{Mucuna} fallows, acid-tolerant crop varieties).
    \item \textbf{Local Governance:} Decentralisation laws transfer natural resource management to Councils. Community-based watershed management committees (e.g., in the Upper Mifi, Menoua, Logone basins).
\end{itemize}

\begin{savaistu}[The "Green Sahel" and the Great Green Wall]
Cameroon is a signatory to the \textbf{Great Green Wall for the Sahara and the Sahel Initiative (GGWSSI)}. In the Far North, this translates into large-scale planting of \emph{Acacia senegal} (gum arabic), \emph{Balanites aegyptiaca}, and \emph{Ziziphus mauritiana} — species that stabilize dunes, fix nitrogen, and provide non-timber forest products (gum, fruit, fodder) for local income. It links soil protection directly to livelihoods.
\end{savaistu}

\courssub{Advantages of Sustainable Soil Management}

Investing in soil health yields multiple dividends — the "triple win" of Climate-Smart Agriculture.

\begin{aretenir}[Advantages of Sustainable Soil Management]
\begin{enumerate}
    \item \textbf{Sustained / Increased Yields:} Healthy soils produce more reliably, year after year, reducing the need for extensification into forests.
    \item \textbf{Food and Nutrition Security:} Higher, stable yields $\rightarrow$ availability, access, utilisation. Diversified rotations improve diet quality.
    \item \textbf{Water Resource Protection:} High infiltration $\rightarrow$ groundwater recharge $\rightarrow$ sustained baseflow in rivers (dry season water supply). Low sedimentation $\rightarrow$ longer dam life, cleaner water, lower treatment costs.
    \item \textbf{Climate Change Mitigation:} Soils are the largest terrestrial carbon pool ($\approx 2500$ Gt C). Increasing SOC by $0.4\%$ per year (the "4 per 1000" initiative) could offset a significant fraction of anthropogenic CO$_2$ emissions.
    \item \textbf{Climate Change Adaptation:} Soils with high OM and structure hold more plant-available water, buffering crops against drought spells (increasingly frequent in the Sudano-Sahelian zone).
    \item \textbf{Biodiversity Conservation:} Soil biota (earthworms, termites, microbes, mycorrhizae) thrive under low-disturbance, high-OM systems.
    \item \textbf{Economic Resilience:} Reduced input costs (less fertiliser needed per unit yield), reduced risk of crop failure, diversified income (agroforestry products).
\end{enumerate}
\end{aretenir}

\begin{exemplebox}[Worked Exercise: Diagnosing a Degraded Farm in the Western Highlands]
\textbf{Scenario:} A farmer in Bafou (West Region) cultivates maize and beans on a $25\%$ slope. He ploughs up-and-down the slope with a hoe, burns residues, uses no fertiliser or manure. Yields have dropped from 3 t/ha to 0.8 t/ha in 10 years. Deep gullies are cutting across the field.

\textbf{Question:} Propose an integrated rehabilitation plan, justifying each measure.

\textbf{Solution:}
\begin{enumerate}
    \item \textbf{Immediate Physical Stabilisation:} Construct \textbf{bench terraces} (or \emph{fanya juu} terraces if labour-limited) along the contour. \emph{Why:} Slope $> 15\%$ makes contour ploughing insufficient; terraces reduce effective slope length to near zero, stopping gully progression. Stabilise risers with \emph{Vetiver grass} (deep roots, stiff stems) or \emph{Pennisetum purpureum} (Napier grass — also fodder).
    \item \textbf{Runoff Management:} Install a \textbf{cut-off drain} (diversion ditch) above the field to divert external runoff. Grassed waterways for safe disposal of excess runoff from terraces.
    \item \textbf{Agronomic Rehabilitation:}
        \begin{itemize}
            \item \textbf{Minimum tillage / Planting pits (Zai):} Disturb soil only at planting stations to preserve structure and OM.
            \item \textbf{Legume Rotation / Intercrop:} Maize – \emph{Mucuna} relay crop, or Maize + \emph{Phaseolus} beans (climbing varieties). \emph{Mucuna} fixes $\approx 100–150$ kg N/ha, suppresses \emph{Imperata} grass.
            \item \textbf{Mulching:} Retain all \emph{Mucuna} biomass and maize stover as mulch ($\geq 5$ t/ha). Restores OM, feeds termites/earthworms to rebuild macroporosity.
        \end{itemize}
    \item \textbf{Agroforestry Integration:} Plant contour hedgerows of \emph{Calliandra calothyrsus} or \emph{Leucaena leucocephala} on terrace risers/edges. Prune 3–4 times/year for mulch/fodder. Deep roots capture leached nutrients.
    \item \textbf{Soil Amendment:} Apply \textbf{dolomitic lime} (CaMg(CO$_3$)$_2$) at 1–2 t/ha based on soil test (pH likely $< 4.5$, Al saturation high). Follow with \textbf{balanced NPK + Manure} (e.g., 5 t/ha compost + 200 kg/ha NPK 20-10-10). Lime first, wait 2–3 weeks, then fertilise.
    \item \textbf{Social/Economic:} Secure land tenure (long-term lease or title) to incentivise investment. Access to credit for terrace construction tools (A-frame, hoe, pickaxe). Farmer field school for peer learning.
\end{enumerate}
\textbf{Expected Outcome:} Erosion reduced to $< 10$ t/ha/yr. SOC rebuilt from $< 1\%$ to $> 2.5\%$ in 3–5 years. Maize yield restored to $> 3$ t/ha sustainably.
\end{exemplebox}

\begin{exoresolu}[GCE Style Essay Question]
\textbf{Statement:} "Soil is a non-renewable resource that is under severe threat in Cameroon. Discuss the major causes of soil degradation in the Far North Region and outline a sustainable management strategy for the region."

\tcblower

\textbf{Detailed Solution:}

\textbf{Introduction:} Define soil degradation. State the non-renewable nature (1 cm = 100–400 years). Identify Far North as a high-risk zone (Sahelian fringe).

\textbf{1. Major Causes of Degradation in the Far North:}
\begin{itemize}
    \item \textbf{Climatic Aridity:} Low ($< 800$ mm), erratic rainfall; high evapotranspiration; harmattan winds. Low biomass production $\rightarrow$ low OM return.
    \item \textbf{Water Erosion:} High-intensity storms on bare, crusted soils (surface sealing on sandy/loamy soils). Rill/gully formation on slopes of the Mandara Mountains and diamictons.
    \item \textbf{Wind Erosion:} Harmattan (Nov–Mar) on bare fields after harvest. Loss of fine silt/clay/OM (most fertile fraction). Dune reactivation (Logone plain).
    \item \textbf{Overgrazing:} High livestock density (transhumance + sedentary). Trampling destroys crust/structure; selective grazing removes palatable perennials; dung removed for fuel.
    \item \textbf{Deforestation / Fuelwood:} Cutting of \emph{Acacia}, \emph{Balanites}, \emph{Adansonia} for fuelwood (urban demand: Maroua, Kousséri) and construction. Loss of windbreaks, nutrient pumps, root binding.
    \item \textbf{Inappropriate Cultivation:} Expansion of rainfed cotton/sorghum/millet onto marginal shallow soils. Clean weeding, no residues, no fallow (population pressure). Monoculture cotton $\rightarrow$ nutrient mining, pest buildup.
    \item \textbf{Salinisation/Sodification:} In irrigated perimeters (SEMRY, Maga), poor drainage + high evaporation $\rightarrow$ salt accumulation.
\end{itemize}

\textbf{2. Sustainable Management Strategy (Integrated Approach):}
\begin{itemize}
    \item \textbf{Wind Erosion Control:} Establish \textbf{multi-species windbreaks} (perimeter: \emph{Casuarina}, \emph{Azadirachta}; middle: \emph{Acacia senegal}, \emph{Balanites}; inner: fruit trees \emph{Ziziphus}, \emph{Tamarindus}). Spacing $\approx 20H$ (H = height). \textbf{Mulching} with millet/sorghum stalks (do not burn/graze all).
    \item \textbf{Water Erosion Control:} \textbf{Contour stone bunds} (pierres) / \emph{banquettes} on gentle slopes. \textbf{Half-moons (demi-lunes)} on degraded plateaus for runoff harvesting + tree planting. \textbf{Gully plugging} with gabions/vegetation in Mandara foothills.
    \item \textbf{Fertility Restoration (ISFM):} \textbf{Cereal–Legume rotation} (Millet/Sorghum – Cowpea/Groundnut). \textbf{Improved fallows} with \emph{Acacia colei}, \emph{Senna siamea} (2–3 years). \textbf{Manure management:} Composting pen manure + crop residues (reduce N loss). \textbf{Micro-dosing} fertiliser (5–10 g/hill NPK) at planting — affordable, efficient.
    \item \textbf{Agroforestry / Silvopastoralism:} \textbf{Faidherbia albida} (Apple-ring acacia) in parklands — reverse phenology (leaves in dry season = fodder/shade; bare in wet season = no crop competition). High value gum arabic (\emph{Acacia senegal}) for income.
    \item \textbf{Governance / Policy:} Secure pastoral corridors (reduce conflict/trampling). Community management of natural regeneration (FMNR — Farmer Managed Natural Regeneration) — protect spontaneous tree seedlings. MINEPDED/MINADER extension for climate-smart practices. Great Green Wall integration.
\end{itemize}

\textbf{Conclusion:} Degradation is driven by climate–human interaction. Reversal requires physical structures + biological regeneration + socio-economic incentives. Success in the Far North secures the northern breadbasket and halts desertification.
\end{exoresolu}

\begin{exercice}[Revision Questions]
\begin{enumerate}
    \item Explain why soil is considered a non-renewable resource on a human timescale. Calculate the approximate time to form a 30 cm topsoil under humid tropical conditions.
    \item Distinguish between \emph{contour ploughing} and \emph{terracing}. For each, state the slope range where it is most effective and the mechanism by which it reduces erosion.
    \item A farmer in the South-West Region applies 400 kg/ha of Urea annually to his cocoa farm but never applies lime. After 5 years, yields decline and cocoa trees show dieback. Explain the soil chemical processes causing this and prescribe a remedy.
    \item Describe the "vicious cycle" linking deforestation, soil erosion, declining yields, and poverty in the Western Highlands. How does agroforestry break this cycle?
    \item List four functions of soil organic matter (SOM) that are critical for sustainable agriculture in Cameroon. For each, give a specific management practice that enhances it.
    \item Outline the roles of MINADER, MINEPDED, and IRAD in soil protection in Cameroon.
\end{enumerate}
\end{exercice}

\begin{corrige}[Exercise 1]
\textbf{1.} Soil formation (pedogenesis) is extremely slow: weathering of rock + organic matter accumulation + horizon differentiation. Rate $\approx 0.01–0.025$ cm/yr (100–400 yr/cm). For 30 cm: $30 \times 100 = 3000$ years to $30 \times 400 = 12\,000$ years. Human civilisation operates on decades/centuries $\rightarrow$ effectively non-renewable.
\end{corrige}

\begin{corrige}[Exercise 2]
\textbf{2.}
\begin{itemize}
    \item \textbf{Contour Ploughing:} Tillage/furrows follow contour lines (constant elevation). \textbf{Slope:} $5–15\%$. \textbf{Mechanism:} Furrows act as mini-barriers $\rightarrow$ reduce runoff velocity $\rightarrow$ increase infiltration time $\rightarrow$ trap sediment in furrow bottoms.
    \item \textbf{Terracing:} Construction of level platforms (benches) separated by vertical/steep risers. \textbf{Slope:} $> 15\%$ (steep). \textbf{Mechanism:} Reduces slope length to terrace width (near zero) $\rightarrow$ eliminates runoff velocity on tread $\rightarrow$ risers (vegetated) dissipate energy. Converts steep slope to series of flat fields.
\end{itemize}
\end{corrige}

\begin{corrige}[Exercise 3]
\textbf{3.} Urea $\ce{CO(NH2)2}$ hydrolyses to $\ce{NH4+}$ $\xrightarrow{\text{nitrification}}$ $\ce{NO3-} + 2\ce{H+}$. Continuous nitrification acidifies soil. Ferralitic soils (South-West) have low CEC, low base saturation, high Al/Fe oxides. pH drops $\rightarrow$ Al$^{3+}$ solubilised (toxic to roots) $\rightarrow$ Ca/Mg leached $\rightarrow$ base saturation $\rightarrow$ 0. \textbf{Remedy:} Apply \textbf{dolomitic lime} (CaMg(CO$_3$)$_2$) at 2–3 t/ha to raise pH to 5.5–6.0, precipitate Al, supply Ca/Mg. Split Urea applications. Add organic matter (compost) to buffer pH and complex Al.
\end{corrige}

\begin{corrige}[Exercise 4]
\textbf{4.} Cycle: Population pressure $\rightarrow$ Deforestation (slash-and-burn) on steep slopes $\rightarrow$ Loss of root binding + canopy interception $\rightarrow$ High kinetic energy rainfall hits bare soil $\rightarrow$ Splash erosion $\rightarrow$ Surface seal $\rightarrow$ High runoff $\rightarrow$ Rill/Gully erosion $\rightarrow$ Loss of A horizon (OM, nutrients, seeds) $\rightarrow$ Exposed B horizon (low fertility, high Al, poor structure) $\rightarrow$ Yields plummet $\rightarrow$ Farmer cannot afford inputs $\rightarrow$ Clears more steep forest $\rightarrow$ Cycle repeats. \textbf{Agroforestry breaks it:} Trees provide permanent cover (interception), deep roots (binding, nutrient pump), prunings (mulch $\rightarrow$ OM, nutrients), products (fruit, timber, medicine $\rightarrow$ income $\rightarrow$ reduces need to clear new land).
\end{corrige}

\begin{corrige}[Exercise 5]
\textbf{5.}
\begin{enumerate}
    \item \textbf{Nutrient reservoir (N, P, S):} Mineralisation feeds crops. \emph{Practice:} Green manures (\emph{Mucuna}), compost application.
    \item \textbf{Cation Exchange Capacity (CEC):} Humus has pH-dependent CEC $\gg$ kaolinite. \emph{Practice:} Mulching, reduced tillage, biochar.
    \item \textbf{Soil Structure / Aggregation:} Polysaccharides, fungal hyphae, roots bind microaggregates $\rightarrow$ macropores $\rightarrow$ infiltration, aeration, root growth. \emph{Practice:} No-till / minimum till, cover crops, organic amendments.
    \item \textbf{Water Holding Capacity:} Humus holds $\times 20$ its weight in water. \emph{Practice:} Mulching, compost incorporation, agroforestry.
    \item \textbf{(Bonus) Biological Activity / Disease Suppression:} Diverse microbiome. \emph{Practice:} Crop rotation, reduced pesticides, organic inputs.
\end{enumerate}
\end{corrige}

\begin{corrige}[Exercise 6]
\textbf{6.}
\begin{itemize}
    \item \textbf{MINADER:} Policy implementation (fertiliser subsidy, seed distribution), agricultural extension (vulgarisation) of ISFM/conservation agriculture, soil fertility mapping (with IRAD), watershed management projects (PADFA, etc.).
    \item \textbf{MINEPDED:} Environmental policy (EIA enforcement), UNCCD focal point (National Action Programme to Combat Desertification), climate change adaptation (NAP), pollution control, coordination of Great Green Wall, environmental education.
    \item \textbf{IRAD:} Research generation: soil survey/classification, fertiliser response trials, variety screening (acid tolerance, drought), development of conservation technologies (improved fallows, Zai, composting methods), training of technicians.
\end{itemize}
\end{corrige}

\newpage

\coursec{Integration activity}

\courssub{Aims of the activity}

By the end of this integration activity, you should be able to:

\begin{itemize}
\item apply your knowledge of \cle{soil formation} and the \cle{soil-forming factors} (parent rock, climate, relief, organisms, time) to interpret a real soil profile;
\item identify, describe and label the main \cle{soil horizons} (O, A, E, B, C, R) in a soil pit or road cut;
\item use simple field and laboratory techniques (feel test, jar settling test, pH paper, colour and moisture observation) to determine the \cle{physical} and \cle{chemical characteristics} of a soil;
\item relate soil characteristics to \cle{land use}, and propose suitable crops and \cle{sustainable soil management} measures;
\item compare soils from two different Cameroonian environments and explain their differences scientifically;
\item work in a team, record data rigorously and communicate findings in a short scientific report, as expected at the GCE Advanced Level.
\end{itemize}

\courssub{Situation}

The Government Bilingual High School of Mbankomo (Centre Region) has received a plot of land behind the school to create a school farm. Before planting, the Principal asks the Upper Sixth Geology class to carry out a \cle{mini soil survey} to find out whether the soil is suitable for agriculture and how it should be protected. The school lies in the humid equatorial forest zone, with an annual rainfall of about \SI{1600}{mm}, a mean annual temperature of about \SI{25}{\ensuremath{{}^{\circ}\mathrm{C}}}, and a gently undulating relief on weathered \cle{granitic} basement rocks of the Congo craton margin.

Your group digs a small soil pit (\SI{50}{cm} $\times$ \SI{50}{cm}, about \SI{80}{cm} deep) on a gentle slope behind the classrooms. You observe and record the following:

\begin{center}
\rowcolors{2}{popblueL}{white}
\begin{tabularx}{\linewidth}{|l|X|X|}
\hline
\rowcolor{popblueL} \textbf{Layer} & \textbf{Depth (cm)} & \textbf{Field observations} \\
\hline
Layer 1 & 0--5 & Dark brown, loose, many rootlets, leaf litter on top, earthworms present \\
\hline
Layer 2 & 5--30 & Reddish brown, crumbly, moist, fine roots, feels slightly gritty \\
\hline
Layer 3 & 30--65 & Bright red, sticky when wet, hard clods when dry, few roots \\
\hline
Layer 4 & 65--80 & Yellowish red mottled, mixed with pale weathered rock fragments \\
\hline
\end{tabularx}
\end{center}

Your group also performs three simple tests:

\begin{enumerate}
\item \textbf{Feel test:} the moist soil from Layer 2 can be rolled into a short ribbon but breaks easily; it feels both gritty and slightly sticky.
\item \textbf{Jar settling test:} a sample of Layer 2 is shaken with water in a transparent jar and left to settle for 24 hours. Measured layers: sand \SI{6}{cm}, silt \SI{2.5}{cm}, clay \SI{1.5}{cm} (total settled thickness \SI{10}{cm}).
\item \textbf{pH test:} indicator paper dipped in a soil--water mixture gives pH $\approx 5.5$.
\end{enumerate}

You also notice that on the steeper part of the slope, small \cle{rills} (tiny erosion channels) are forming and the topsoil is being washed towards the stream below, especially after heavy rains.

For comparison, your teacher gives you the description of a soil profile from the \textbf{Far North Region} (near Maroua), developed in a hot semi-arid climate (rainfall about \SI{800}{mm}, long dry season) on sandy deposits: a thin greyish-brown sandy topsoil poor in organic matter, over a pale, weakly developed subsoil, with little clay and low moisture.

\textbf{Problem:} Is the school soil fertile and suitable for farming? Which crops should be planted, and how can the soil be protected from degradation?

\courssub{Tasks}

\begin{enumerate}
\item Draw a labelled sketch of the soil pit, and assign each observed layer (1--4) to its correct master horizon (O/A, B, C). Justify each assignment with one observation.
\item Using the jar settling test results, calculate the percentage of sand, silt and clay in Layer 2 and name the \cle{textural class} of the soil. Confirm your answer with the feel test.
\item State the pH of the soil and classify it as acidic, neutral or alkaline. Name one crop that tolerates this pH and one amendment that could correct it.
\item Give two pieces of evidence from the profile showing that \cle{living organisms} and \cle{climate} have influenced the formation of this soil.
\item Identify the type of \cle{soil erosion} observed on the slope, state its cause, and explain why it is dangerous for soil fertility.
\item Recommend \textbf{two suitable crops} for the school farm and \textbf{two soil-protection measures}, justifying each choice with the soil characteristics you observed.
\item Compare the school soil with the Far North profile: give two differences and explain each one using the \cle{soil-forming factors}.
\item Write a one-page group report (introduction, methods, results, conclusions, recommendations) to be presented to the Principal.
\end{enumerate}

\courssub{Model answer}

\textbf{1. Horizon identification.} Layer 1 (0--5 cm), dark, rich in rootlets, litter and earthworms, is the \cle{O/A horizon} (organic matter and humus-rich topsoil). Layer 2 (5--30 cm), reddish brown and crumbly with fine roots, is the lower \cle{A horizon} (mineral topsoil enriched in humus). Layer 3 (30--65 cm), bright red, sticky and clay-rich, is the \cle{B horizon} (subsoil of accumulation of clay and iron oxides, which give the red colour). Layer 4 (65--80 cm), mottled with weathered rock fragments, is the \cle{C horizon} (weathered parent material). Solid granite (the \cle{R horizon}) was not reached.

\begin{popfigure}
\begin{tikzpicture}[scale=0.9]
\fill[popgreenL] (0,6.5) rectangle (4,7);
\fill[brown!60!black] (0,6) rectangle (4,6.5);
\fill[brown!45] (0,4.5) rectangle (4,6);
\fill[red!55!orange] (0,2.5) rectangle (4,4.5);
\fill[orange!40!yellow] (0,1) rectangle (4,2.5);
\fill[gray!50] (0,0) rectangle (4,1);
\draw (0,0) rectangle (4,7);
\node[right, align=left, text width=5cm] at (4.2,6.75) {\textbf{O/A}: litter, humus, roots};
\node[right, align=left, text width=5cm] at (4.2,5.25) {\textbf{A}: humus-rich mineral topsoil};
\node[right, align=left, text width=5cm] at (4.2,3.5) {\textbf{B}: clay + iron oxides (red)};
\node[right, align=left, text width=5cm] at (4.2,1.75) {\textbf{C}: weathered granite fragments};
\node[right, align=left, text width=5cm] at (4.2,0.5) {\textbf{R}: bedrock (not reached)};
\draw[<->] (-0.3,7) -- (-0.3,0) node[midway, left, align=center]{depth\\ \SI{80}{cm}};
\end{tikzpicture}
\legende{Labelled sketch of the school soil pit (idealised profile of a ferrallitic soil).}
\end{popfigure}

\textbf{2. Texture.} Total settled thickness $= \SI{10}{cm}$.
\begin{align*}
\text{Sand} &= \frac{6}{10}\times 100 = 60\ \%\\
\text{Silt} &= \frac{2.5}{10}\times 100 = 25\ \%\\
\text{Clay} &= \frac{1.5}{10}\times 100 = 15\ \%
\end{align*}
With 60 \% sand, 25 \% silt and 15 \% clay, the soil is a \cle{sandy loam} (according to the USDA textural triangle). This agrees with the feel test: gritty (sand dominates) but slightly sticky and able to form a short ribbon (some clay present). A sandy loam is well drained, easy to work and warms up quickly, but it does not retain water and nutrients as well as a clay soil.

\textbf{3. pH.} pH $\approx 5.5$: the soil is \cle{slightly acidic}, typical of ferrallitic soils under humid equatorial climates where bases are leached. Crops such as cassava, groundnut, pineapple or cocoyam tolerate this pH. To correct the acidity, the soil can be \cle{limed} (addition of agricultural lime, \ce{CaCO3}), which raises the pH and supplies calcium.

\textbf{4. Evidence of soil-forming factors.}
\begin{itemize}
\item \emph{Organisms:} leaf litter, abundant rootlets and earthworms in Layer 1 show biological activity; decomposition of this litter forms the \cle{humus} that darkens the topsoil.
\item \emph{Climate:} the deep, red, clay-rich B horizon results from intense \cle{chemical weathering} (hydrolysis) and \cle{leaching} under high rainfall and temperature; iron oxides remaining after leaching give the red colour typical of \cle{ferrallitic soils}.
\end{itemize}

\textbf{5. Erosion.} The small channels on the slope are \cle{rill erosion}, a form of water erosion caused by runoff of heavy rain on bare or sparsely vegetated sloping ground. It is dangerous because it removes the \cle{topsoil} (horizons O and A), which is the most fertile layer, richest in humus and nutrients; continued erosion leads to gullies, loss of fertility and silting of the stream below.

\textbf{6. Recommendations.}
\begin{itemize}
\item \emph{Suitable crops:} \textbf{cassava} and \textbf{groundnut}. Both tolerate slightly acidic, sandy-loam, well-drained soils; groundnut, a legume, also fixes atmospheric nitrogen and improves soil fertility.
\item \emph{Protection measures:} (i) \textbf{contour planting and grass strips} on the slope to slow runoff and trap sediment, stopping the rills from growing; (ii) \textbf{mulching and addition of compost/manure} to cover the soil, keep moisture, feed the sandy loam with organic matter and protect it from raindrop impact. (Planting cover crops and avoiding bush burning are also acceptable.)
\end{itemize}

\textbf{7. Comparison with the Far North profile.}

\begin{center}
\rowcolors{2}{popblueL}{white}
\begin{tabularx}{\linewidth}{|X|X|X|}
\hline
\rowcolor{popblueL} \textbf{Difference} & \textbf{School soil (Mbankomo)} & \textbf{Far North (Maroua)} \\
\hline
Organic matter & Thick dark humus-rich topsoil & Thin topsoil, poor in humus \\
\hline
Colour and depth & Deep, red, clay-rich B horizon & Pale, thin, weakly developed, sandy \\
\hline
Moisture & Moist, leached, acidic & Dry, little leaching \\
\hline
\end{tabularx}
\end{center}

\emph{Explanation using soil-forming factors:} the differences are controlled mainly by \cle{climate} and \cle{organisms}. At Mbankomo, high rainfall and temperature cause strong chemical weathering and leaching, producing a deep red ferrallitic soil, while dense forest vegetation supplies abundant litter, giving a humus-rich topsoil. In the Far North, the semi-arid climate (low rainfall, long dry season) limits chemical weathering and leaching, so the profile stays thin and pale; sparse savanna vegetation provides little litter, hence the low organic matter content. \cle{Parent material} (granite versus sandy deposits) and \cle{time} reinforce the contrast.

\begin{attention}[Common mistakes to avoid]
Do not confuse the B horizon (accumulation of clay and iron) with the C horizon (weathered parent rock). Percentages from the jar test must be calculated relative to the \emph{total settled thickness}, not to the jar height. Finally, always justify a crop or protection measure with an \emph{observed soil characteristic} — an unjustified recommendation earns no marks.
\end{attention}

\begin{aretenir}[What this activity consolidates]
A soil profile is read from top to bottom (O, A, E, B, C, R); texture is estimated by the feel and jar tests; pH, colour, moisture and organic matter reveal the chemical and biological state; and every land-use decision must follow from the soil's characteristics and respect sustainable management: protect the topsoil, control erosion, and return organic matter to the soil.
\end{aretenir}

\newpage

\coursec{Methods and worked examples}

\courssub{Methods and Worked Examples}

This section presents the essential skills you must master for the GCE Advanced Level Geology examination on Soil. Each method is broken down into clear, numbered steps, followed by a fully worked GCE-style examination question. Study the reasoning carefully; the examiner rewards clear methodology and correct geological terminology as much as the final answer.

\begin{methode}[Identifying Soil Forming Factors and Processes (Jenny's State Factor Equation)]
\begin{enumerate}
\item \textbf{List the five state factors} using the mnemonic \textbf{CLORPT}: \textbf{C}limate, \textbf{O}rganisms (biota), \textbf{R}elief (topography), \textbf{P}arent Material, \textbf{T}ime.
\item \textbf{Analyse the specific role of each factor} in the given context:
      \begin{itemize}
      \item \textbf{Climate:} Temperature and rainfall control the \emph{rate} of chemical weathering and leaching. High rainfall + high temperature $\rightarrow$ intense leaching, laterisation. Low rainfall $\rightarrow$ accumulation of salts/carbonates.
      \item \textbf{Organisms:} Vegetation type determines organic matter input (humus), root action, and microbial activity. Forest $\rightarrow$ acid humus (mull/mor); Grassland $\rightarrow$ base-rich humus.
      \item \textbf{Relief:} Slope angle and aspect affect drainage, erosion, deposition, and microclimate (sun-facing vs. shade-facing). Steep slopes $\rightarrow$ thin, eroded soils; valley bottoms $\rightarrow$ deep, colluvial/alluvial soils.
      \item \textbf{Parent Material:} Mineralogical composition and texture of the underlying rock determine the initial nutrient reserve and texture of the soil. Basalt $\rightarrow$ clay-rich, nutrient-rich; Granite $\rightarrow$ sandy, acidic; Limestone $\rightarrow$ calcareous.
      \item \textbf{Time:} Duration of weathering. Young soils (e.g., volcanic ash, alluvium) show weak horizonation; old soils (e.g., African peneplains) show deep, strongly developed profiles (Oxisols/Ferralsols).
      \end{itemize}
\item \textbf{Identify the dominant pedogenic processes} resulting from the factor combination: \emph{Laterisation} (Fe/Al oxide accumulation), \emph{Podzolisation} (Fe/Al/organic matter translocation in acid, cool/wet climates), \emph{Gleying} (reduction in waterlogged conditions), \emph{Calcification} (CaCO$_3$ accumulation in semi-arid), \emph{Salinisation} (soluble salt accumulation), \emph{Humification/Mineralisation}.
\item \textbf{Name the resulting Soil Order} (using FAO/WRB or USDA Soil Taxonomy names common in Cameroon syllabus: Ferralsols, Acrisols, Nitisols, Andosols, Vertisols, Fluvisols, Cambisols, Arenosols, Regosols, Leptosols, Gleysols, Plinthosols).
\end{enumerate}
\end{methode}

\begin{exoresolu}[Soil Formation on the Cameroon Volcanic Line]
\textbf{Statement:} Mount Cameroon is an active stratovolcano composed mainly of basaltic lavas and pyroclastic deposits. The western flank receives over \SI{4000}{mm} of rainfall annually, supporting dense tropical rainforest. The upper slopes (> \SI{2000}{m}) have a cooler, misty climate with grassland vegetation. Recent lava flows (1999, 2000) are colonised by pioneer species. 
\begin{enumerate}
\item Using the CLORPT framework, explain why the soils on the lower western slopes are deep, red, clay-rich Ferralsols, while soils on the recent lava flows are shallow, stony Leptosols/Andosols.
\item Identify the dominant pedogenic process operating on the lower slopes and the process initiating on the recent flows.
\end{enumerate}

\tcblower
\textbf{Detailed Solution:}

\noindent\textbf{1. CLORPT Analysis:}

\begin{center}
\begin{tabularx}{\linewidth}{|l|X|X|}\hline
\rowcolor{popblueL}
\textbf{Factor} & \textbf{Lower Western Slopes (Ferralsols)} & \textbf{Recent Lava Flows (Leptosols/Andosols)} \\ \hline
\textbf{C}limate & High rainfall (\textgreater 4000 mm), warm temperatures year-round. Intense chemical weathering and leaching. & Same macro-climate, but high altitude on upper slopes brings cooler temperatures, reducing weathering rates locally. \\ \hline
\textbf{O}rganisms & Dense tropical rainforest: high biomass, rapid nutrient cycling, acidic humus (mor/moder) promoting acid leaching. & Pioneer vegetation (lichens, mosses, ferns): very low organic matter input, slow nutrient cycling. \\ \hline
\textbf{R}elief & Moderate to gentle slopes on lower flanks: good drainage, low erosion, allows deep profile development. & Irregular, rough topography of aa/pahoehoe flows; steep flow fronts. High runoff, unstable surfaces. \\ \hline
\textbf{P}arent Material & Deeply weathered basalt (saprolite) derived from older volcanic cycles. Basalt is rich in Fe, Mg, Ca $\rightarrow$ high nutrient reserve initially. & Fresh, unweathered basaltic lava (rock). Glassy, porous texture. High primary mineral content (olivine, pyroxene, plagioclase). \\ \hline
\textbf{T}ime & Very long (hundreds of thousands of years on stable peneplain remnants). Multiple weathering cycles. & Extremely short (years to decades since 1999/2000 flows). Negligible weathering time. \\ \hline
\end{tabularx}
\end{center}

\noindent\textbf{2. Dominant Pedogenic Processes:}
\begin{itemize}
\item \textbf{Lower slopes: Laterisation (Ferralitisation).} Intense hydrolysis of primary silicates (olivine, pyroxene, feldspar) under hot, wet conditions releases Fe$^{3+}$ and Al$^{3+}$ which precipitate as insoluble oxides/hydroxides (haematite, goethite, gibbsite). Bases (Ca, Mg, K, Na) and silica are leached downwards and lost. Result: deep, red, kaolinitic/oxidic, low CEC Ferralsols (Oxisols).
\item \textbf{Recent flows: Initial Weathering / Andosolisation.} Physical breakdown (thermal stress, root wedging) and incipient chemical weathering of volcanic glass release Al-humus complexes and short-range-order minerals (allophane, imogolite). High phosphate fixation. Organic matter accumulation begins (Andic properties). If vegetation establishes well $\rightarrow$ Andosols; if erosion dominates $\rightarrow$ Leptosols.
\end{itemize}

\noindent\textbf{Key Exam Point:} Always link the \emph{process} to the specific \emph{factor combination}. "High rainfall causes leaching" is insufficient; "High rainfall \emph{combined with} high temperature and \emph{permeable} parent material causes intense laterisation" is a full answer.
\end{exoresolu}

\begin{methode}[Describing and Interpreting a Soil Profile (Horizon Notation)]
\begin{enumerate}
\item \textbf{Draw the profile to scale} (depth in cm). Use standard horizon symbols (Master horizons: O, A, E, B, C, R; Transitional: AB, BA, BC, CB; Subordinate distinctions: lower case suffixes).
\item \textbf{For each horizon, describe systematically:}
      \begin{itemize}
      \item \textbf{Depth/Thickness:} Upper and lower boundaries (clear, gradual, diffuse; smooth, wavy, irregular, broken).
      \item \textbf{Colour:} Munsell notation (Hue Value/Chroma) moist and dry (e.g., 10YR 3/2 moist). Note mottles (redoximorphic features).
      \item \textbf{Texture:} Field estimate (sand, silt, clay percentages $\rightarrow$ textural class: sand, loam, clay loam, clay).
      \item \textbf{Structure:} Type (granular, blocky, prismatic, platy, massive, single grain), Grade (weak, moderate, strong), Size (fine, medium, coarse).
      \item \textbf{Consistence:} Dry (loose, soft, hard), Moist (loose, friable, firm), Wet (non-sticky/non-plastic $\rightarrow$ very sticky/very plastic).
      \item \textbf{Roots/Pores:} Abundance (few, common, many) and size (fine, medium, coarse).
      \item \textbf{Special Features:} Concretions/nodules (Fe/Mn, CaCO$_3$, Gibbsite), cutans (clay skins), slickensides, charcoal, artefacts.
      \end{itemize}
\item \textbf{Assign Horizon Designations:} Match observations to diagnostic criteria (e.g., \textbf{mollic} = thick, dark, base-rich A; \textbf{umbric} = thick, dark, base-poor A; \textbf{oxic} = B horizon with $\ge$ 15\% clay, low CEC, $\le$ 10\% weatherable minerals; \textbf{argic} = illuvial clay accumulation (clay skins); \textbf{spodic} = illuvial Fe/Al/organic matter; \textbf{plinthic} = Fe-rich, humus-poor mixture that hardens irreversibly).
\item \textbf{Classify the Soil} (WRB Reference Soil Group + qualifiers, or USDA Order/Suborder).
\end{enumerate}
\end{methode}

\begin{exoresolu}[Profile Description and Classification: A Soil from the Bamenda Highlands]
\textbf{Statement:} A soil pit is dug on a gentle mid-slope (\SI{5}{\%}) under grassland (Sudano-Guinea savanna) at \SI{1500}{m} altitude on the Bamenda Highlands (trachytic/basaltic parent material). Annual rainfall \SI{2000}{mm}, distinct dry season. The profile description is:
\begin{itemize}
\item \textbf{0--25 cm:} Dark reddish brown (5YR 3/3 moist), clay loam; strong, fine, subangular blocky; friable; many fine/medium roots; clear smooth boundary.
\item \textbf{25--60 cm:} Dark red (2.5YR 3/6 moist), clay; moderate, medium, angular blocky; firm; common fine roots; few thin clay cutans on ped faces; gradual smooth boundary.
\item \textbf{60--120 cm:} Red (2.5YR 4/8 moist), clay; weak, coarse, angular blocky; firm; very few roots; common distinct Fe/Mn nodules (5--10 mm); diffuse boundary.
\item \textbf{120--180 cm+:} Red (2.5YR 5/8 moist), clay; massive; very firm; abundant Fe/Mn concretions hardening to plinthite on exposure; weathered trachyte fragments.
\end{itemize}
Laboratory data: Topsoil pH(H$_2$O) 5.8, Base Saturation 65\%, CEC 18 cmol$_c$/kg clay. Subsoil pH 5.2, Base Saturation 35\%, CEC 12 cmol$_c$/kg clay. Clay content increases from 30\% to 55\% with depth.
\begin{enumerate}
\item Assign horizon designations (Master + subordinate) for each layer.
\item Identify the diagnostic horizons and properties present.
\item Classify the soil to the Reference Soil Group (WRB/FAO) level. Justify.
\end{enumerate}

\tcblower
\textbf{Detailed Solution:}

\noindent\textbf{1. Horizon Designations:}
\begin{itemize}
\item \textbf{0--25 cm: Ah} (Mineral surface horizon, dark colour from organic matter, base saturation > 50\% $\rightarrow$ \textbf{mollic} qualifier possible; thickness 25 cm meets WRB minimum for mollic; structure $\rightarrow$ granular/blocky).
\item \textbf{25--60 cm: Bt} (\textbf{B} horizon with evidence of \textbf{t} = illuvial clay accumulation: clay increase 30\% $\rightarrow$ 55\%, presence of \emph{clay cutans} (argillans) on ped faces).
\item \textbf{60--120 cm: Bt} (Continuation of argic horizon, weaker structure, Fe/Mn nodules = \emph{concretionary} feature).
\item \textbf{120--180 cm+: BC} (Transitional to parent material; massive structure; abundant plinthite = \emph{plinthic} feature; weathered rock fragments).
\end{itemize}
\textit{Note: No E horizon (no eluviation of clay/Fe/Al leaving a pale layer). No O horizon (grassland, rapid mineralisation).}

\noindent\textbf{2. Diagnostic Horizons/Properties:}
\begin{itemize}
\item \textbf{Mollic horizon (surface):} Ah is $\ge$ 25 cm thick, dark, base-rich (BS 65\%), structured. Qualifies as \textbf{Mollic}.
\item \textbf{Argic horizon (subsurface):} Bt horizon (25--120 cm) shows distinct clay increase (factor $\ge$ 1.2 or absolute $\ge$ 3\% over 30 cm vertical) and clay cutans. Starts at 25 cm (within 100 cm).
\item \textbf{Ferralic properties:} Subsoil CEC (12 cmol$_c$/kg clay) is low (low activity clay). Weatherable minerals likely low in trachyte/basalt saprolite. However, presence of \textbf{argic horizon} overrides ferralic for classification at RSG level.
\item \textbf{Plinthic material:} Abundant Fe/Mn concretions hardening irreversibly (plinthite) in BC horizon (starting at 120 cm).
\item \textbf{Base saturation:} Decreases with depth (65\% $\rightarrow$ 35\%), typical of leaching.
\end{itemize}

\noindent\textbf{3. Classification (WRB 2015/2022):}
\begin{itemize}
\item \textbf{Reference Soil Group (RSG):} \textbf{Lixisol}. 
      \begin{itemize}
      \item Reason: Presence of an \textbf{argic horizon} (Bt) starting within 100 cm.
      \item Low activity clay (CEC $\le$ 24 cmol$_c$/kg clay $\rightarrow$ 12 cmol$_c$/kg clay).
      \item Base saturation $\ge$ 50\% in the major part of the argic horizon between 20--100 cm (Upper Bt 25--60 cm likely BS > 50\%, lower part drops). Fits Lixisol better than Luvisol (high activity clay) or Alisol (low BS).
      \item Surface horizon is \textbf{mollic}.
      \item No nitic horizon (requires shiny ped faces, gradual boundaries).
      \item No plinthic horizon starting $\le$ 50 cm (plinthite deep at 120 cm).
      \end{itemize}
\item \textbf{Principal Qualifiers:} \textbf{Chromic} (hue redder than 5YR in Bt, i.e., 2.5YR), \textbf{Eutric} (base saturation $\ge$ 50\% in upper part of argic), \textbf{Clayic} (clay $\ge$ 35\% in argic), \textbf{Endoplinthic} (plinthite $\ge$ 15\% vol within 150 cm).
\item \textbf{Full Name:} \textbf{Chromic Eutric Clayic Endoplinthic Lixisol}. (Simplified for exam: \textbf{Chromic Lixisol}).
\end{itemize}

\noindent\textbf{Exam Tip:} Always check for \emph{Argic} vs \emph{Nitic} vs \emph{Ferralic}. Nitic requires shiny ped faces (pressure faces) and gradual horizon boundaries. Ferralic requires very low CEC and no clay cutans. Argic requires clay cutans or significant clay increase. This profile has clear cutans $\rightarrow$ Argic. Low activity clay + High BS $\rightarrow$ \textbf{Lixisol} (WRB) $\approx$ \textbf{Ustalf/Haplustalf} (USDA) or \textbf{Ultisol} with high BS.
\end{exoresolu}

\begin{methode}[Determining Soil Texture Using the Textural Triangle (USDA/FAO)]
\begin{enumerate}
\item \textbf{Obtain particle size distribution data} (\% Sand \SIrange{2.0}{0.05}{mm}, \% Silt \SIrange{0.05}{0.002}{mm}, \% Clay < \SI{0.002}{mm}) from laboratory analysis (pipette/hydrometer method) or field estimation (hand-feel method).
\item \textbf{Verify the sum} = 100\% (or very close). If not, recalculate percentages on the < \SI{2}{mm} fraction basis (exclude gravel > \SI{2}{mm}).
\item \textbf{Locate the point on the Textural Triangle:}
      \begin{itemize}
      \item Read \% Clay on the horizontal axis (left to right).
      \item Read \% Silt on the right axis (bottom left to top right).
      \item Read \% Sand on the left axis (top left to bottom right).
      \item The intersection of the three lines gives the texture class.
      \end{itemize}
\item \textbf{Name the texture class} (e.g., Sandy Loam, Silty Clay Loam, Clay). Use the standard 12 classes (USDA) or the FAO triangle (slightly different boundaries).
\item \textbf{Interpret the texture} for management: Water holding capacity (Clay > Loam > Sand), Drainage/Infiltration (Sand > Loam > Clay), Nutrient retention (CEC: Clay > Loam > Sand), Workability (Loam best), Erodibility (Silt > Fine Sand > Clay > Coarse Sand).
\end{enumerate}
\end{methode}

\begin{exoresolu}[Texture Determination and Management Implications]
\textbf{Statement:} A soil survey in the Upper Nun Valley (Bamenda Highlands) yields the following particle size distribution for a topsoil sample (fine earth fraction < \SI{2}{mm}):
\begin{center}
\begin{tabular}{|c|c|c|c|}\hline
\rowcolor{popblueL}
\textbf{Fraction} & \textbf{Size Range (mm)} & \textbf{Mass (g)} & \textbf{Percentage (\%)}\\\hline
Coarse Sand & 2.0 -- 0.6 & 12.5 & \\
Fine Sand & 0.6 -- 0.2 & 18.3 & \\
Very Fine Sand & 0.2 -- 0.05 & 9.2 & \\
Coarse Silt & 0.05 -- 0.02 & 15.0 & \\
Fine Silt & 0.02 -- 0.002 & 25.0 & \\
Clay & < 0.002 & 20.0 & \\\hline
Total & & 100.0 & \\\hline
\end{tabular}
\end{center}
\begin{enumerate}
\item Calculate the percentages for the three main USDA fractions: Sand, Silt, Clay.
\item Plot the point on the USDA Textural Triangle (describe the location) and determine the Texture Class.
\item Discuss the agricultural potential and management constraints of this texture for maize cultivation.
\end{enumerate}

\tcblower
\textbf{Detailed Solution:}

\noindent\textbf{1. Calculation of Main Fractions:}
\begin{align*}
\textbf{Sand} &= \text{Coarse Sand} + \text{Fine Sand} + \text{Very Fine Sand} \\
&= 12.5 + 18.3 + 9.2 = \mathbf{40.0\%} \\[4pt]
\textbf{Silt} &= \text{Coarse Silt} + \text{Fine Silt} \\
&= 15.0 + 25.0 = \mathbf{40.0\%} \\[4pt]
\textbf{Clay} &= \mathbf{20.0\%} \\
\text{Sum} &= 40.0 + 40.0 + 20.0 = 100.0\% \quad \checkmark
\end{align*}

\noindent\textbf{2. Texture Class Determination:}
\begin{itemize}
\item On the USDA triangle: Locate 20\% Clay (horizontal axis). Draw line parallel to Sand axis (slanting up-right).
\item Locate 40\% Silt (right axis). Draw line parallel to Clay axis (horizontal left).
\item Locate 40\% Sand (left axis). Draw line parallel to Silt axis (slanting down-right).
\item The three lines intersect in the region labelled \textbf{Loam} (specifically near the centre, bordering \textbf{Silt Loam} and \textbf{Sandy Clay Loam}).
\item \textbf{Texture Class: Loam.} (Equal sand/silt, moderate clay).
\end{itemize}

\noindent\textbf{3. Agricultural Potential and Management for Maize:}
\begin{center}
\begin{tabularx}{\linewidth}{|l|X|}\hline
\rowcolor{popblueL}
\textbf{Property} & \textbf{Implication for Maize} \\ \hline
\textbf{Water Holding Capacity} & \textbf{Good.} Loam holds plant-available water well (approx. \SI{150--200}{mm/m} depth). Reduces drought stress during dry spells in the growing season. \\ \hline
\textbf{Drainage / Aeration} & \textbf{Good.} Sufficient macropores (from sand/structure) for root respiration. Maize is sensitive to waterlogging. \\ \hline
\textbf{Nutrient Retention (CEC)} & \textbf{Moderate.} 20\% clay + organic matter provides reasonable CEC. Fertiliser (NPK) applications will be retained but split applications advised to minimise leaching of N (nitrate). \\ \hline
\textbf{Workability / Tilth} & \textbf{Excellent.} "Loam" is the ideal texture for cultivation. Easy to till at correct moisture content. Good seedbed preparation for maize emergence. \\ \hline
\textbf{Erodibility} & \textbf{Moderate to High.} High silt content (40\%) makes it susceptible to splash erosion and crusting if bare. \textbf{Management:} Mulching, cover crops, contour ridging, minimum tillage essential on slopes. \\ \hline
\textbf{Compaction Risk} & \textbf{Moderate.} Avoid heavy machinery on wet soil. Controlled traffic recommended. \\ \hline
\end{tabularx}
\end{center}

\noindent\textbf{Conclusion:} High potential for maize \emph{if} erosion is controlled on slopes and soil organic matter is maintained to stabilise structure.
\end{exoresolu}

\begin{methode}[Calculating Soil Physical Properties: Bulk Density, Porosity, Water Content]
\begin{enumerate}
\item \textbf{Identify the known variables:} Mass of oven-dry soil ($M_s$), Volume of soil core/sample ($V_t$), Mass of wet soil ($M_{wet}$), Particle density ($\rho_s$, typically \SI{2.65}{g.cm^{-3}} for mineral soils, lower for organic/volcanic).
\item \textbf{Apply the fundamental definitions:}
      \begin{align*}
      \text{Bulk Density } (\rho_b) &= \frac{M_s}{V_t} \quad (\SI{}{g.cm^{-3}} \text{ or } \SI{}{Mg.m^{-3}}) \\
      \text{Porosity } (f) &= 1 - \frac{\rho_b}{\rho_s} \quad (\text{decimal or \%}) \\
      \text{Gravimetric Water Content } (\theta_g) &= \frac{M_{wet} - M_s}{M_s} \quad (\text{g/g or \%}) \\
      \text{Volumetric Water Content } (\theta_v) &= \theta_g \times \rho_b \quad (\text{cm}^3/\text{cm}^3 \text{ or \%}) \\
      \text{Air-filled Porosity } (f_a) &= f - \theta_v
      \end{align*}
\item \textbf{Interpret the values:}
      \begin{itemize}
      \item $\rho_b$: $<$ \SI{1.1}{g.cm^{-3}} (good, loose/organic), \SI{1.1}{g.cm^{-3}}--\SI{1.4}{g.cm^{-3}} (ideal for root growth), $>$ \SI{1.6}{g.cm^{-3}} (compacted, restricts roots), $>$ \SI{1.8}{g.cm^{-3}} (very compact/dense).
      \item $f$: $>$ 50\% (good), 40--50\% (moderate), $<$ 40\% (poor aeration/drainage).
      \item $\theta_v$ at Field Capacity (FC) $\approx$ 0.2--0.4; at Permanent Wilting Point (PWP) $\approx$ 0.05--0.2. Available Water Capacity (AWC) = FC - PWP.
      \item $f_a$: $>$ 10--15\% needed for adequate gas exchange (O$_2$ diffusion).
      \end{itemize}
\item \textbf{Field Capacity / Wilting Point estimation:} Use texture-based pedotransfer functions if lab data unavailable (e.g., Saxton \& Rawls equations).
\end{enumerate}
\end{methode}

\begin{exoresolu}[Soil Physical Quality Assessment for Irrigation Scheduling]
\textbf{Statement:} A farmer in the Logone Valley (Far North Region) wants to schedule irrigation for onions on a clay soil. An undisturbed core sample (\SI{100}{cm^3} volume) is taken at \SI{30}{cm} depth. 
\begin{itemize}
\item Mass of wet core + soil = \SI{285}{g}
\item Mass of oven-dry core + soil = \SI{260}{g}
\item Mass of empty core = \SI{85}{g}
\item Particle density ($\rho_s$) = \SI{2.65}{g.cm^{-3}}
\end{itemize}
Laboratory analysis gives: Field Capacity (FC) = 38\% (vol), Permanent Wilting Point (PWP) = 22\% (vol).
\begin{enumerate}
\item Calculate Bulk Density ($\rho_b$), Total Porosity ($f$), Gravimetric ($\theta_g$) and Volumetric ($\theta_v$) water content of the sample.
\item Determine if irrigation is needed immediately (compare $\theta_v$ to FC and PWP).
\item Calculate the Available Water Capacity (AWC) in \SI{}{mm} for a \SI{60}{cm} root zone. How many mm of water can be depleted before irrigation is required (Management Allowed Depletion, MAD = 50\% of AWC)?
\end{enumerate}

\tcblower
\textbf{Detailed Solution:}

\noindent\textbf{1. Mass and Volume Calculations:}
\begin{align*}
M_s \text{ (dry soil)} &= 260 - 85 = \SI{175}{g} \\
M_{wet} \text{ (wet soil)} &= 285 - 85 = \SI{200}{g} \\
M_w \text{ (water)} &= M_{wet} - M_s = 200 - 175 = \SI{25}{g} \\
V_t &= \SI{100}{cm^3}
\end{align*}

\noindent\textbf{Bulk Density:}
\[
\rho_b = \frac{M_s}{V_t} = \frac{175}{100} = \mathbf{1.75\ g.cm^{-3}}
\]
\textit{Interpretation: High. Typical for compacted clay or vertic soil. Root restriction likely.}

\noindent\textbf{Total Porosity:}
\[
f = 1 - \frac{\rho_b}{\rho_s} = 1 - \frac{1.75}{2.65} = 1 - 0.660 = \mathbf{0.340 \text{ or } 34.0\%}
\]
\textit{Interpretation: Low porosity. Limited air/water storage.}

\noindent\textbf{Gravimetric Water Content:}
\[
\theta_g = \frac{M_w}{M_s} = \frac{25}{175} = \mathbf{0.143\ g/g \text{ (or } 14.3\%)}
\]

\noindent\textbf{Volumetric Water Content:}
\[
\theta_v = \theta_g \times \rho_b = 0.143 \times 1.75 = \mathbf{0.250\ cm^3/cm^3 \text{ (or } 25.0\%)}
\]
(Alternatively: $\theta_v = \frac{V_w}{V_t} = \frac{25\ cm^3}{100\ cm^3} = 0.25$ since \SI{1}{g} water $\approx$ \SI{1}{cm^3}).

\noindent\textbf{2. Irrigation Decision:}
\begin{itemize}
\item Current $\theta_v$ = 25.0\%.
\item PWP = 22\%. FC = 38\%.
\item Current status: \textbf{Above PWP (25 > 22), but well below FC (25 < 38).}
\item Soil is in the "Available Water" range but only at $\frac{25-22}{38-22} = \frac{3}{16} = 18.75\%$ of available capacity.
\item \textbf{Decision: Irrigation is NEEDED soon.} The crop is experiencing mild water stress. For onions (shallow roots, sensitive), maintain $\theta_v$ closer to FC.
\end{itemize}

\noindent\textbf{3. Available Water Capacity (AWC) and MAD:}
\begin{align*}
\text{AWC (volumetric)} &= FC - PWP = 38\% - 22\% = 16\% = 0.16\ cm^3/cm^3 \\
\text{AWC (depth)} &= 0.16 \times \SI{60}{cm} = \mathbf{9.6\ cm = 96\ mm} \\
\text{MAD (50\%)} &= 0.50 \times 96\ mm = \mathbf{48\ mm}
\end{align*}
\textit{Current depletion} = $(FC - \theta_v) \times \text{depth} = (0.38 - 0.25) \times 60 = 0.13 \times 60 = \SI{7.8}{cm} = \SI{78}{mm}$.
\textit{Since current depletion (78 mm) > MAD (48 mm), irrigation is OVERDUE.}

\noindent\textbf{Exam Tip:} Always convert to consistent units (cm or mm). 1\% vol $\times$ 1 cm depth = 0.1 mm. 16\% $\times$ 60 cm = 9.6 cm = 96 mm. High bulk density here is a major constraint; deep ripping or organic matter addition needed long-term.
\end{exoresolu}

\begin{methode}[Interpreting Soil Chemical Data: pH, CEC, Base Saturation, Nutrient Ratios]
\begin{enumerate}
\item \textbf{pH (H$_2$O and KCl):} 
      \begin{itemize}
      \item pH $<$ 4.5: Very strongly acid (Al toxicity risk, Mo deficiency).
      \item pH 4.5--5.5: Strongly acid (Al toxicity possible, P fixation high).
      \item pH 5.5--6.5: Slightly acid to neutral (Optimal for most crops, max microbial activity).
      \item pH 6.5--7.5: Neutral to slightly alkaline.
      \item pH $>$ 7.5: Alkaline (Fe, Mn, Zn, Cu, P deficiency risk; Na toxicity if sodic).
      \item $\Delta$pH = pH(KCl) - pH(H$_2$O). Negative $\Delta$pH $\rightarrow$ net negative charge (variable charge soils, oxides/kaolinite). Positive $\Delta$pH $\rightarrow$ net positive charge (rare, high anion exchange).
      \end{itemize}
\item \textbf{Cation Exchange Capacity (CEC):} 
      \begin{itemize}
      \item Units: cmol$_c$/kg (meq/100g). 
      \item Low $<$ 10 (sandy, kaolinitic/oxidic), Medium 10--25 (illitic/mixed), High $>$ 25 (smectitic, vermiculite, high OM).
      \item \textbf{Effective CEC (ECEC)} = Sum of exchangeable bases (Ca, Mg, K, Na) + Exchangeable Acidity (Al + H). Better for acid soils.
      \end{itemize}
\item \textbf{Base Saturation (BS\%):} 
      \[
      BS\% = \frac{\text{Sum of Exchangeable Bases (Ca+Mg+K+Na)}}{\text{CEC (or ECEC)}} \times 100
      \]
      \begin{itemize}
      \item $<$ 35\%: Low (Acid, leached, Al saturation high).
      \item 35--60\%: Moderate.
      \item 60--80\%: High (Good fertility).
      \item $>$ 80\%: Very high (Approaching saturation, risk of nutrient imbalance).
      \end{itemize}
\item \textbf{Exchangeable Cation Ratios (Ideal ranges for crop growth):}
      \begin{itemize}
      \item Ca:Mg $\approx$ 3:1 to 5:1 (up to 10:1 for peanuts/calciophiles). $<$ 2:1 $\rightarrow$ Mg-induced Ca deficiency.
      \item (Ca+Mg):K $\approx$ 10:1 to 20:1. High K $\rightarrow$ Mg/Ca uptake inhibition.
      \item Mg:K $\approx$ 3:1 to 5:1.
      \item ESP (Exchangeable Sodium Percentage) = $\frac{Na}{CEC} \times 100$. ESP $>$ 15\% $\rightarrow$ Sodic soil (dispersion, crusting).
      \end{itemize}
\item \textbf{Phosphorus (Available P - Bray/Olsen):} Low $<$ 10 mg/kg, Medium 10--20, High $>$ 20 (Bray-1). Acid soils $\rightarrow$ Bray; Alkaline/Calcareous $\rightarrow$ Olsen.
\end{enumerate}
\end{methode}

\begin{exoresolu}[Chemical Fertility Diagnosis for a Cocoa Farm in the Centre Region]
\textbf{Statement:} Soil analysis from a mature cocoa plantation (15 years old) on a Ferralsol (Oxisol) near Yaoundé gives the following results for the 0--20 cm layer:
\begin{center}
\begin{tabular}{|l|c|}\hline
\rowcolor{popblueL}
\textbf{Parameter} & \textbf{Value} \\\hline
pH (H$_2$O) & 4.8 \\
pH (KCl) & 3.9 \\
CEC (NH$_4$OAc, pH 7) & 8.5 cmol$_c$/kg \\
Exchangeable Ca & 1.2 cmol$_c$/kg \\
Exchangeable Mg & 0.4 cmol$_c$/kg \\
Exchangeable K & 0.15 cmol$_c$/kg \\
Exchangeable Na & 0.05 cmol$_c$/kg \\
Exchangeable Al & 1.8 cmol$_c$/kg \\
Exchangeable H & 0.5 cmol$_c$/kg \\
Available P (Bray-1) & 4 mg/kg \\
Organic C & 1.8\% \\\hline
\end{tabular}
\end{center}
\begin{enumerate}
\item Calculate the Base Saturation (BS\%) using CEC (pH 7) and the Effective CEC (ECEC). Calculate Al Saturation.
\item Interpret the pH, $\Delta$pH, CEC, and nutrient status.
\item Recommend a liming and fertilisation strategy for sustained cocoa yield.
\end{enumerate}

\tcblower
\textbf{Detailed Solution:}

\noindent\textbf{1. Calculations:}
\begin{align*}
\text{Sum of Bases (SBS)} &= Ca + Mg + K + Na = 1.2 + 0.4 + 0.15 + 0.05 = \mathbf{1.80\ cmol_c/kg} \\
\text{Exchangeable Acidity (EA)} &= Al + H = 1.8 + 0.5 = \mathbf{2.3\ cmol_c/kg} \\
\text{ECEC} &= SBS + EA = 1.80 + 2.30 = \mathbf{4.10\ cmol_c/kg} \\
\text{BS\% (on CEC pH 7)} &= \frac{1.80}{8.5} \times 100 = \mathbf{21.2\%} \\
\text{BS\% (on ECEC)} &= \frac{1.80}{4.10} \times 100 = \mathbf{43.9\%} \\
\text{Al Saturation} &= \frac{Exch.\ Al}{ECEC} \times 100 = \frac{1.8}{4.1} \times 100 = \mathbf{43.9\%}
\end{align*}

\noindent\textbf{2. Interpretation:}
\begin{itemize}
\item \textbf{pH 4.8 (H$_2$O):} Strongly acidic. Suboptimal for cocoa (prefers 5.5--6.5).
\item \textbf{$\Delta$pH = 3.9 - 4.8 = -0.9:} Strongly negative. Confirms \textbf{variable charge} mineralogy (kaolinite, Fe/Al oxides typical of Ferralsols). CEC at pH 7 (8.5) overestimates field CEC; ECEC (4.1) is the real operational CEC.
\item \textbf{CEC (8.5) / ECEC (4.1):} Low. Typical of highly weathered Ferralsols. Low nutrient retention capacity. Nutrients must be supplied via organic matter cycling and frequent, small fertiliser applications.
\item \textbf{Base Saturation (21\% on CEC7, 44\% on ECEC):} Low. High leaching environment.
\item \textbf{Al Saturation (44\%):} \textbf{CRITICAL.} $>$ 20--30\% is toxic for most crops, including cocoa. Al$^{3+}$ damages root tips, inhibits Ca/Mg uptake. This is the primary yield-limiting factor.
\item \textbf{Ca:Mg Ratio:} $1.2 / 0.4 = 3.0$. Acceptable (ideal 3--5).
\item \textbf{K Status:} 0.15 cmol$_c$/kg is low. Cocoa requires high K for pod filling.
\item \textbf{P (4 mg/kg):} Very low (Bray-1). High P fixation by Fe/Al oxides.
\item \textbf{Org C (1.8\%):} Low for a 15-yr cocoa system under shade. Indicates rapid mineralisation/low residue return.
\end{itemize}

\noindent\textbf{3. Management Recommendations:}
\begin{enumerate}
\item \textbf{LIMING (URGENT):} Apply \textbf{Dolomitic Limestone (CaMg(CO$_3$)$_2$)} to raise pH to 5.5--6.0 and supply Ca + Mg.
      \begin{itemize}
      \item Lime Requirement (LR) estimate: Target $\Delta$pH = 1.0 unit. Buffer capacity of Ferralsol $\approx$ 0.5--1.0 t/ha per 0.1 pH unit? Better: LR (t/ha) $\approx$ (Target Al Sat $<$ 10\%) $\rightarrow$ Need to neutralise excess Al. 
      \item Practical rate: \textbf{2--3 t/ha} Dolomite broadcast and incorporated (if possible) or surface applied under mulch. Split application (1.5 t/ha $\times$ 2 years) safer.
      \item Effect: Raises pH, precipitates Al(OH)$_3$, increases BS\%, improves P availability, enhances microbial activity/nodulation (shade trees).
      \end{itemize}
\item \textbf{PHOSPHORUS:} \textbf{Rock Phosphate (e.g., Minjingu PR)} is better than TSP on acid Ferralsols. Acid soil dissolves PR slowly. Apply \textbf{50--100 kg P$_2$O$_5$/ha} as PR every 2--3 years. TSP would be fixed instantly.
\item \textbf{POTASSIUM:} Apply \textbf{Muriate of Potash (KCl) or Sulfate of Potash (SOP)}. Rate: \textbf{100--150 kg K$_2$O/ha/yr} split (pre-flowering, pod filling). Avoid KCl if Cl$^-$ sensitivity suspected (rare for cocoa), SOP preferred.
\item \textbf{NITROGEN:} \textbf{Urea} or \textbf{Ammonium Sulfate} (acidifying, use with lime). \textbf{50--80 kg N/ha/yr} split. Shade legumes (Gliricidia, Leucaena) can fix N.
\item \textbf{ORGANIC MATTER:} Return cocoa pod husks (composted), prune residues. Mulching conserves moisture, adds K, Ca, reduces soil temperature.
\end{enumerate}

\noindent\textbf{Key Exam Concept:} On variable charge soils (Ferralsols, Acrisols), \textbf{ECEC and Al Saturation} are the critical diagnostic values, not CEC(pH 7). Liming targets Al toxicity, not just pH.
\end{exoresolu}

\begin{methode}[Assessing Land Capability / Suitability for Agriculture (FAO Framework)]
\begin{enumerate}
\item \textbf{Determine the Land Use Requirements (LURs)} for the specific crop/system (e.g., Maize: deep, well-drained, pH 5.5--7.0, moderate fertility, slope $<$ 15\%; Rice (irrigated): flat, impermeable subsoil, high water holding; Cocoa: deep, well-drained, pH 5.0--6.5, high organic matter, shade).
\item \textbf{Match Land Qualities (LQs) to LURs:} Rate each relevant land quality (e.g., moisture availability, oxygen availability, nutrient availability, rooting conditions, erosion hazard, workability) on a suitability scale: \textbf{S1} (Highly Suitable), \textbf{S2} (Moderately Suitable), \textbf{S3} (Marginally Suitable), \textbf{N1} (Currently Not Suitable), \textbf{N2} (Permanently Not Suitable).
\item \textbf{Apply the "Limiting Factor Principle" (Liebig's Law of the Minimum):} The overall suitability class for a Land Utilization Type (LUT) is determined by the \textbf{most limiting} land quality. No averaging.
\item \textbf{Identify the specific limitations} (subclass symbols): 
      \begin{itemize}
      \item \textbf{m} - moisture deficiency, \textbf{w} - excess water/oxygen deficiency, \textbf{n} - nutrient deficiency, \textbf{r} - rooting restriction (depth, gravel), \textbf{t} - topography/slope/erosion, \textbf{s} - salinity/sodicity, \textbf{c} - climate (temp), \textbf{f} - flooding.
      \end{itemize}
\item \textbf{Propose Improvements (Input Levels):} Can the limitation be corrected? 
      \begin{itemize}
      \item \textbf{Level A (Low inputs):} No major amendments (traditional).
      \item \textbf{Level B (Intermediate):} Fertilisers, lime, simple drainage, contour farming.
      \item \textbf{Level C (High inputs):} Irrigation, deep drainage, terracing, heavy fertilisation, soil amendment.
      \end{itemize}
      Re-classify assuming improvements (Potential Suitability).
\end{enumerate}
\end{methode}

\begin{exoresolu}[Land Suitability Classification for Rice Expansion in the Ndop Plain]
\textbf{Statement:} The Ndop Plain (Northwest Region) is a floodplain with Vertisols (smectitic clay) and Fluvisols (alluvial). A proposal exists to expand irrigated rice cultivation. Site characteristics for a representative Vertisol profile:
\begin{itemize}
\item Topography: Flat ($<$ 1\% slope), extensive.
\item Climate: \SI{2000}{mm} rainfall (May--Oct), distinct dry season (Nov--Apr). Mean temp \SI{22}{\ensuremath{{}^{\circ}\mathrm{C}}}.
\item Soil: Vertisol. 0--30 cm: Dark grey clay, strong coarse prismatic breaking to angular blocky, very hard dry, very sticky/plastic wet, slickensides at 40 cm. 30--100 cm+: Grey clay, slickensides, few CaCO$_3$ nodules.
\item Drainage: Imperfect to poor. Flooded 3--4 months/year (Jun--Sep). Water table at surface in wet season, \SI{1.5}{m} deep in dry season.
\item Chemical: pH 6.5--7.2, CEC 45 cmol$_c$/kg (smectite), Base Saturation 85\%, ESP 5\%. Available P low (5 mg/kg Olsen).
\item Current Land Use: Grazing, recession sorghum/maize on residual moisture.
\end{itemize}
Assess the \textbf{Current Suitability} and \textbf{Potential Suitability (with High Inputs - Level C)} for \textbf{Irrigated Paddy Rice (two crops/year)}. Use FAO Framework classes S1, S2, S3, N1, N2.

\tcblower
\textbf{Detailed Solution:}

\noindent\textbf{1. Land Use Requirements (LURs) for Irrigated Rice (High Input):}
\begin{itemize}
\item \textbf{Topography (t):} Flat (0--1\%) essential for water control. $\rightarrow$ S1.
\item \textbf{Wetness/Oxygen (w):} Requires controlled flooding (5--10 cm) during vegetative/reproductive stages; drainage at harvest. Tolerates poor drainage \emph{if} managed. $\rightarrow$ S1 (with water control).
\item \textbf{Soil Physical (r):} Deep ($>$ 50 cm), heavy texture (clay) preferred to minimise percolation losses. High water holding. $\rightarrow$ S1 (Vertisol is deep clay).
\item \textbf{Fertility (n):} High CEC, high BS good. N, P, K responsive. Zn deficiency common in high pH flooded soils. $\rightarrow$ S1/S2 (P, Zn limitation correctable).
\item \textbf{Salinity/Sodicity (s):} ESP $<$ 15\%, EC low. $\rightarrow$ S1.
\item \textbf{Workability (k):} Very hard when dry, very sticky when wet. Difficult for land prep (puddling) without heavy machinery/timing. $\rightarrow$ S2/S3 limitation.
\end{itemize}

\noindent\textbf{2. Current Suitability (Rainfed / Low Input - Level A):}
\begin{center}
\begin{tabularx}{\linewidth}{|l|l|l|X|}\hline
\rowcolor{popblueL}
\textbf{Land Quality} & \textbf{Rating} & \textbf{Class} & \textbf{Reason} \\ \hline
Moisture (m) & S3 / N1 & \textbf{N1} & No irrigation infrastructure. Dry season crop impossible. Wet season flooding uncontrolled (too deep/long). \\ \hline
Oxygen (w) & S3 & \textbf{S3} & Uncontrolled flooding $\rightarrow$ deep water $\rightarrow$ only floating rice possible. \\ \hline
Rooting (r) & S1 & S1 & Deep clay. \\ \hline
Nutrients (n) & S2 & S2 & Low P, adequate N (mineralisation), high K/Ca/Mg. \\ \hline
Workability (k) & S3 & S3 & Vertisol: narrow moisture range for tillage. Manual labour impossible when wet/dry. \\ \hline
Erosion (t) & S1 & S1 & Flat. \\ \hline
\multicolumn{3}{|l|}{\textbf{Overall Current Suitability}} & \textbf{N1 (Currently Not Suitable)} for \emph{irrigated double-crop rice}. \\ \hline
\end{tabularx}
\end{center}
\textit{Note: Currently suitable only for recession agriculture (S2/S3) or grazing.}

\noindent\textbf{3. Potential Suitability (High Inputs - Level C: Irrigation + Drainage + Mechanisation + Fertiliser):}
\begin{itemize}
\item \textbf{Major Improvements:} Construction of irrigation canals, drainage ditches, land levelling, bunding. Provision of pumps/weirs. Mechanised puddling (tractors). Fertiliser (NPK + Zn).
\item \textbf{Re-rating:}
      \begin{itemize}
      \item Moisture (m): Controlled $\rightarrow$ \textbf{S1}.
      \item Oxygen (w): Controlled flooding/drainage $\rightarrow$ \textbf{S1}.
      \item Workability (k): Mechanised puddling at correct moisture $\rightarrow$ \textbf{S2} (Vertisol workability remains a minor constraint: timing critical, power requirement high).
      \item Nutrients (n): Fertiliser + Zn $\rightarrow$ \textbf{S1}.
      \item Salinity risk: Monitor with irrigation $\rightarrow$ \textbf{S1} (if water quality good).
      \end{itemize}
\end{itemize}

\noindent\textbf{Overall Potential Suitability:} \textbf{S2 (Moderately Suitable)}.
\textbf{Limiting Subclass:} \textbf{S2k} (Workability/Management constraint due to Vertisol properties).
\textbf{Justification:} Even with full water control and inputs, the Vertisol's extreme consistence (very hard dry / very sticky wet) requires high energy input for puddling and precise timing. Risk of structural damage if worked too wet. This prevents S1 classification.

\noindent\textbf{Exam Note:} Always distinguish \textbf{Current} (existing conditions) vs \textbf{Potential} (after specified improvements). State the input level assumed for Potential. The "Limiting Factor" rule means one S2 drops the whole map unit to S2.
\end{exoresolu}

\begin{methode}[Designing Soil Conservation Measures (Erosion Control)]
\begin{enumerate}
\item \textbf{Estimate Soil Loss Risk (USLE/RUSLE factors):} $A = R \times K \times LS \times C \times P$.
      \begin{itemize}
      \item \textbf{R} (Rainfall Erosivity): High in Cameroon coastal/mountains (\textgreater 500 MJ mm ha$^{-1}$ h$^{-1}$ yr$^{-1}$).
      \item \textbf{K} (Soil Erodibility): High for high silt, low OM, weak structure (e.g., 0.3--0.4). Low for clay/OM rich (0.1--0.2).
      \item \textbf{LS} (Slope Length/Steepness): Increases exponentially with slope \% and length.
      \item \textbf{C} (Cover/Management): 0.001 (dense forest) to 1.0 (bare fallow). Key management lever.
      \item \textbf{P} (Support Practice): 1.0 (up-down slope) to 0.1--0.5 (contouring, terracing, strips).
      \end{itemize}
\item \textbf{Compare Predicted Loss (A) to Soil Loss Tolerance (T):} T = \SIrange{2}{11}{t.ha^{-1}.yr^{-1}} (deep soils \num{11}, shallow \num{2}). If $A > T \rightarrow$ Conservation needed.
\item \textbf{Select Measures based on Slope, Soil, Land Use, Socio-economics:}
      \begin{itemize}
      \item \textbf{Agronomic:} Cover crops, mulching, residue retention, agroforestry, contour farming (C factor reduction).
      \item \textbf{Mechanical/Structural:} Contour bunds/ridges, graded bunds/channel terraces, bench terraces (steep), grass strips (Vetiver), check dams (gullies).
      \item \textbf{Biological:} Alley cropping (hedgerows), windbreaks, reforestation.
      \end{itemize}
\item \textbf{Design Parameters (Rules of Thumb):}
      \begin{itemize}
      \item \textbf{Contour Bunds/Vetiver Strips:} Vertical Interval (VI) in cm. Horizontal Interval (HI) in m. $\text{HI} = \frac{\text{VI}}{S\%} \times 100$. Typical VI = \SIrange{60}{100}{cm} for erosion control; \SIrange{100}{150}{cm} for Vetiver.
      \item \textbf{Graded Channel Terraces:} Grade 0.1--0.3\%. Capacity for 10--20 yr storm.
      \item \textbf{Bench Terraces:} Slope $>$ 20--25\%. Width $\ge$ \SI{3}{m}.
      \end{itemize}
\end{enumerate}
\end{methode}

\begin{exoresolu}[Erosion Control Design for Steep Coffee Farms in the West Region]
\textbf{Statement:} Smallholder coffee farms on the western slopes of the Bamboutos Mountains (West Region) are on \SI{35}{\%} slopes. Soils are deep Humic Nitisols (clay loam, high OM, stable structure). Rainfall erosivity $R = \SI{800}{MJ.mm.ha^{-1}.h^{-1}.yr^{-1}}$. Soil erodibility $K = \SI{0.15}{t.ha.h.ha^{-1}.MJ^{-1}.mm^{-1}}$ (low due to clay/OM). Current practice: Clean weeding (bare soil between coffee rows), up-and-down slope cultivation. Coffee rows spaced \SI{3}{m} apart. Slope length \SI{50}{m}.
\begin{enumerate}
\item Calculate the current annual soil loss (A) using USLE. (LS factor for 35\% slope, 50 m length $\approx$ 12.0).
\item Determine if conservation is needed (Tolerance T = \SI{10}{t.ha^{-1}.yr^{-1}} for deep soil).
\item Propose a specific, integrated conservation plan (agronomic + structural + biological) suitable for smallholders. Calculate the expected soil loss with your plan (estimate new C and P factors).
\end{enumerate}

\tcblower
\textbf{Detailed Solution:}

\noindent\textbf{1. Current Soil Loss (USLE):}
\[
A = R \times K \times LS \times C \times P
\]
\begin{itemize}
\item $R = 800$
\item $K = 0.15$
\item $LS = 12.0$ (given for 35\%, 50 m)
\item $C = 0.35$ (Clean weeded perennial crop, row crop equivalent? Coffee with bare inter-row $\approx$ 0.3--0.4. Use \textbf{0.35}).
\item $P = 1.0$ (Up-and-down slope cultivation).
\end{itemize}
\[
A = 800 \times 0.15 \times 12.0 \times 0.35 \times 1.0 = \mathbf{504\ t.ha^{-1}.yr^{-1}}
\]

\noindent\textbf{2. Comparison with Tolerance:}
$T = \SI{10}{t.ha^{-1}.yr^{-1}}$.
$A (504) \gg T (10)$. \textbf{Severe erosion risk.} Unsustainable. Current loss $\approx$ \SI{3.7}{mm/yr} (assuming $\rho_b = 1.35$). Entire topsoil lost in decades.

\noindent\textbf{3. Integrated Conservation Plan:}

\begin{center}
\begin{tabularx}{\linewidth}{|l|X|l|}\hline
\rowcolor{popblueL}
\textbf{Measure Type} & \textbf{Specific Action} & \textbf{Effect on Factors} \\ \hline
\textbf{Agronomic (C)} & \textbf{Mulching:} Apply coffee pulp/husk mulch (\SI{5--10}{t/ha}) in inter-rows. \textbf{Cover crop:} Plant \textit{Pueraria phaseoloides} or \textit{Desmodium} in inter-rows (suppresses weeds, fixes N). \textbf{Pruning residues:} Spread on contour lines. & \textbf{C $\downarrow$ 0.35 $\rightarrow$ 0.05} (Good cover + mulch). \\ \hline
\textbf{Structural (P)} & \textbf{Contour Planting:} Align new coffee rows on contour. \textbf{Vetiver Grass Strips:} Plant on contour at \textbf{Vertical Interval (VI) = 1 m}. 
      \begin{itemize}
      \item HI = $\frac{VI}{S\%} \times 100 = \frac{100\ cm}{35} \times 100 \approx \SI{2.86}{m}$.
      \item Plant strips every \SI{3}{m} horizontal (matches coffee row spacing).
      \end{itemize}
      \textbf{Infiltration pits (Zai):} \SI{30x30x30}{cm} pits on contour between rows to harvest runoff. & \textbf{P $\downarrow$ 1.0 $\rightarrow$ 0.30} (Contouring + Vetiver strips + pits). \\ \hline
\textbf{Biological} & \textbf{Shade Trees:} Maintain/plant \textit{Albizia}, \textit{Gliricidia} (nitrogen fixing) at \SI{10x10}{m}. Deep roots stabilise soil, litter improves OM. \textbf{Grassed waterways:} Stabilise natural drainage lines with \textit{Brachiaria} / Vetiver. & Long-term structure/OM improvement $\rightarrow$ K may decrease further. \\ \hline
\end{tabularx}
\end{center}

\noindent\textbf{4. Predicted Soil Loss with Plan:}
\[
A_{new} = 800 \times 0.15 \times 12.0 \times \mathbf{0.05} \times \mathbf{0.30} = \mathbf{21.6\ t.ha^{-1}.yr^{-1}}
\]
\textit{Still slightly above T (10). Further reduction needed.}

\noindent\textbf{Refinement:} Increase mulch/cover to achieve $C = 0.02$ (excellent cover). Vetiver strips mature $\rightarrow$ $P = 0.20$.
\[
A_{opt} = 800 \times 0.15 \times 12.0 \times 0.02 \times 0.20 = \mathbf{5.76\ t.ha^{-1}.yr^{-1}} < T.
\]
\textbf{Achievable with good management.}

\noindent\textbf{Key Exam Points:} 
\begin{itemize}
\item Show USLE calculation clearly.
\item VI/HI calculation for Vetiver/Contour bunds is a favourite exam question: $HI = \frac{VI}{S} \times 100$ (VI in cm, S in \%, HI in m).
\item Integrated approach (Agronomic + Structural + Biological) scores higher than just listing structures.
\item Mention socio-economic feasibility: Vetiver is cheap, uses local labour, provides thatch/fodder.
\end{itemize}
\end{exoresolu}

\begin{methode}[Diagnosing and Managing Soil Salinity / Sodicity (Irrigated Lands)]
\begin{enumerate}
\item \textbf{Obtain saturation extract data:} EC$_e$ (Electrical Conductivity, \SI{}{dS.m^{-1}}), SAR (Sodium Adsorption Ratio), pH$_e$, Exchangeable Na\% (ESP).
\item \textbf{Classify the Soil:}
      \begin{center}
      \begin{tabular}{|l|c|c|c|}\hline
      \rowcolor{popblueL}
      \textbf{Class} & \textbf{EC$_e$ (dS/m)} & \textbf{ESP / SAR} & \textbf{pH} \\\hline
      Normal / Non-saline & $<$ 4 & $<$ 15 / $<$ 13 & $<$ 8.5 \\\hline
      \textbf{Saline} & $\ge$ 4 & $<$ 15 / $<$ 13 & $<$ 8.5 (usually) \\\hline
      \textbf{Sodic (Alkali)} & $<$ 4 & $\ge$ 15 / $\ge$ 13 & $>$ 8.5 (often 8.5--10) \\\hline
      \textbf{Saline-Sodic} & $\ge$ 4 & $\ge$ 15 / $\ge$ 13 & Variable ($<$ 8.5 if saline dominant) \\\hline
      \end{tabular}
      \end{center}
\item \textbf{Assess Crop Tolerance:} Match EC$_e$ to crop threshold (e.g., Rice \SI{3}{dS/m}, Maize \num{1.7}, Beans \num{1.0}, Cotton \num{7.7}, Barley \num{8.0}). Yield decline = slope $\times$ (EC$_e$ - threshold).
\item \textbf{Reclamation Strategy:}
      \begin{itemize}
      \item \textbf{Saline Soil:} \textbf{Leaching.} Apply excess water (Leaching Requirement LR = $\frac{EC_{iw}}{EC_{dw}}$ or $LR = \frac{EC_e}{EC_{dw}}$). Good drainage essential. No chemical amendment needed.
      \item \textbf{Sodic Soil:} \textbf{Chemical Amendment + Leaching.} Add Ca source (Gypsum CaSO$_4\cdot$2H$_2$O, or Sulphur/acid if calcareous) to replace Na$^+$ on exchange complex: $2Na^+_{ex} + Ca^{2+}_{sol} \rightarrow Ca^{2+}_{ex} + 2Na^+_{sol}$. Then leach Na$^+$ out. Deep ploughing to improve structure.
      \item \textbf{Saline-Sodic:} \textbf{Amendment first (Gypsum), then Leaching.} Gypsum prevents dispersion during leaching.
      \end{itemize}
\item \textbf{Calculate Gypsum Requirement (GR):} 
      \[
      GR (\text{t/ha, 30 cm}) = \frac{(ESP_{initial} - ESP_{target}) \times CEC \times 0.86 \times \text{depth(cm)} \times \rho_b}{100 \times \text{Purity}\%}
      \]
      (Constant 0.86 derives from molar mass: 1 cmol$_c$ Na $\rightarrow$ 0.086 kg pure gypsum; factor 100 for ha-cm conversion).
\item \textbf{Long-term Management:} Salt-tolerant varieties, drip irrigation (minimises leaching fraction but avoids surface accumulation), organic matter, drainage maintenance.
\end{enumerate}
\end{methode}

\begin{exoresolu}[Reclamation of a Sodic Soil in the Logone Floodplain]
\textbf{Statement:} A farmer in the Logone Valley (Far North) wants to cultivate maize on a soil affected by sodicity. Soil analysis (0--30 cm) shows:
\begin{itemize}
\item pH (paste) = 9.2
\item EC$_e$ = \SI{2.5}{dS.m^{-1}}
\item CEC = \SI{35}{cmol_{c}/kg} (Smectitic clay)
\item Exchangeable Na = \SI{8.0}{cmol_{c}/kg}
\item Exchangeable Ca + Mg = \SI{25.0}{cmol_{c}/kg}
\item ESP = 23\%
\item SAR (saturation extract) = 25
\end{itemize}
Maize threshold EC$_e$ = \SI{1.7}{dS.m^{-1}}, ESP tolerance $<$ 10\%.
\begin{enumerate}
\item Classify the soil salinity/sodicity status.
\item Calculate the Gypsum Requirement (tonnes/ha) to reduce ESP to 5\% in the 0--30 cm layer. (Assume Gypsum purity 85\%, Soil bulk density \SI{1.4}{g.cm^{-3}}).
\item Outline the reclamation procedure steps.
\end{enumerate}

\tcblower
\textbf{Detailed Solution:}

\noindent\textbf{1. Classification:}
\begin{itemize}
\item EC$_e$ = 2.5 dS/m $<$ 4 $\rightarrow$ \textbf{Non-saline}.
\item ESP = 23\% $\ge$ 15\% $\rightarrow$ \textbf{Sodic}.
\item pH = 9.2 $>$ 8.5 $\rightarrow$ Confirms sodic (hydrolysis of Na$_2$CO$_3$ $\rightarrow$ OH$^-$).
\item \textbf{Diagnosis: SODIC SOIL (Non-saline, Sodic).}
\end{itemize}

\noindent\textbf{2. Gypsum Requirement (GR) Calculation:}
\textbf{Method 1: Molar Stoichiometry (Rigorous)}
\begin{itemize}
\item Target ESP = 5\%. Current ESP = 23\%. Reduction needed = 18\% of CEC.
\item Na to replace = $0.18 \times CEC = 0.18 \times 35 = \SI{6.3}{cmol_{c}/kg\ soil}$.
\item Reaction: $2Na^+_{ex} + CaSO_4 \rightarrow Ca^{2+}_{ex} + 2Na^+ + SO_4^{2-}$.
\item 1 mole CaSO$_4$ replaces 2 moles Na$^+$ (2 cmol$_c$).
\item Moles CaSO$_4$ needed per kg soil = $\frac{6.3}{2} = \SI{3.15}{mmol/kg} = \SI{3.15}{cmol/kg}$.
\item Molar mass Gypsum (CaSO$_4\cdot$2H$_2$O) = \SI{172}{g/mol}.
\item Mass pure gypsum per kg soil = $3.15 \times 10^{-3}\ mol/kg \times 172\ g/mol = \SI{0.542}{g/kg}$.
\item Soil mass in 30 cm layer per ha = $10,000\ m^2 \times 0.30\ m \times 1.4\ Mg/m^3 = \SI{4200}{Mg} = \SI{4200000}{kg}$.
\item Total pure gypsum = $0.542\ g/kg \times 4,200,000\ kg = \SI{2276400}{g} = \SI{22.76}{t}$.
\item Adjust for purity (85\%): $GR = \frac{22.76}{0.85} = \mathbf{26.8\ t/ha}$.
\end{itemize}

\textbf{Method 2: Standard Formula (Verification)}
\[
GR = \frac{(23 - 5) \times 35 \times 0.86 \times 30 \times 1.4}{100 \times 0.85} = \frac{18 \times 35 \times 0.86 \times 30 \times 1.4}{85} = \frac{22755.6}{85} = \mathbf{26.8\ t/ha}
\]

\noindent\textbf{3. Reclamation Procedure:}
\begin{enumerate}
\item \textbf{Drainage Assessment:} Ensure adequate natural/artificial drainage (outlet for leachate). Sodic soils disperse and seal without drainage $\rightarrow$ failure.
\item \textbf{Land Preparation:} Deep ploughing (\SIrange{30}{40}{cm}) to break any hardpan, mix amendment.
\item \textbf{Gypsum Application:} Broadcast \textbf{10 t/ha} (Year 1), incorporate by disc harrowing. Irrigate heavily (\SIrange{100}{150}{mm}) to dissolve and move Ca$^{2+}$ into exchange complex.
\item \textbf{Leaching:} Pond water (\SIrange{200}{300}{mm}) to flush displaced Na$^+$ below root zone. Monitor drainage water EC/SAR.
\item \textbf{Green Manure:} Grow salt-tolerant green manure (Sesbania, Dhaincha) during reclamation period. Incorporate before flowering $\rightarrow$ adds OM, produces CO$_2$/organic acids $\rightarrow$ helps dissolve native CaCO$_3$ if present, improves structure.
\item \textbf{Second Gypsum Application:} \textbf{10 t/ha} (Year 2) if ESP still $>$ 10\%.
\item \textbf{Crop Establishment:} Plant maize (moderately tolerant) on ridges. Use starter fertiliser (DAP/NPK). Avoid urea (ammonium $\rightarrow$ NH$_3$ volatilisation at high pH). Use Nitrate forms or incorporate urea deep.
\item \textbf{Maintenance:} Regular organic matter addition. Monitor ESP/pH annually. Avoid irrigation water with high SAR (RSC - Residual Sodium Carbonate).
\end{enumerate}

\noindent\textbf{Exam Tip:} Always check \textbf{Drainage} first. Reclaiming sodic soil without drainage is wasting gypsum. The reaction $2Na_{ex} + Ca_{sol} \rightarrow Ca_{ex} + 2Na_{sol}$ is the core chemistry. Know the molar mass of gypsum (172) and the 2:1 stoichiometry.
\end{exoresolu}

\begin{methode}[Evaluating Soil Organic Matter Dynamics and Carbon Sequestration]
\begin{enumerate}
\item \textbf{Understand the Carbon Cycle in Soil:} Inputs (Root residues, Shoot residues, Manure, Compost, Biochar) $\rightarrow$ Decomposition (Microbial respiration $\rightarrow$ CO$_2$) $\rightarrow$ Stabilisation (Humification $\rightarrow$ Humus, Organo-mineral complexes, Physical protection in aggregates) $\rightarrow$ Outputs (CO$_2$, CH$_4$, DOC leaching, Erosion).
\item \textbf{Key Metrics:}
      \begin{itemize}
      \item \textbf{SOC Stock (t C/ha)} = $\%C \times \rho_b (g/cm^3) \times Depth (cm) \times 100$ (for 1 ha = 10,000 m$^2$). Formula: $Stock = \%C \times \rho_b \times d \times 100$.
      \item \textbf{C:N Ratio:} $<$ 20:1 $\rightarrow$ Net mineralisation (N release). $>$ 25:1 $\rightarrow$ Net immobilisation (N lock-up). Optimal for compost $\approx$ 25--30:1.
      \item \textbf{Decomposition Rate (k):} First order kinetics: $C_t = C_0 e^{-kt}$. Half-life $t_{1/2} = \ln(2)/k$. $k$ depends on climate (T, W), quality (lignin, polyphenols), soil protection.
      \item \textbf{4 per 1000 Initiative:} Target 0.4\% annual increase in SOC stock for climate mitigation/food security.
      \end{itemize}
\item \textbf{Management Practices to Increase SOC:}
      \begin{itemize}
      \item \textbf{Maximise Inputs:} High biomass crops, cover crops, agroforestry, residue retention, organic amendments (manure, compost, biochar).
      \item \textbf{Minimise Losses:} No-till / Reduced tillage (reduces oxidation/erosion), Continuous cover, Agroforestry (deep roots), Erosion control.
      \item \textbf{Enhance Stabilisation:} Clay mineral management (less relevant), Promote aggregate formation (fungal hyphae, roots, OM), Biochar (recalcitrant C).
      \end{itemize}
\item \textbf{Trade-offs:} N$_2$O emissions (denitrification in wet, high N soils), CH$_4$ (waterlogged), Nutrient availability timing (immobilisation).
\end{enumerate}
\end{methode}

\begin{exoresolu}[Soil Organic Carbon Stock Calculation and 4p1000 Target]
\textbf{Statement:} A research plot in the humid forest zone (Dja Reserve buffer) compares two systems on a Ferralsol (Oxisol):
\begin{itemize}
\item \textbf{System A:} Continuous maize with residue burning. $\rho_b = \SI{1.35}{g.cm^{-3}}$, SOC = 1.2\% (0--30 cm).
\item \textbf{System B:} Maize-Legume rotation (Mucuna cover crop) + No-till + Residue retention. $\rho_b = \SI{1.15}{g.cm^{-3}}$, SOC = 2.1\% (0--30 cm).
\end{itemize}
\begin{enumerate}
\item Calculate the SOC stock (t C/ha) in the 0--30 cm layer for each system.
\item Calculate the additional C sequestered in System B compared to A.
\item If System A adopts System B practices, estimate the time (years) to reach System B's SOC level, assuming a net sequestration rate of \SI{0.5}{t C.ha^{-1}.yr^{-1}} (typical for humid tropics with cover crops/no-till).
\item Discuss the co-benefits of this transition for soil physical/chemical properties in a Ferralsol.
\end{enumerate}

\tcblower
\textbf{Detailed Solution:}

\noindent\textbf{1. SOC Stock Calculation (0--30 cm):}
Formula: $Stock (t C/ha) = \%C \times \rho_b (g/cm^3) \times Depth (cm) \times 100$.
(Derivation: $1 ha = 10^8 cm^2$. Vol = $10^8 \times d (cm) = 10^8 d\ cm^3$. Mass = Vol $\times \rho_b = 10^8 d \rho_b\ g$. C mass = Mass $\times \%C/100 = 10^6 d \rho_b \%C\ g = d \rho_b \%C \times 100\ t$).

\begin{itemize}
\item \textbf{System A:} $Stock_A = 1.2 \times 1.35 \times 30 \times 100 = \mathbf{48.6\ t\ C/ha}$.
\item \textbf{System B:} $Stock_B = 2.1 \times 1.15 \times 30 \times 100 = \mathbf{72.45\ t\ C/ha}$.
\end{itemize}

\noindent\textbf{2. Additional C Sequestered:}
$\Delta Stock = 72.45 - 48.6 = \mathbf{23.85\ t\ C/ha}$.
\textit{Note: System B also has lower bulk density, meaning more soil volume per hectare, amplifying the stock difference beyond just the \%C difference.}

\noindent\textbf{3. Time to Reach Equilibrium:}
$\Delta Stock = \SI{23.85}{t C/ha}$.
Rate = \SI{0.5}{t C/ha/yr}.
$Time = \frac{23.85}{0.5} = \mathbf{47.7\ years} \approx \mathbf{48\ years}$.
\textit{Reality check:} Sequestration rate slows as soil approaches new saturation (C saturation limit $\approx$ f(clay, mineralogy)). First 10--20 years fastest. 48 years is a reasonable long-term estimate for full equilibration.

\noindent\textbf{4. Co-benefits for Ferralsol (System B vs A):}
\begin{center}
\begin{tabularx}{\linewidth}{|l|X|X|}\hline
\rowcolor{popblueL}
\textbf{Property} & \textbf{System A (Burn/Till)} & \textbf{System B (No-till/Cover)} \\ \hline
\textbf{Aggregation / Structure} & Weak, massive, prone to slaking/crusting. Low porosity. & \textbf{Strong granular/blocky.} Fungal hyphae, roots, polysaccharides bind microaggregates. High macroporosity. \\ \hline
\textbf{Water Infiltration / Storage} & Low infiltration, high runoff/erosion. Low plant-available water. & \textbf{High infiltration,} reduced runoff. Higher porosity $\rightarrow$ more storage. Residue mulch reduces evaporation. \\ \hline
\textbf{Nutrient Retention (CEC)} & Low effective CEC (variable charge). High leaching of N, K. & \textbf{OM contributes 50--70\% of CEC} in Ferralsols. Higher ECEC, better nutrient holding. \\ \hline
\textbf{pH / Al Toxicity} & Acidification from nitrification/leaching. Al saturation high. & \textbf{OM buffers pH,} complexes Al$^{3+}$ (detoxification). Higher pH. \\ \hline
\textbf{Biological Activity} & Low (food/habitat scarce, high temp fluctuations). & \textbf{High.} Earthworms, termites, microbes thrive. Nutrient cycling rapid. \\ \hline
\textbf{Erosion} & Severe (bare soil, crusting, steep slopes common). & \textbf{Minimal.} Continuous cover + stable aggregates + macropores. \\ \hline
\end{tabularx}
\end{center}

\noindent\textbf{Conclusion:} Transition to Conservation Agriculture (System B) is a \textbf{win-win-win}: Climate mitigation (C sequestration), Adaptation (resilience to drought/flood), and Productivity (sustained yields with lower external inputs). This aligns with Cameroon's commitments to UNFCCC and LDN (Land Degradation Neutrality).
\end{exoresolu}

\begin{aretenir}[Summary of Key Methods for GCE Exam]
\begin{itemize}
\item \textbf{Soil Formation:} Master CLORPT $\rightarrow$ Process $\rightarrow$ Soil Type. Use Cameroon examples (Mount Cameroon, Bamenda, Logone, Douala).
\item \textbf{Profile Description:} Systematic observation (Depth, Colour, Texture, Structure, Consistence, Roots, Boundary) $\rightarrow$ Horizon Symbols (Master + Subordinate) $\rightarrow$ Diagnostic Horizons $\rightarrow$ Classification (WRB: Ferralsol, Nitisol, Luvisol/Lixisol, Acrisol, Vertisol, Gleysol, Andosol, Fluvisol, Cambisol, Arenosol, Regosol, Leptosol, Plinthosol).
\item \textbf{Texture:} Calculate Sand/Silt/Clay \% $\rightarrow$ Textural Triangle $\rightarrow$ Class $\rightarrow$ Management implications.
\item \textbf{Physical Properties:} $\rho_b = M_s/V_t$; $f = 1 - \rho_b/\rho_s$; $\theta_g = M_w/M_s$; $\theta_v = \theta_g \rho_b$; AWC = FC - PWP; MAD = fraction $\times$ AWC $\times$ Depth.
\item \textbf{Chemical Interpretation:} pH, $\Delta$pH, CEC vs ECEC, BS\%, Al Sat, ESP, Ca:Mg, (Ca+Mg):K, P availability. \textbf{Ferralsols/Acrisols: Use ECEC and Al Sat.}
\item \textbf{Land Evaluation:} FAO Framework. Match LQ to LUR. Limiting Factor Principle. Current vs Potential Suitability. Subclass limitations (m, w, n, r, t, s, k).
\item \textbf{Conservation:} USLE ($A=RKLSCP$). $T$ values. Agronomic (C), Structural (P), Biological measures. Vetiver VI/HI calculation.
\item \textbf{Salinity/Sodicity:} EC$_e$, ESP, SAR, pH classification. Saline $\rightarrow$ Leach. Sodic $\rightarrow$ Gypsum + Leach. GR calculation.
\item \textbf{SOC:} Stock = $\%C \times \rho_b \times d \times 100$. 4p1000. Management: Inputs $\uparrow$, Losses $\downarrow$, Stabilisation $\uparrow$. Co-benefits.
\end{itemize}
\end{aretenir}

\newpage

\coursec{Exercises}

\courssub{I check my knowledge}

\begin{exercice}[Multiple choice questions]
For each question, choose the single correct answer.
\begin{enumerate}
\item The layer of the soil profile in which soluble minerals are washed downwards by percolating water is:
\begin{enumerate}
\item the A horizon (eluviation);
\item the B horizon (illuviation);
\item the C horizon (weathered parent rock);
\item the R horizon (bedrock).
\end{enumerate}
\item A soil containing approximately 40\% sand, 40\% silt and 20\% clay is best described as:
\begin{enumerate}
\item a sandy soil;
\item a clay soil;
\item a loam;
\item a peat.
\end{enumerate}
\item The dark colour of the topsoil in a forest soil is mainly due to:
\begin{enumerate}
\item iron oxides;
\item humus;
\item calcium carbonate;
\item quartz grains.
\end{enumerate}
\item Lateritic soils of the South Region of Cameroon owe their red colour to:
\begin{enumerate}
\item organic matter;
\item hydrated iron oxides (goethite, lepidocrocite);
\item gypsum;
\item volcanic ash.
\end{enumerate}
\end{enumerate}
\end{exercice}

\begin{exercice}[True or false]
State whether each statement is true or false, and justify your answer in one or two sentences.
\begin{enumerate}
\item Soil formation is a rapid process: one centimetre of topsoil can form in a single year.
\item The B horizon is enriched in clay, iron and aluminium compounds leached from the horizon above.
\item A sandy soil retains water better than a clay soil.
\item Bush fallowing and crop rotation help to maintain soil fertility.
\item Once severely eroded, a soil can be fully restored within a few farming seasons.
\end{enumerate}
\end{exercice}

\begin{exercice}[Key definitions]
Define each of the following terms in one or two clear sentences, as you would in the GCE examination:
\begin{enumerate}
\item soil;
\item humus;
\item eluviation;
\item illuviation;
\item soil profile;
\item lateritic soil.
\end{enumerate}
\end{exercice}

\begin{exercice}[Direct calculations]
A soil sample of mass $250\ \mathrm{g}$ was collected from a farm near Bafoussam. After oven-drying at $105^\circ\mathrm{C}$ until constant mass, its mass was $210\ \mathrm{g}$.
\begin{enumerate}
\item Calculate the mass of water lost by the sample.
\item Calculate the percentage moisture content of the fresh soil (on a dry-mass basis, as standard in soil science).
\item After the jar settling test on $100\ \mathrm{g}$ of the dried soil, the layers measured were: sand $55\ \mathrm{g}$, silt $30\ \mathrm{g}$, clay $15\ \mathrm{g}$. Calculate the percentage of each particle-size fraction.
\item Using the textural triangle studied in class, state the textural class of this soil.
\end{enumerate}
\end{exercice}

\courssub{I practise}

\begin{exercice}[Interpreting a soil profile]
The table below describes the layers of a soil profile observed in a road cutting near Ebolowa, but the horizons have been scrambled.
\begin{center}
\begin{tabularx}{\linewidth}{|c|X|}
\hline
\rowcolor{popblueL} Layer & Description \\
\hline
1 & Unweathered solid bedrock, no biological activity \\
\hline
2 & Dark grey layer rich in humus and living organisms (mineral matter present) \\
\hline
3 & Partly weathered rock fragments mixed with little mineral matter \\
\hline
4 & Red-brown layer enriched in clay and iron oxides washed from above \\
\hline
5 & Light-coloured layer strongly leached of soluble minerals and clay \\
\hline
\end{tabularx}
\end{center}
\begin{enumerate}
\item Identify the horizon (O, A, E, B, C or R) corresponding to each layer.
\item Redraw the profile in its correct vertical order, from the surface to the bedrock.
\item State one characteristic of each horizon.
\item Name and define the two processes responsible for the difference between layers 5 and 4.
\end{enumerate}
\end{exercice}

\begin{exercice}[Textural analysis of three soils]
Three soil samples were analysed in the school laboratory. The results are given below.
\begin{center}
\begin{tabularx}{\linewidth}{|l|X|X|X|}
\hline
\rowcolor{popblueL} Sample & Sand (\%) & Silt (\%) & Clay (\%) \\
\hline
A (Maroua, Far North) & 75 & 15 & 10 \\
\hline
B (Mbalmayo, Centre) & 20 & 30 & 50 \\
\hline
C (Buea, South-West) & 40 & 40 & 20 \\
\hline
\end{tabularx}
\end{center}
\begin{enumerate}
\item Locate each sample on the textural triangle and give its textural class.
\item For each sample, state its drainage and water-retention behaviour.
\item Which sample is best suited for maize cultivation? Justify your answer with two reasons.
\item Suggest one amendment that could improve sample A for farming.
\end{enumerate}
\end{exercice}

\begin{exercice}[Simple soil tests]
A student carries out two tests on a garden soil from Yaoundé.
\begin{enumerate}
\item Describe, step by step, the procedure of the jar settling test and state what each settled layer represents.
\item The pH of the soil, measured with a pH meter in a soil--water suspension (1:2.5), is 5.2. Is this soil acidic, neutral or alkaline?
\item Name two crops that tolerate such a pH and one amendment a farmer could apply to raise the pH.
\item Explain why knowing soil pH is important before planting.
\end{enumerate}
\end{exercice}

\begin{exercice}[Diagnosing a soil problem]
A farmer in the Far North Region observes declining millet yields, a hard surface crust after the rare rains, and dust clouds carrying away topsoil during the dry season.
\begin{enumerate}
\item Identify the two main degradation processes affecting this farm.
\item Give two likely human causes of this degradation.
\item Propose four sustainable management techniques adapted to this semi-arid environment, and explain briefly how each one works.
\end{enumerate}
\end{exercice}

\courssub{I prepare for the GCE}

\begin{exercice}[Structured question: the soil profile (20 marks)]
The diagram below (to be drawn by the candidate) represents a mature soil profile developed on granite in the Centre Region of Cameroon.
\begin{enumerate}
\item Draw a labelled soil profile showing the O, A, E, B, C and R horizons. \hfill (5 marks)
\item For each horizon, state one physical or chemical characteristic. \hfill (5 marks)
\item Define the processes of eluviation and illuviation, and indicate the horizon in which each is dominant. \hfill (4 marks)
\item Explain the role of (i) the parent rock and (ii) climate in the development of this profile. \hfill (4 marks)
\item State one reason why the A horizon is the most important layer for agriculture. \hfill (2 marks)
\end{enumerate}
\end{exercice}

\begin{exercice}[Essay: lateritic soils of the South Region (20 marks)]
\emph{``Explain how climate and living organisms contribute to the formation of lateritic soils in the South Region of Cameroon.''}

Your essay should:
\begin{enumerate}
\item describe the climatic conditions of the South Region relevant to soil formation; \hfill (4 marks)
\item explain the chemical weathering processes (hydrolysis, leaching of silica and bases, concentration of iron and aluminium oxides) that produce laterite; \hfill (8 marks)
\item explain the role of vegetation, micro-organisms and soil fauna in providing humus and mixing the soil; \hfill (5 marks)
\item state two agricultural limitations of lateritic soils and one way of managing them. \hfill (3 marks)
\end{enumerate}
\end{exercice}

\begin{exercice}[GCE-style discussion essay (25 marks)]
\emph{``Soil is a non-renewable resource on a human timescale.'' Discuss this statement with reference to soil formation, soil degradation and soil protection in Cameroon.}

In your answer you are expected to:
\begin{enumerate}
\item explain, with an order of magnitude, the time required for the natural formation of a fertile topsoil; \hfill (5 marks)
\item describe at least three degradation processes threatening Cameroonian soils (for example water erosion, wind erosion, loss of fertility, lateritic crusting, pollution), with a named region for each; \hfill (9 marks)
\item present at least four protection and sustainable management measures, explaining how each works; \hfill (8 marks)
\item conclude with a justified personal position on the statement. \hfill (3 marks)
\end{enumerate}
\end{exercice}

\newpage

\coursec{Solutions to the exercises}

\begin{corrige}[Exercise 1 — Multiple choice questions]
\begin{enumerate}
\item \textbf{Answer: (b) the B horizon (illuviation).}  
  \textit{Justification:} Eluviation is the removal of soluble minerals and fine particles from the upper horizons (A and E) by percolating water. Illuviation is the deposition of these materials in the B horizon, which therefore becomes enriched in clay, iron and aluminium compounds.

\item \textbf{Answer: (c) a loam.}  
  \textit{Justification:} A loam is a soil with roughly equal proportions of sand and silt and a smaller proportion of clay (typically 40\% sand, 40\% silt, 20\% clay). Sandy soil has >70\% sand; clay soil has >40\% clay; peat is organic.

\item \textbf{Answer: (b) humus.}  
  \textit{Justification:} Humus, the dark-coloured decomposed organic matter, gives the topsoil (A horizon) its dark grey to black colour. Iron oxides give red/brown colours, calcium carbonate gives white, quartz is colourless.

\item \textbf{Answer: (b) hydrated iron oxides (goethite, lepidocrocite).}  
  \textit{Justification:} Lateritic soils in the humid tropical South Region form under intense leaching. Silica and bases are removed, leaving behind insoluble iron and aluminium oxides. Hydrated iron oxides (goethite, lepidocrocite) impart the characteristic red-brown colour.
\end{enumerate}
\end{corrige}

\begin{corrige}[Exercise 2 — True or false]
\begin{enumerate}
\item \textbf{False.} Soil formation is an extremely slow process; it takes hundreds to thousands of years to form 1 cm of topsoil under natural conditions.
\item \textbf{True.} The B horizon (illuviation horizon) accumulates clay, iron oxides, aluminium oxides and organic matter translocated from the overlying A and E horizons.
\item \textbf{False.} Clay soils retain water much better than sandy soils because their tiny particles and pore spaces hold water by capillary forces; sandy soils have large pores and drain rapidly.
\item \textbf{True.} Bush fallowing (resting land under natural vegetation) and crop rotation (alternating crops with different nutrient demands) both help restore organic matter and nutrient balance, maintaining fertility.
\item \textbf{False.} Severe erosion removes the topsoil which took centuries to form; full restoration of soil structure, organic matter and fertility takes decades to centuries, not a few seasons.
\end{enumerate}
\end{corrige}

\begin{corrige}[Exercise 3 — Key definitions]
\begin{enumerate}
\item \textbf{Soil:} The unconsolidated mineral and organic material on the Earth's surface that supports plant growth, formed by weathering of parent rock under the influence of climate, organisms, topography and time (CLORPT factors).
\item \textbf{Humus:} Dark, amorphous, stable organic matter resulting from the decomposition of plant and animal residues by soil micro-organisms; it improves soil structure, water retention, cation exchange capacity and nutrient supply.
\item \textbf{Eluviation:} The downward movement of dissolved or suspended material (clay, iron, aluminium, humus, soluble salts) from upper soil horizons by percolating water.
\item \textbf{Illuviation:} The accumulation of materials (clay, iron/aluminium oxides, organic matter) that have been leached from upper horizons, occurring mainly in the B horizon.
\item \textbf{Soil profile:} A vertical section through the soil showing its succession of horizons (O, A, E, B, C, R) from the surface down to the unweathered parent rock.
\item \textbf{Lateritic soil:} A highly weathered, red tropical soil formed under hot humid climates, characterized by intense leaching of silica and bases (desilication, debasification), concentration of iron and aluminium oxides (ferralitization), low cation exchange capacity, and often a hard indurated layer (lateritic crust or petroplinthite).
\end{enumerate}
\end{corrige}

\begin{corrige}[Exercise 4 — Direct calculations]
\begin{enumerate}
\item \textbf{Mass of water lost} = initial mass – dry mass = $250\ \mathrm{g} - 210\ \mathrm{g} = \boxed{40\ \mathrm{g}}$.
\item \textbf{Moisture content (dry-mass basis)} = $\dfrac{\text{mass of water}}{\text{dry mass}} \times 100\% = \dfrac{40}{210} \times 100\% = \boxed{19.0\%}$.
\item \textbf{Particle-size percentages} (based on $100\ \mathrm{g}$ dried sample):  
  Sand = $\dfrac{55}{100} \times 100\% = \boxed{55\%}$;  
  Silt = $\dfrac{30}{100} \times 100\% = \boxed{30\%}$;  
  Clay = $\dfrac{15}{100} \times 100\% = \boxed{15\%}$.
\item \textbf{Textural class:} With 55\% sand, 30\% silt, 15\% clay, the soil plots in the \textbf{sandy loam} region of the USDA textural triangle.
\end{enumerate}
\end{corrige}

\begin{corrige}[Exercise 5 — Interpreting a soil profile]
\begin{enumerate}
\item \textbf{Horizon identification (matching exercise key):}
\begin{center}
\begin{tabular}{|c|c|}
\hline
\rowcolor{popblueL} Layer (from description) & Horizon \\
\hline
1 & R \\
2 & A (or O/A if organic layer separate) \\
3 & C \\
4 & B \\
5 & E \\
\hline
\end{tabular}
\end{center}
\item \textbf{Correct vertical order (surface $\to$ bedrock) with labelled diagram:}
\begin{popfigure}
\begin{tikzpicture}[scale=0.85]
\draw[thick] (0,0) rectangle (6,-7);
% O/A horizon
\fill[popgreenL] (0,0) rectangle (6,-1);
\node[align=center] at (3,-0.5) {O/A \\ (dark, humus-rich)};
% E horizon
\fill[popblueL] (0,-1) rectangle (6,-2);
\node[align=center] at (3,-1.5) {E \\ (light, leached)};
% B horizon
\fill[poporange] (0,-2) rectangle (6,-4);
\node[align=center] at (3,-3) {B \\ (red-brown, clay/Fe oxides)};
% C horizon
\fill[popmuted] (0,-4) rectangle (6,-6);
\node[align=center] at (3,-5) {C \\ (weathered rock fragments)};
% R horizon
\fill[popdark] (0,-6) rectangle (6,-7);
\node[align=center] at (3,-6.5) {R \\ (unweathered bedrock)};
\end{tikzpicture}
\legende{Idealised soil profile showing horizon sequence O/A–E–B–C–R.}
\end{popfigure}
\item \textbf{Characteristics of each horizon:}
\begin{itemize}
\item \textbf{O/A:} Dark colour, high organic matter (humus), intense biological activity (roots, fauna, microbes), granular/crumb structure, high nutrient availability.
\item \textbf{E:} Light colour (grey/white), leached of clay, Fe/Al oxides and organic matter (eluviation), platy structure, low organic matter, quartz-rich.
\item \textbf{B:} Red-brown to yellow-brown, accumulation of clay, Fe/Al oxides and organic matter (illuviation), blocky/prismatic structure, higher bulk density, lower porosity than A.
\item \textbf{C:} Partially weathered parent material, rock structure still visible (e.g., saprolite), little biological activity, mineral composition reflects parent rock.
\item \textbf{R:} Consolidated unweathered bedrock, no soil structure, no biological activity.
\end{itemize}
\item \textbf{Processes:}
\begin{itemize}
\item \textbf{Eluviation:} Downward translocation of soluble and suspended materials (clay, Fe, Al, humus, bases) from the A and E horizons by percolating water.
\item \textbf{Illuviation:} Deposition of those translocated materials in the B horizon, causing its enrichment in clay (argillic horizon) or Fe/Al oxides (oxic horizon).
\end{itemize}
\end{enumerate}
\end{corrige}

\begin{corrige}[Exercise 6 — Textural analysis of three soils]
\begin{enumerate}
\item \textbf{Textural classes (using the USDA textural triangle):}
\begin{itemize}
\item Sample A (75\% sand, 15\% silt, 10\% clay) $\to$ \textbf{Loamy sand} (sand >70\%, clay <15\%).
\item Sample B (20\% sand, 30\% silt, 50\% clay) $\to$ \textbf{Clay}.
\item Sample C (40\% sand, 40\% silt, 20\% clay) $\to$ \textbf{Loam}.
\end{itemize}
\item \textbf{Drainage and water retention:}
\begin{itemize}
\item Sample A (Loamy sand): Rapid drainage, low water-holding capacity, prone to drought stress; nutrients leach easily.
\item Sample B (Clay): Very slow drainage, high total water-holding capacity, but much water held tightly in micropores (unavailable to plants); risk of waterlogging and poor aeration.
\item Sample C (Loam): Moderate drainage, good water-holding capacity with a high proportion of plant-available water (mesopores); good aeration; ideal balance for most crops.
\end{itemize}
\item \textbf{Best for maize: Sample C (Loam).}  
  \textit{Reasons:} (1) Good aeration and root penetration due to balanced texture and stable structure; (2) Adequate water retention without waterlogging, providing steady moisture for maize growth; (3) Good nutrient retention (CEC) compared to sandy soils.
\item \textbf{Amendment for Sample A:} Incorporate \textbf{organic matter (compost, farmyard manure, green manure)} to increase water-holding capacity, improve aggregate stability, enhance cation exchange capacity and nutrient retention.
\end{enumerate}
\end{corrige}

\begin{corrige}[Exercise 7 — Simple soil tests]
\begin{enumerate}
\item \textbf{Jar settling test procedure (Bouyoucos method simplified):}
\begin{enumerate}
\item Take a representative soil sample, air-dry and sieve through 2 mm mesh.
\item Place $50\text{--}100\ \mathrm{g}$ in a graduated cylinder or tall jar (1 L).
\item Add water to about 2/3 full, add a dispersing agent (e.g., sodium hexametaphosphate, Calgon) and shake vigorously for 5--10 minutes to disperse aggregates.
\item Allow to settle undisturbed on a level surface.
\item After $\approx 40$ seconds: sand fraction has settled (coarse particles $>50\ \mu\mathrm{m}$). Mark level.
\item After $\approx 2$ hours: silt fraction has settled (medium particles $2\text{--}50\ \mu\mathrm{m}$). Mark level.
\item After $\approx 24\text{--}48$ hours: clay fraction settles (fine particles $<2\ \mu\mathrm{m}$); water clears. Mark final level.
\item Measure the thickness of each layer and calculate percentages: $\% = \frac{\text{layer thickness}}{\text{total sediment thickness}} \times 100$.
\end{enumerate}
\textit{Each layer represents:} bottom = sand, middle = silt, top = clay.

\item \textbf{pH 5.2} is \textbf{acidic} (pH < 7; strongly acidic if < 5.5).
\item \textbf{Crops tolerating pH 5.2:} Coffee, oil palm, cassava, sweet potato, pineapple, rubber.  
  \textbf{Amendment to raise pH:} Agricultural lime (calcium carbonate, $\ce{CaCO3}$) or dolomitic lime (calcium magnesium carbonate, $\ce{CaMg(CO3)2}$) applied based on lime requirement test.
\item \textbf{Importance of pH:} Soil pH controls nutrient availability (e.g., P, Mo low in acid soils; Al/Mn toxicity possible; N, K, S, Ca, Mg optimal near 6.0--7.0), microbial activity (nitrogen fixation, nitrification, decomposition), root growth and herbicide efficacy. Knowing pH guides liming and fertilizer choices.
\end{enumerate}
\end{corrige}

\begin{corrige}[Exercise 8 — Diagnosing a soil problem]
\begin{enumerate}
\item \textbf{Main degradation processes observed:} (1) \textbf{Wind erosion} (dust clouds carrying topsoil, especially during dry harmattan season); (2) \textbf{Surface crusting / sealing} (hard crust forms after rain due to dispersion of clay by raindrop impact, reducing infiltration, increasing runoff).
\item \textbf{Human causes:} (1) Overgrazing / removal of vegetation cover leaving soil exposed to wind and raindrop impact; (2) Continuous cultivation without organic matter return, destroying soil structure (loss of aggregation) and reducing aggregate stability.
\item \textbf{Sustainable management techniques (semi-arid context, e.g., Far North Region):}
\begin{enumerate}
\item \textbf{Mulching / crop residues:} Cover soil surface to reduce evaporation, prevent crusting by absorbing raindrop energy, protect from wind erosion, add organic matter upon decomposition.
\item \textbf{Agroforestry / windbreaks:} Plant trees/shrubs (e.g., \textit{Faidherbia albida}, \textit{Acacia senegal}) in rows perpendicular to prevailing winds to reduce wind speed, trap dust, improve fertility via nitrogen fixation and leaf litter (reverse phenology of \textit{F. albida} benefits crops).
\item \textbf{Contour bunds / zaï pits / half-moons:} Construct small earth ridges or planting pits along contours to capture runoff, increase infiltration, concentrate water and nutrients at planting stations (traditional zaï in Sahel).
\item \textbf{Crop rotation with legumes / intercropping:} Legumes (cowpea, groundnut, bambara nut) fix atmospheric nitrogen, improve soil structure via taproots, enhance fertility, reduce need for external nitrogen inputs.
\end{enumerate}
\end{enumerate}
\end{corrige}

\begin{corrige}[Exercise 9 — Structured question: the soil profile (20 marks)]
\textbf{Indicative mark scheme and examiner expectations:}
\begin{enumerate}
\item \textbf{Draw labelled profile (5 marks):}  
  - Correct vertical sequence O–A–E–B–C–R (1 mark).  
  - Distinct boundaries, reasonable thickness proportions (1 mark).  
  - Labels O, A, E, B, C, R correctly placed horizontally (3 marks).  
  \textit{Examiner tip:} Use a ruler; label horizons horizontally to the right; show organic layer (O) at top if present; do not merge A and E.

\item \textbf{Characteristic per horizon (5 marks, 1 each):}  
  O: Fresh/partially decomposed litter, dark, loose.  
  A: Mineral + humus, dark, granular, high biological activity, root mat.  
  E: Light, leached (albic), platy, low organic matter, quartz skeleton.  
  B: Clay/Fe/Al accumulation (argillic/oxic), red-brown, blocky/prismatic, cutans.  
  C: Weathered parent rock fragments (saprolite), geogenic structure, little alteration.  
  R: Unweathered bedrock, continuous, hard.

\item \textbf{Eluviation \& illuviation (4 marks):}  
  - Definition of eluviation: removal/translocation of material from upper horizons (1 mark) + dominant in E horizon (1 mark).  
  - Definition of illuviation: accumulation/deposition of translocated material in lower horizon (1 mark) + dominant in B horizon (1 mark).

\item \textbf{Role of parent rock and climate (4 marks):}  
  - Parent rock (granite): Provides mineral composition (quartz $\to$ sand/silt skeleton; feldspars $\to$ clay minerals + K/Na/Ca; micas $\to$ K + clay); influences profile depth, nutrient reserve, texture (2 marks).  
  - Climate (hot, humid): High temperature and rainfall accelerate chemical weathering (hydrolysis of feldspars) and leaching (desilication, debasification), driving deep profile development and ferralitization/lateritization (2 marks).

\item \textbf{Why A horizon most important for agriculture (2 marks):}  
  Contains most organic matter, nutrients (N, P, S), microbial life, root density, best structure (granular) and water-holding capacity; directly supports seed germination and plant growth.
\end{enumerate}
\end{corrige}

\begin{corrige}[Exercise 10 — Essay: lateritic soils of the South Region (20 marks)]
\textbf{Indicative mark scheme and examiner expectations:}
\begin{enumerate}
\item \textbf{Climatic conditions (4 marks):}  
  - Equatorial climate (Guinean type): high mean annual temperature ($\approx 25^\circ\mathrm{C}$), high rainfall ($>1500\ \mathrm{mm}$, often $>2000\ \mathrm{mm}$), short or no dry season.  
  - High relative humidity, intense insolation $\to$ year-round intense chemical weathering and leaching.

\item \textbf{Chemical weathering processes (8 marks):}  
  - \textbf{Hydrolysis:} Primary silicates (feldspars, micas) react with $\ce{H+}$/$\ce{H2O}$ $\to$ secondary clay minerals (kaolinite, halloysite) + soluble silica ($\ce{H4SiO4}$) + bases ($\ce{K+}$, $\ce{Na+}$, $\ce{Ca^2+}$, $\ce{Mg^2+}$).  
  - \textbf{Leaching (Desilication \& Debasification):} High rainfall percolates, removes soluble silica and bases from the profile $\to$ silica/alumina ratio decreases.  
  - \textbf{Concentration (Ferralitization):} Insoluble iron and aluminium oxides/hydroxides (hematite $\ce{Fe2O3}$, goethite $\ce{FeOOH}$, gibbsite $\ce{Al(OH)3}$) accumulate residually $\to$ laterite (oxic horizon).  
  - Mention: Ferralitization is the dominant pedogenic process forming Ferralsols (Oxisols).

\item \textbf{Role of vegetation, micro-organisms, fauna (5 marks):}  
  - Dense forest $\to$ abundant litter $\to$ humus formation (though rapid decomposition limits accumulation).  
  - Micro-organisms decompose organic matter, release $\ce{CO2}$ and organic acids enhancing mineral weathering.  
  - Soil fauna (earthworms, termites) mix organic/mineral matter (bioturbation), create macropores, improve structure and drainage.  
  - Nutrient cycling: rapid uptake by vegetation keeps nutrients in biomass (closed cycle), not in soil solution; slash-and-burn disrupts this.

\item \textbf{Agricultural limitations and management (3 marks):}  
  - Limitations: (1) Low cation exchange capacity (kaolinite, oxides) $\to$ low nutrient retention, leaching of applied fertilizers; (2) Phosphorus fixation by Fe/Al oxides $\to$ P deficiency; (3) Acidity (Al$^{3+}$ toxicity, Ca/Mg deficiency); (4) Hard lateritic crust (petroplinthite) limits root depth and mechanization.  
  - Management: Liming to raise pH + reduce Al toxicity + supply Ca/Mg; organic matter addition (compost, mulch) to improve CEC, structure, water retention; phosphate rock or TSP placement near roots; agroforestry (alley cropping) to recycle nutrients from depth.
\end{enumerate}
\end{corrige}

\begin{corrige}[Exercise 11 — GCE-style discussion essay (25 marks)]
\textbf{Indicative mark scheme and examiner expectations:}
\begin{enumerate}
\item \textbf{Time for fertile topsoil formation (5 marks):}  
  - Order of magnitude: \textbf{200–1000 years per cm} of topsoil under temperate conditions; slower in tropics due to rapid organic matter decomposition and intense leaching.  
  - Emphasize: Soil is \textbf{non-renewable on human timescale} (generations); a "renewable" resource only if managed sustainably.

\item \textbf{Degradation processes with named regions (9 marks, 3 processes $\times$ 3 marks each):}  
  - \textbf{Water erosion:} Sheet, rill, gully erosion on steep slopes in \textbf{North West Region (Bamenda highlands)} and \textbf{South West Region (Mount Cameroon slopes)}; loss of topsoil, sedimentation in rivers (e.g., Wouri, Mungo).  
  - \textbf{Wind erosion:} \textbf{Far North Region (Maroua, Lake Chad basin)} – dry harmattan winds remove fine topsoil particles, dust storms, dune encroachment.  
  - \textbf{Loss of fertility / nutrient mining:} \textbf{Centre/South Regions} – continuous cropping (cocoa, coffee, food crops) without fallow or fertilizer depletes N, P, K; declining yields.  
  - \textbf{Lateritic crusting (irreversible hardening):} \textbf{South Region} – repeated drying/wetting of exposed plinthite (in B horizon) forms ironstone crust (petroplinthite), sealing soil.  
  - \textbf{Pollution:} Industrial zones (\textbf{Douala, Limbe}) – heavy metals, hydrocarbons; agricultural areas – pesticide/fertilizer residues contaminating groundwater.

\item \textbf{Protection and sustainable management measures (8 marks, 4 measures $\times$ 2 marks each):}  
  - \textbf{Contour farming / terracing / bunds:} Reduces runoff velocity, increases infiltration on slopes (Bamenda, Mt Cameroon).  
  - \textbf{Agroforestry / alley cropping:} Trees (e.g., \textit{Leucaena leucocephala}, \textit{Calliandra calothyrsus}, \textit{Gliciridia sepium}) provide mulch, N fixation, windbreaks, root binding, nutrient pumping.  
  - \textbf{Conservation tillage / no-till / minimum tillage:} Minimizes soil disturbance, preserves structure, organic matter, macrofauna, reduces erosion.  
  - \textbf{Integrated Soil Fertility Management (ISFM):} Combines organic inputs (manure, compost, crop residues), mineral fertilizers (targeted), improved seeds, crop rotation/legumes.  
  - \textbf{Controlled grazing / rotational grazing / fodder banks:} Prevents overgrazing, allows vegetation recovery, reduces compaction (Adamawa, North, Far North).  
  - \textbf{Reforestation / afforestation / assisted natural regeneration:} Restores protective cover on degraded lands, watershed protection.

\item \textbf{Conclusion with justified position (3 marks):}  
  - Clear statement: \textbf{Agree} – soil formation is far slower than degradation rates; once lost (especially topsoil or via crusting), it is practically irreplaceable within human lifetimes.  
  - Justification: Cite timescales (centuries per cm), irreversibility of lateritic crusting, food security dependence on thin topsoil layer.  
  - Call for urgent policy enforcement (land-use planning, erosion control laws, fertilizer subsidies), farmer education (extension services), and research (soil mapping, suitable varieties).
\end{enumerate}
\end{corrige}

\newpage

\coursec{Summary sheet}

\begin{popfigure}
\begin{center}\tcbox[colback=popink,colframe=popink,coltext=white,arc=9pt,boxsep=4pt,left=6pt,right=6pt]{\bfseries\small THE SOIL (GEO04)}\end{center}
\vspace{-2pt}
\begin{tcbraster}[raster columns=3,raster equal height=rows,raster column skip=6pt,raster row skip=6pt,enhanced,blankest]
\begin{tcolorbox}[enhanced,colback=white,colframe=popgreen,boxrule=0.9pt,arc=3pt,title={Soil formation},fonttitle=\bfseries\scriptsize,coltitle=white,colbacktitle=popgreen,left=4pt,right=4pt,top=2pt,bottom=2pt,boxsep=1pt,toptitle=1.5pt,bottomtitle=1.5pt,halign=left]
\begin{itemize}[nosep,leftmargin=*,itemsep=1pt,font=\scriptsize]
\item Weathering: physical, chemical, biological
\item Factors: climate, rock, relief, organisms, time
\end{itemize}
\end{tcolorbox}
\begin{tcolorbox}[enhanced,colback=white,colframe=poporange,boxrule=0.9pt,arc=3pt,title={Soil profile \& layers},fonttitle=\bfseries\scriptsize,coltitle=white,colbacktitle=poporange,left=4pt,right=4pt,top=2pt,bottom=2pt,boxsep=1pt,toptitle=1.5pt,bottomtitle=1.5pt,halign=left]
\begin{itemize}[nosep,leftmargin=*,itemsep=1pt,font=\scriptsize]
\item Horizons O, A, E, B, C, R
\item Zonal, intrazonal, azonal soils
\end{itemize}
\end{tcolorbox}
\begin{tcolorbox}[enhanced,colback=white,colframe=popgold,boxrule=0.9pt,arc=3pt,title={Physical characteristics},fonttitle=\bfseries\scriptsize,coltitle=white,colbacktitle=popgold,left=4pt,right=4pt,top=2pt,bottom=2pt,boxsep=1pt,toptitle=1.5pt,bottomtitle=1.5pt,halign=left]
\begin{itemize}[nosep,leftmargin=*,itemsep=1pt,font=\scriptsize]
\item Texture \& structure
\item Porosity, colour, depth
\end{itemize}
\end{tcolorbox}
\begin{tcolorbox}[enhanced,colback=white,colframe=poppurple,boxrule=0.9pt,arc=3pt,title={Chemical \& biological},fonttitle=\bfseries\scriptsize,coltitle=white,colbacktitle=poppurple,left=4pt,right=4pt,top=2pt,bottom=2pt,boxsep=1pt,toptitle=1.5pt,bottomtitle=1.5pt,halign=left]
\begin{itemize}[nosep,leftmargin=*,itemsep=1pt,font=\scriptsize]
\item pH, humus, CEC, nutrients
\item Soil organisms
\end{itemize}
\end{tcolorbox}
\begin{tcolorbox}[enhanced,colback=white,colframe=poppink,boxrule=0.9pt,arc=3pt,title={Usage \& management},fonttitle=\bfseries\scriptsize,coltitle=white,colbacktitle=poppink,left=4pt,right=4pt,top=2pt,bottom=2pt,boxsep=1pt,toptitle=1.5pt,bottomtitle=1.5pt,halign=left]
\begin{itemize}[nosep,leftmargin=*,itemsep=1pt,font=\scriptsize]
\item Agriculture, construction, mining
\item Erosion control, rotation, fallow
\end{itemize}
\end{tcolorbox}
\end{tcbraster}

\legende{Mind map of the chapter ``Economic Geology: The Soil''.}
\end{popfigure}

\begin{aretenir}[Summary Sheet 1 -- Key Definitions]
\begin{itemize}
\item \cle{Soil}: the thin, unconsolidated superficial layer of the Earth's crust, produced by the \cle{weathering} of bedrock and the decomposition of organic matter, capable of supporting plant life.
\item \cle{Weathering}: the in-situ breakdown and alteration of rocks at or near the Earth's surface. \textbf{Physical (mechanical)} weathering: freeze--thaw (frost shattering), thermal expansion and contraction, pressure release (unloading). \textbf{Chemical} weathering: hydrolysis, oxidation, carbonation, solution. \textbf{Biological} weathering: root pressure, burrowing animals, organic acids released by organisms.
\item \cle{Soil profile}: a vertical section from the surface down to the unaltered bedrock, showing the superposed \cle{horizons} (layers) of the soil.
\item \cle{Humus}: dark, stable, decomposed organic matter formed from plant and animal remains; the main reservoir of soil fertility.
\item \cle{Soil texture}: the relative proportions of sand ($2$--$0.05$ mm), silt ($0.05$--$0.002$ mm) and clay ($<0.002$ mm) particles in a soil.
\item \cle{Soil structure}: the arrangement of soil particles into aggregates (crumb, blocky, prismatic, platy structures).
\item \cle{CEC} (cation exchange capacity): the ability of clay minerals and humus to retain and exchange nutrient cations such as \ce{Ca^{2+}}, \ce{Mg^{2+}} and \ce{K^{+}}; a high CEC indicates a fertile soil.
\end{itemize}
\end{aretenir}

\begin{aretenir}[Summary Sheet 2 -- The Soil Profile (learn the sequence!)]
\begin{center}
\begin{tabularx}{\linewidth}{|l|X|}
\hline
\rowcolor{popblueL} \textbf{Horizon} & \textbf{Description} \\
\hline
O & Organic litter: fresh and partly decomposed plant and animal remains \\
\hline
A & Topsoil: mineral matter mixed with humus; zone of \cle{eluviation} (leaching, washing out) \\
\hline
E & Bleached, strongly leached layer (not always present) \\
\hline
B & Subsoil: zone of \cle{illuviation} (accumulation of clay, iron and aluminium oxides washed down from above) \\
\hline
C & Weathered parent rock (regolith) \\
\hline
R & Unaltered bedrock \\
\hline
\end{tabularx}
\end{center}
\end{aretenir}

\begin{aretenir}[Summary Sheet 3 -- Soil Types to Know]
\begin{itemize}
\item \textbf{Zonal soils} (climate-controlled, mature profiles): lateritic/ferrallitic soils of the humid tropics (e.g.\ southern Cameroon), chernozems (black earths of temperate grasslands), podzols (cool, humid regions).
\item \textbf{Intrazonal soils} (a local factor dominates over climate): hydromorphic soils (waterlogged), saline soils, calcareous soils (rendzinas on limestone).
\item \textbf{Azonal soils} (immature, no developed profile): alluvial soils (Wouri and Benue plains), volcanic soils (fertile slopes of Mount Cameroon), regosols and lithosols (thin soils on steep slopes).
\end{itemize}
\end{aretenir}

\begin{aretenir}[Summary Sheet 4 -- Factors of Soil Formation (CLORPT)]
\textbf{Cl}imate (the dominant factor), \textbf{O}rganisms, \textbf{R}elief, \textbf{P}arent rock, \textbf{T}ime. High rainfall combined with high temperature $\Rightarrow$ intense chemical weathering and strong leaching $\Rightarrow$ deep, iron-rich, red lateritic soils, as in the humid forest zone of southern Cameroon.
\end{aretenir}

\begin{aretenir}[Summary Sheet 5 -- Properties and Their Consequences]
\begin{itemize}
\item Sandy soil: well drained, warms up quickly, but poor in nutrients and water retention; clay soil: fertile and water-retentive but heavy, sticky and poorly drained; \textbf{loam} (balanced sand--silt--clay mixture) is the ideal agricultural texture.
\item Soil \cle{pH}: controls nutrient availability; most crops prefer $5.5$--$7$; excessive acidity is corrected by liming (adding \ce{CaCO3}).
\item Porosity and permeability govern aeration, drainage and root penetration.
\end{itemize}
\end{aretenir}

\begin{aretenir}[Summary Sheet 6 -- Sustainable Management and Protection]
\begin{itemize}
\item Threats: erosion (by water and wind), leaching, overgrazing, deforestation, bush fires, pollution, salinisation, urbanisation.
\item Conservation methods: contour ploughing, terracing, strip cropping, cover crops, mulching, agroforestry, crop rotation, fallowing, application of organic manure, controlled grazing, reforestation, windbreaks.
\item Advantages of soil management: sustained fertility and crop yields, food security, erosion and flood control, protection of groundwater, conservation of biodiversity, stable farm incomes.
\end{itemize}
\end{aretenir}

\begin{attention}[Common Mistakes to Avoid]
\begin{itemize}
\item Confusing \textbf{eluviation} (washing out, horizons A and E) with \textbf{illuviation} (washing in, horizon B).
\item Saying lateritic soils are ``rich'': they are deeply weathered and strongly \emph{leached}; their fertility depends on humus content and careful management.
\item Confusing texture (particle sizes) with structure (aggregation of particles).
\item Believing soil forms quickly: a few centimetres of topsoil may need centuries to develop.
\end{itemize}
\end{attention}

\begin{methode}[GCE Examination Tips]
\begin{enumerate}
\item Always \textbf{draw and label} a soil profile when asked about horizons -- a neat, fully labelled diagram earns marks.
\item In essays on soil management, structure your answer: problem $\rightarrow$ cause $\rightarrow$ method $\rightarrow$ advantage.
\item Use Cameroonian examples: volcanic soils of Mount Cameroon, lateritic soils of the South, alluvial soils of the Benue trough, gully erosion in the North-West highlands.
\item Define every technical term the first time you use it.
\end{enumerate}
\end{methode}

\findecours

\end{document}
